\documentclass[12pt]{article}

\usepackage[margin=1in]{geometry}
\usepackage{amsmath,amssymb,amsthm}
\usepackage{ulem}
\usepackage{booktabs}
\usepackage{tabularx}
\usepackage{array}
\usepackage{enumitem}
\usepackage{lmodern}
\usepackage[T1]{fontenc}
\usepackage{microtype}
\usepackage[hidelinks]{hyperref}

% --- framework annotation system ---
\newcommand{\mtag}[1]{\textbf{[#1]}}
\newcommand{\up}[1]{\ensuremath{^{#1}}}
\newcommand{\down}[1]{\ensuremath{_{#1}}}

\title{Matter--Antimatter Pair Storm: An Emergent Statistical Framework for the Unification of Fundamental Forces}
\author{AUTHOR NAME\thanks{Author affiliation.} \\ \small Author affiliation}
\date{\today}

\begin{document}
\maketitle


\begin{quote}
\textbf{Draft v2.8 EN} (English translation of the Chinese v2.8, incorporating responses to seven referee issues M1--M7) \ensuremath{\cdot} 2026-08-24 \ensuremath{\cdot} Heuristic unification framework (reference architecture, not established physics)
\textbf{Notation convention}: \mtag{Axiom} = fundamental postulate; \mtag{Assumption} = physical assumption introduced by this framework; \mtag{Derivation} = derivation carried out within the framework; \mtag{Citation} = established physical/mathematical result; \mtag{To be verified} = conjecture requiring further verification. \textbf{Term-strength convention (v2.2)}: \textbf{Theorem} = mathematically rigorous, framework-independent (including established results); \textbf{Proposition} = deduction valid under the framework axioms; \textbf{Corollary} = direct consequence of a proposition; \textbf{Conjecture} = unproved guess; \textbf{Argument} = heuristic/structural reasoning; \mtag{Numerical} = numerical verification result.
\end{quote}



\begin{abstract}

This paper proposes a unified physical framework based on a single underlying mechanism, the \textbf{matter--antimatter pair storm}, providing an emergent statistical-mechanical account of the unification of the four fundamental forces (gravity, electromagnetism, strong and weak nuclear forces). The framework begins with four axioms: (i) the universe at its deepest level is an event flow of continuously meeting--annihilating--regenerating matter--antimatter pairs, isotropic, 50:50 balanced, at Planck-scale frequency; (ii) annihilation has an extremely small probability of "leakage," whose energy-scale-dependent form unifies the baryon asymmetry and dark energy; (iii) the arrow of time bifurcates into two manifolds that are mutual CPT mirrors, with time = frequency of annihilation events; (iv) space can be curled into time (the limit of black-hole space--time interchange).

On top of the axioms, we establish the following chain of arguments:

\begin{enumerate}
\item \textbf{Emergence of space--time}: The event flow constructs a smooth Lorentzian manifold via causal-set sprinkling and the Malament recovery theorem, giving rigorous mathematical language to "manifold = statistical average of micro-variations"; the six-direction structure of the discrete grid spans 3+3 dimensional space--time (three spatial + three temporal), and observers live on a two-dimensional "time surface" (\S{}3.3).
\end{enumerate}

\begin{enumerate}
\item \textbf{Emergence of gravity} (\S{}4, core): mass-shielded flux produces a negative pressure difference, and the flux density decays according to the \textbf{Boltzmann distribution} \ensuremath{\rho}\ensuremath{\propto}e\up{-\Phi{}/c^{2}} of the gravitational potential (derived from ergodicity + detailed balance + chaotic propagation, Proposition 2.1), giving g\down{00} = \ensuremath{-}e\up{-2\Phi{}/c^{2}} and hence \textbf{PPN parameters \ensuremath{\gamma}=1, \ensuremath{\beta}=1, all passing experimental verdicts}; the "micro-variation weight conservation" proposition fixes the spatial components; via the Jacobson thermodynamic path, the Einstein equation emerges as a horizon-thermodynamics identity. Black-hole entropy S = A/(4l\down{P}\up{2}) acquires a complete argument chain:
\begin{itemize}
\item \textbf{Microstate counting} (\S{}4.8): flux-tube coherent blocks (1 bit per 4 Planck areas, k = 2\up{1/4}), with G, \ensuremath{\hbar} appearing as composites of the Planck area;
\item \textbf{Packing proposition package} (\S{}4.10): geometric origin of the 1/4 coefficient --- \textbf{Theorem 1 (k\ensuremath{\geq}4 lower bound, new proof)} + Tammes regular tetrahedron + exact Voronoi partition + \textbf{Proposition 6 (interior-space encoding priority, framework-internal deduction)} --- five occurrences of "4" cross-validated as one source;
\item \textbf{Observer model} (\S{}4.9): observer = present-plane soliton (O1--O4), accessible information = mutual information I(R;\ensuremath{\Sigma}), holography = geometric necessity of observer limitation; "present-plane selection" (N1--N4) closed by vortex dynamics (4\ensuremath{\pi} phase locking) + spontaneous symmetry breaking; initial perturbation = \ensuremath{\varepsilon} leakage event (P1--P4, "necessary contingency").
\item \textbf{v1.21 completion of the gravity side}: (f) \textbf{isotropic exponential metric identification} --- Boltzmann g\down{00}=\ensuremath{-}e\up{-2x} + micro-variation weight conservation g\down{rr}=e\up{2x} + Axiom 1 isotropy \ensuremath{\implies} g\down{\mu{}\nu{}}=diag(\ensuremath{-}e\up{-2x}, e\up{2x}, r\up{2}e\up{2x}, r\up{2}sin\up{2}\ensuremath{\theta}\ensuremath{\cdot}e\up{2x}) (the literature exponential metric), \textbf{complete PPN parameter set \ensuremath{\gamma}=\ensuremath{\beta}=1, all others 0 (indistinguishable from GR at 1PN)}, 2PN g\down{rr} bifurcation 33\%, 3PN g\down{00} bifurcation 11\% (B3); (g) \textbf{shadow compatible with EHT} --- in the isotropic convention the framework shadow b\down{min}=2e\ensuremath{\cdot}GM/c\up{2}=5.436 GM/c\up{2} is compatible with EHT 5.2\ensuremath{\pm}0.3 (+4.55\%), correcting v1.14's "strong-field conflict with EHT" (an artifact of the standard coordinate convention) --- the framework's strong field is upgraded from "no metric commitment" to "geometric predictions compatible with observation"; (h) \textbf{direction discretization} --- Planck-tube diameter \mtag{Framework assumption: direction discretization}\ensuremath{\implies} direction cluster k=4 (Wang Hong/Tammes) + regular-tetrahedron convex covering + statistical average yielding continuous directions \ensuremath{\implies} \textbf{Bekenstein bound S\ensuremath{\leq}A/(4l\down{P}\up{2}) recovered from direction discretization}, doubly closed with the k=2\up{1/4} microstate counting ("=" directions) (B4); (i) \textbf{observer mutual-information area law} --- bulk = surface projection \ensuremath{\implies} H(R|\ensuremath{\Sigma})=0 \ensuremath{\implies} I(R;\ensuremath{\Sigma})=H(\ensuremath{\Sigma}) \ensuremath{\propto} A; the observer's raw 4 bits/plaquette compressed coherently by a factor 16 (4 plaquettes \ensuremath{\times} 4 directions) \ensuremath{\implies} entropy density 1/4 bit/plaquette = 1/(4l\down{P}\up{2}) (B5).
\end{itemize}
\end{enumerate}

\begin{enumerate}
\item \textbf{Gauge-group structure}: The Standard Model gauge group SU(3)\ensuremath{\times}SU(2)\ensuremath{\times}U(1) corresponds to the representation-theoretic structure of "pairing orders" n=3,2,1 (\textbf{v2.5: uniqueness via three-criterion structural screening} --- d\down{min} matching / same-sign pairing order / third-order antisymmetric tensor, 10 candidates eliminated one by one, \S{}5.2); color confinement acquires its mathematical root in the third-order antisymmetric tensor; the three fermion generations correspond to three winding modes of the vortex; the multi-representation unified field gives the four forces as a partial-wave decomposition of flux-field fluctuations. \textbf{v1.19 gauge-side emergence arguments} (\S{}5.6--5.8): (a) \textbf{quantum-double algebra proposition package} --- the framework's three event-flow operations \ensuremath{\leftrightarrow} Hopf structure, non-Abelian gauge structure emerging exactly from discrete algebra (Kitaev theorem, Z\down{2}/Z\down{3}/S\down{3} numerical verification); (b) \textbf{framework Hamiltonian} H\down{F} = \ensuremath{-}J\ensuremath{\Sigma}A\down{v} \ensuremath{-} K\ensuremath{\Sigma}B\down{p} --- [H\down{F},A\down{v}]=[H\down{F},B\down{p}]=0 proves gauge symmetry is a dynamical symmetry (emergent), 2+1D positive energy gap (topological order) is a low-dimensional rigorous realization of the mass gap, torus degeneracy = |G|\up{2}; (c) \textbf{Wen string-net 3+1D route} --- quantum lines = string-net edges (Framework assumption: cosmic strings = dynamically emergent network), four forces = emergent channels of string-net order (EM\ensuremath{\to}U(1) string-net photon, gravity\ensuremath{\to}string-net emergent spin-2); \textbf{v1.28--1.30 rigorization}: rigorous proof of \ensuremath{\lambda}=0 string-net condensation (E0=\ensuremath{-}N\down{c}, degeneracy 2\up{n-N\down{c}}, \ensuremath{\Delta}=2) + rigorous duality mapping + \textbf{Z\down{2} \ensuremath{\subset} U(1) group-embedding theorem + Maxwell linearization proposition + numerical verification of Coulomb-phase photon dispersion} (continuum-limit renormalization group = world-class open problem) + \textbf{spin 2 = rigorous verification of second-order tensor rotation eigenvalues \ensuremath{\pm}2} (framework spin-2 mechanism); \textbf{v1.34/1.36 deepening}: \textbf{criterion (4) of the string-net low-energy spectrum numericalized} (C4 helicity decomposition, spin-2 mode gapped \ensuremath{\Delta}\ensuremath{\pm}2\ensuremath{\approx}0.49, carried by second-order tensor sector) + \textbf{3+1D \ensuremath{\lambda}\down{c}\ensuremath{\approx}0.14} (locked by four system sizes, forming a dimension--critical-coupling decreasing sequence with 2D toric code) + \textbf{Wilson RG path} (U(1) Gaussian fixed point + 2D Ising block RG flowing to the continuum limit --- Theorem A + Proposition B + Numerical verification C, four-layer closure) --- the string-net path upgraded from semi-rigorous to partially rigorous, cross-validated with the gravitational statistical-mechanics emergence (\S{}4) as two independent paths. \textbf{v1.20 confined string-tension numerics} (\S{}5.7): (d) \textbf{X-field charge pair} \ensuremath{\sigma}\down{E}\ensuremath{\approx}2\ensuremath{\kappa} --- three system sizes (3\ensuremath{\times}3/3\ensuremath{\times}4/4\ensuremath{\times}4) thermodynamic-limit convergence <0.02\%, \ensuremath{\kappa}=0 strictly deconfined; (e) \textbf{Z-field flux pair} \ensuremath{\sigma}\down{m}\ensuremath{\approx}2\ensuremath{\lambda} (constrained-sector method: B\down{p}-sector invariant subspace) --- 0.21/1.09/3.04 at \ensuremath{\lambda}=0.5/1.0/2.0 (3\ensuremath{\times}4), \ensuremath{\lambda}=0 strictly deconfined --- the framework's "confinement = flux-tube linear potential" \mtag{Framework mechanism: confinement = flux-tube linear potential} acquires quantitative numerical realization (early numerics superseded after P0-1 review); (f) \textbf{transition-point localization (P0-3 open-boundary FSS \ensuremath{\to} v1.23 torus refinement)} --- order parameter m=\ensuremath{\langle}X\ensuremath{\rangle} crossing: open-boundary FSS \ensuremath{\kappa}\down{c}\ensuremath{\approx}0.25\ensuremath{\pm}0.05 (v1.22); \textbf{torus (periodic boundary) refinement: half-polarization point \ensuremath{\kappa}*\ensuremath{\approx}0.20\ensuremath{\pm}0.01} (2\ensuremath{\times}2/2\ensuremath{\times}3/2\ensuremath{\times}4/3\ensuremath{\times}3 size-independent, boundary pollution removed); \textbf{criterion correction} --- the m=0.5 crossing is a half-polarization point, not the transition point (torus measurement m(h=0.32)\ensuremath{\approx}0.84), true transition \textbf{\ensuremath{\kappa}\down{c}\ensuremath{\approx}0.25--0.32, consistent with the literature X-field (\ensuremath{\sigma}\up{x} field) value h\down{xc}\ensuremath{\approx}0.32} (note: h\down{c}=0.328 is the \ensuremath{\sigma}\up{z-field} value, mixed up in P0-3, corrected in v1.23) --- topological-order\ensuremath{\to}confinement transition confirmed, confinement phase diagram complete (\S{}5.7).
\end{enumerate}

\begin{enumerate}
\item \textbf{Foundations of quantum mechanics}: The classical event flow cannot produce Bell violation; the framework must be non-commutativized --- the event flow is upgraded to an operator algebra, with CHSH = 2\ensuremath{\sqrt{2}} and complementarity V\up{2}+I\up{2}\ensuremath{\leq}1 proven; the positioning is 't Hooft-like: providing an ontological motivation ("event flow does not commute") for the Hilbert-space structure of quantum mechanics.
\end{enumerate}

\begin{enumerate}
\item \textbf{Cosmology}: The energy-scale dependence of the leakage probability \ensuremath{\varepsilon}(E) (\ensuremath{\beta}\ensuremath{\approx}0.615) unifies the baryon asymmetry (\ensuremath{\eta}\ensuremath{\approx}6\ensuremath{\times}10\up{-10}) and the dark-energy density (\ensuremath{\rho}\down{\Lambda{}}\ensuremath{\sim}(meV)\up{4}); the statistical decay factor \ensuremath{\varepsilon}\down{B}=(l\down{P}/R\down{U})\up{2} yields the dark-matter density; dark matter = superposition model of this-branch Planck-fluctuation remnants + mirror-branch gravitational projection; \textbf{the initial perturbation (\ensuremath{\varepsilon}\down{0}=7.7\ensuremath{\times}10\up{-20}) shares one root with baryon asymmetry and dark energy --- three sources, one origin} (\S{}4.9 P4).
\end{enumerate}

The framework yields a number of testable predictions, including: \textbf{black-hole strong-field prediction package (B1--B5)} (extreme-Kerr Hawking-radiation angular distribution deviating from axisymmetry, primordial black-hole entropy volume correction \ensuremath{\sim}10\up{-8}, Hawking-radiation--infalling-matter entanglement, ER-bridge packing bound, non-thermal correction 1/S), \textbf{EMRI precession-difference verdict} (S-E1--S-E4: \textbf{after v1.35 correction, \ensuremath{\Delta}lnL criterion 27\ensuremath{\sigma}/28\ensuremath{\sigma} (robust with a-\ensuremath{\varphi}0 fully free, \ensuremath{\sigma}\down{\delta{}}=1.6\ensuremath{\times}10\up{-6})}), \textbf{FQH graviton-mode \ensuremath{\varepsilon} term} (S-E5, framework prediction), accretion-disk firewall non-thermal correction (S4, primordial black-hole window), void--mirror-branch gravitational-projection correlation (void-CPT-\ensuremath{\kappa} anomaly), constant deviation of dark energy w=\ensuremath{-}1+4\ensuremath{\varepsilon}/3, SU(3) \ensuremath{\beta}-function staircase jump, ln 2 factor in the entanglement-entropy area law.

\textbf{Honesty statement}: This framework is a heuristic unification framework, not established physics. All \mtag{Derivation}s in this paper are internally consistent arguments within the framework, and no \mtag{To be verified} item may be used as a basis for derivation. \textbf{The v1.0 strong-field metric extrapolation prediction (g\down{00}=\ensuremath{-}e\up{-2x} neutron-star deviation) has been formally withdrawn} (beyond the domain of validity, \S{}4.6); the framework's strong-field content is the seven-loop mechanism chain + isotropic exponential metric (shadow compatible with EHT, \S{}4.5b/\S{}4.6). Remaining major open problems: complete emergence of the non-Abelian component (4D), concrete values of J/K coupling, string-net emergent U(1) continuum limit and spin 2 (3+1D) (\S{}10.2).



\end{abstract}
\section{1 Introduction}

\subsection{1.1 Motivation and background}

The two pillars of modern physics --- general relativity (GR) and quantum field theory (QFT) --- precisely describe gravity and the other three fundamental forces respectively, yet their unification (quantum gravity) is the deepest unsolved problem of physics. The Standard Model unifies strong, weak and electromagnetic forces via the SU(3)\ensuremath{\times}SU(2)\ensuremath{\times}U(1) gauge group; GR describes gravity via Riemannian geometry. However:

\begin{itemize}
\item Gravity is non-renormalizable and lacks a complete quantum description;
\item The Standard Model contains 19--26 free parameters (hierarchy problem, cosmological-constant problem, strong-CP problem);
\item The relative strengths of the four forces (1 : 10\up{-2} : 10\up{-6} : 10\up{-38}) lack a unified explanation;
\item Candidate paths such as grand unification (GUT), string theory, and loop quantum gravity have all lacked experimental support.
\end{itemize}

This framework proposes a different unification strategy: \textbf{instead of postulating a larger symmetry group or extra dimensions, it seeks a single underlying statistical mechanism (the event flow) such that the four forces become emergent modes of the same mechanism at different scales/configurations}. This strategy is analogous to the relation of statistical mechanics to thermodynamics: GR is the geometric description, and the framework's statistical-fluid description is its effective theory.

\subsection{1.2 Core concept: the matter--antimatter pair storm}

\mtag{Axiom 1}The universe at its deepest level is an event flow of continuously meeting--annihilating--regenerating matter--antimatter pairs, isotropic, 50:50 balanced, at Planck-scale frequency (\ensuremath{\nu} \ensuremath{\sim} t\down{P}\up{-1} \ensuremath{\approx} 10\up{43} s\up{-1}). Each encounter is a "micro-variation"; the space--time manifold = statistical average of micro-variations ("the manifold varies infinitely finely").

The four fundamental forces = four modes of the same storm:

\begin{table}[htbp]
\centering
\small
\begin{tabularx}{\textwidth}{@{}lXX@{}}
\toprule
Force & Mode & Mechanism \\
\midrule
Strong & Shielding harbor & Nearby particles share a low-pressure cavity (intuition for color confinement) \\
Weak & Draft/wind & Extremely short-range pipe coupling (chirality, W/Z masses) \\
Electromagnetic & Selective coupling & Long-range flux-pressure gradient (cavity model) \\
Gravity & Flux negative pressure & Long-range, extremely weak flux-occlusion deficit \\
\bottomrule
\end{tabularx}
\end{table}

\textbf{Unification = disappearance of mode-boundary distances}: the three energy scales --- electroweak unification (100 GeV), GUT (10\up{15} GeV), Planck (10\up{19} GeV) --- are a ladder on which mode boundaries vanish one by one; \textbf{symmetry breaking = mode boundaries re-emerge as flux density decreases}.

\subsection{1.3 Paper structure}

\S{}2 establishes the four axioms and fundamental objects; \S{}3 argues the emergent generation of space--time (causal sets, time surface, 3+3 dimensions); \S{}4 is the core chapter --- the complete emergence argument for gravity (Boltzmann distribution, PPN verdicts, microstate counting, thermal-equilibrium propositions, \textbf{seven-loop strong-field mechanism, packing proposition package, observer model O1--O4, present-plane selection N1--N4, initial-perturbation origin P1--P4}); \S{}5 gauge-group structure; \S{}6 dynamics (least action); \S{}7 foundations of quantum mechanics; \S{}8 cosmology; \S{}9 testable predictions; \S{}10 conclusions and open problems.


\section{2 Axioms and Basic Postulates}

\subsection{2.1 The four axioms}

\textbf{\mtag{Axiom 1 (matter--antimatter pair storm)}}: The universe at its deepest level is an event flow of continuously meeting--annihilating--regenerating matter--antimatter pairs, isotropic, 50:50 balanced, at Planck-scale frequency; the manifold = statistical average of micro-variations.

\textbf{\mtag{Axiom 2 (leakage mechanism)}}: Each annihilation has a small probability \ensuremath{\varepsilon} of binding failure; the positive fraction leaks into the positive-time-arrow branch, the anti fraction into the mirror branch. \ensuremath{\varepsilon} is energy-scale dependent: \ensuremath{\varepsilon}(E) = \ensuremath{\varepsilon}\down{0}(M\down{Pl}/E)\up{\beta{},} \ensuremath{\beta} \ensuremath{\approx} 0.615.

\textbf{\mtag{Axiom 3 (bifurcation of the time arrow)}}: The two branches are mutual CPT mirrors; time = encounter/annihilation frequency, defined by the generation rate of the overlap integral \ensuremath{\tau}(x) = \ensuremath{\int}\ensuremath{\bar{\psi}}\ensuremath{\psi} d\up{4}x.

\textbf{\mtag{Axiom 4 (space curled into time)}}: Space can be curled into time; the space--time interchange inside a black hole = space curled to its limit (radial direction curled into the frame-dragging direction).

\subsection{2.2 Fundamental objects and definitions}

\begin{itemize}
\item \textbf{Micro-variation weight current J\up{\mu{}}}\mtag{Derivation}: the conserved current of the annihilation event flow (\ensuremath{\nabla}\down{\mu{}} J\up{\mu{}} = 0, for \ensuremath{\varepsilon} \ensuremath{\ll} 1). Micro-variation weight W[\ensuremath{\Sigma}] = \ensuremath{\int}\down{\Sigma{}} J\up{\mu{}} d\ensuremath{\Sigma}\down{\mu{}}.
\item \textbf{Flux density \ensuremath{\rho}\down{pair}(x)}: local annihilation-event rate density, determining the flow rate of time.
\item \textbf{Quantum line}: the correlation channel between the two branches established by leakage events of common origin (the underlying ontology of ER=EPR).
\item \textbf{Plaquette}: the fundamental cell constituting the manifold (six-directional connections, spanning three planes).
\end{itemize}

\subsection{2.3 Hierarchical structure of mechanics}

\begin{verbatim}
Event flow (Planck noise)
  |-- Statistical average => manifold / space--time
  |-- Solitonization => particles (stable localized structures)
  |-- Pressure resistance => mass / inertia (equivalence principle = two sides of the same coupling)
  |-- Flux response => four-force modes => rising energy scale => unification / breaking
  `-- Leakage eps => time arrow / baryon asymmetry / dark energy
\end{verbatim}



\section{3 Emergence of Space--Time}

\subsection{3.1 Axiomatization via causal sets}

\mtag{Citation + Derivation} Formalize the Axiom-1 event flow as a \textbf{causal set}: an event set C with partial order \ensuremath{\prec}, satisfying non-reflexivity, transitivity, and local finiteness. Each encounter--annihilation \ensuremath{\in} C; "can influence each other" = \ensuremath{\prec}. Isotropy corresponds to \textbf{Poisson sprinkling}: random points placed on a smooth Lorentzian manifold (M,g) as a Poisson process with density \ensuremath{\rho} = l\down{P}\up{-4}.

\textbf{Malament recovery theorem}\mtag{Citation}: If (M,g) and (M\ensuremath{^\prime},g\ensuremath{^\prime}) are both past/future-distinguishing manifolds with isomorphic causal orders, then M \ensuremath{\cong} M\ensuremath{^\prime} and g\ensuremath{^\prime} = \ensuremath{\Omega}\up{2}g --- the causal structure uniquely determines the conformal class of the metric. The conformal factor is fixed by event counting:

\[g_{\mu\nu}(x) = \Omega^2(x)g_0{}_{\mu\nu}(x),\qquad \Omega^4(x) = \frac{dN(x)}{\rho\,dV_0(x)}\]

\textbf{Three-step proposition chain for "manifold micro-variation"}\mtag{Derivation}:
\begin{enumerate}
\item Poisson sprinkling (law of large numbers) \ensuremath{\to} causal partial order;
\item Malament recovery theorem \ensuremath{\to} conformal class of metric;
\item Counting density \ensuremath{\to} conformal factor \ensuremath{\Omega}.
\end{enumerate}

\subsection{3.2 Time = event frequency}

\mtag{Derivation} For a static metric ds\up{2} = g\down{00} dt\up{2} + h\down{ij} dx\up{i} dx\up{j,} an ideal clock = event counter. Events occur uniformly at Planck frequency f\down{0} in proper time; the coordinate-time event rate f(x):

\[g_{00}(x) = -\left(\frac{f(x)}{f_0}\right)^2,\qquad \frac{f_B}{f_A} = \sqrt{\frac{g_{00}(B)}{g_{00}(A)}} \approx 1 + \frac{\Phi_B - \Phi_A}{c^2}\]

The weak field gives gravitational redshift --- \textbf{frequency gradient corresponds one-to-one to the g\down{00} gradient}.

\subsection{3.3 3+3 dimensional discrete space--time and the time surface}

\mtag{Derivation} The six directions of the plaquette (front/back, left/right, up/down) span three coordinate axes and three planes (xy/xz/yz), corresponding respectively to \textbf{past/present/future --- three temporal dimensions}. We perceive only one temporal dimension because the other two are curled into the normal directions of the spatial planes.

\textbf{Corollary (holography = geometric necessity of the time surface)}\mtag{Derivation, complete proposition chain in \S{}4.9}:
\begin{itemize}
\item Observer = local soliton on the present plane (O1); accessible information = projection onto the time surface = mutual information I(R;\ensuremath{\Sigma}) (O2);
\item Bulk information is redundant, surface information is primary; information capacity \ensuremath{\propto} surface area (O4);
\item The observer's uniform prior is derived from observer geometry (O3) --- no information channel along external spatial directions.
\end{itemize}

\textbf{Why the "present plane" is selected}\mtag{Derivation, proposition chain in \S{}4.9}: the vortex must lock its rotation axis (N1) and exactly one plane (N2); the three planes are equivalent (N3, "present" = definition of the selected plane); the initial perturbation (= \ensuremath{\varepsilon} leakage event, P1--P4) breaks the SO(3) degeneracy (N4).

\subsection{3.4 Open problems of the continuum limit}

\mtag{To be verified} (1) prove "event set = Poisson sprinkling" from annihilation dynamics; (2) convergence of the continuum limit (not fully resolved in standard causal-set theory either); (3) proof of the postulate "each event = 1 proper-time element".


\section{4 Emergence of Gravity (Core Chapter)}

\subsection{4.1 Flux occlusion \ensuremath{\to} gravitational potential}

\mtag{Derivation} An object of mass M "swallows" flux, producing a steady-state deficit. In source-free regions the flux deficit satisfies \ensuremath{\nabla}\up{2}(\ensuremath{\delta}\ensuremath{\rho}) = 0, with spherically symmetric solution \ensuremath{\delta}\ensuremath{\rho} \ensuremath{\propto} 1/r; matching the Newtonian potential via the time component:

\[\frac{\rho_{pair}(r)}{\rho_{pair}(\infty)} = 1 - \frac{GM}{c^2 r} \quad\text{(first order)}\]

Weak-field numerical check: at the solar surface GM/c\up{2}R = 2.1\ensuremath{\times}10\up{-6} \ensuremath{\ll} 1 \checkmark.

\textbf{Factor-2 explanation (v1.12)}: The flux-deficit coefficient is taken as \textbf{1} (GM/c\up{2}r, same form as the Newtonian potential \ensuremath{\Phi}/c\up{2}), not 2GM/c\up{2}r --- \textbf{the "2" in g\down{00} comes from the squared structure g\down{00} = \ensuremath{-}(\ensuremath{\rho}/\ensuremath{\rho}\ensuremath{\infty})\up{2}} (\S{}4.3; Axiom 3: time = event frequency \ensuremath{\implies} \ensuremath{\sqrt{-g\down{00}}}=\ensuremath{\rho}/\ensuremath{\rho}\ensuremath{\infty}). That is:
\begin{itemize}
\item Flux deficit \ensuremath{\rho}\ensuremath{\infty}\ensuremath{-}\ensuremath{\rho}\down{ss} = \ensuremath{\rho}\ensuremath{\infty}\ensuremath{\cdot}GM/(c\up{2}r): \textbf{first-order coefficient 1} (= \ensuremath{\Phi}/c\up{2});
\item g\down{00} = \ensuremath{-}(1\ensuremath{-}GM/(c\up{2}r))\up{2} = \ensuremath{-}(1\ensuremath{-}2GM/(c\up{2}r)+\ldots): \textbf{first-order coefficient 2} (consistent with GR, \ensuremath{\beta}\down{ppn}=1 self-consistent).
\end{itemize}
This also fixes the flux-vanishing radius of the finite-R steady state (v1.11) as \textbf{R = GM/c\up{2}} (not the Schwarzschild 2GM/c\up{2}). Numerical verification: R=GM/c\up{2} agrees with the Boltzmann weak field (deviation <0.6\% for r\ensuremath{\geq}10), while R=2GM/c\up{2} deviates >11\%; the squared g\down{00} structure differs from the GR weak field by 1.0\% at r=10 and 0.01\% at r=100.

\textbf{Strong-field bifurcation verdict (v1.13)}\mtag{Proposition}: the PPN \ensuremath{\beta} parameters of the two candidate paths bifurcate --- the Boltzmann g\down{00}=\ensuremath{-}e\up{-2\Phi{}/c^{2}} expansion has \ensuremath{\beta}=1 (passes observation 1\ensuremath{\pm}3\ensuremath{\times}10\up{-4}), while the squared steady-state g\down{00}=\ensuremath{-}(1\ensuremath{-}R/r)\up{2} expansion has \ensuremath{\beta}=1/2 (\textbf{excluded at 1600\ensuremath{\sigma}}). \ensuremath{\implies} \textbf{The Boltzmann path is the correct description from weak to intermediate field}; the squared steady state is only a first-order approximation (still valid as a flux profile, but not an exact g\down{00} mapping). \textbf{The v1.12 horizon prediction r=GM/c\up{2} is withdrawn} (it was based on the squared steady state with g\down{00}=0, excluded by the \ensuremath{\beta} verdict). Framework strong field: the Boltzmann factor never vanishes \ensuremath{\implies} no g\down{00}=0 geometric horizon; the horizon = \textbf{storm overload / flux-vanishing equipotential surface} (physical mechanism, framework assumption) --- consistent with v1.1's "strong field undefined" and the seven-loop mechanism.

\subsection{4.2 Boltzmann distribution (flux-field thermal equilibrium)}

\mtag{Derivation + thermal-equilibrium proof} The flux density decays with the \textbf{Boltzmann weight} of the gravitational potential:

\[\frac{\rho(r)}{\rho_\infty} = e^{-\Phi(r)/c^2} = 1 - x + \tfrac12 x^2 - \cdots,\qquad x = \frac{GM}{c^2 r}\]

\textbf{Thermal-equilibrium proposition (six-step chain)}\mtag{Derivation}:
\begin{enumerate}
\item The event flow is a stationary random process (SDE dX\down{t} = b dt + \ensuremath{\sigma} dW);
\item Ergodicity (Birkhoff--Khinchin theorem): time average = ensemble average;
\item Microscopic reversibility (annihilation--regeneration CPT symmetry);
\item Chaotic propagation (Boltzmann--Grad limit, linking Hilbert's sixth problem / the Deng--Hani--Ma hard-sphere chain) \ensuremath{\implies} Boltzmann equation;
\item Detailed balance \ensuremath{\implies} Boltzmann distribution \ensuremath{\rho} \ensuremath{\propto} e\up{-E/kT};
\item kT = mc\up{2} (relativistic limit) \ensuremath{\implies} \ensuremath{\rho} \ensuremath{\propto} e\up{-\Phi{}/c^{2}}.
\end{enumerate}

\textbf{Numerical simulation verification}\mtag{Derivation}: detailed-balance Markov chain converges to the Boltzmann distribution (deviation 2.12\%); ergodicity (first/second-half time averages differ <0.2\%).

\textbf{Important limitation}\mtag{To be verified}: the Boltzmann distribution of a pure 1/r potential diverges (r\ensuremath{\to}0) --- gravitational systems have no global thermal equilibrium (consistent with standard statistical mechanics); e\up{-\Phi{}/c^{2}} is valid only locally / in the weak field / with cutoff, and its weak-field expansion agrees with \S{}4.1.

\textbf{\mtag{v2.6 supplement: numerical proof of Boltzmann emergence + response to the circularity objection (referee M1)}}:
The referee objected that the Boltzmann "proof" is circular (assuming thermal equilibrium to prove thermal equilibrium). \textbf{Decisive numerical response} (`\down{m1}\down{boltzmann}.py`, inputs only conservation laws + collisions, initial extreme non-equilibrium single-energy-shell): (1) \textbf{event-flow language (DSMC, N=2000)} --- temperature emergence 0.00\% deviation, velocity distribution converges to Maxwell--Boltzmann (weighted deviation 7.0\%), \textbf{Boltzmann H-function net irreversible decrease \ensuremath{\Delta}H=\ensuremath{-}2.27}, ergodicity 0.00\% --- \textbf{thermal equilibrium is an emergent product of the long-time limit, not an input assumption}; (2) \textbf{deterministic hard disks (no molecular-chaos assumption)} --- temperature 0.00\%, net H decrease \ensuremath{\Delta}H=\ensuremath{-}2.38 (2D long-time tail causes \ensuremath{\sim}18\% distribution deviation, honestly noted); (3) \textbf{explicit preconditions of the derivation chain} --- stationarity (Axiom 1 isotropy), ergodicity/chaotic propagation (dynamical properties, measured 0.00\%), detailed balance (time-reversal reversibility of the event flow) --- all from the Axiom-1 level, \textbf{not thermal-equilibrium assumptions}.

\textbf{Candidate event-flow Hamiltonian (H\down{event})}: the event flow is upgraded from metaphor to a computable model --- classical layer H=\ensuremath{\Sigma}p\up{2}/2m+\ensuremath{\Sigma}V\down{contact} (collision = encounter--annihilation--regeneration); statistical layer S=Tr f(D\up{2}/\ensuremath{\Lambda}\up{2}) (spectral action, \S{}5.4); quantum layer open-system Lindblad d\ensuremath{\rho}/dt=\ensuremath{-}i[H\down{0},\ensuremath{\rho}]+\ensuremath{\Sigma}(L\ensuremath{\rho}L\ensuremath{^\dagger}\ensuremath{-}\ensuremath{\tfrac12}{L\ensuremath{^\dagger}L,\ensuremath{\rho}}) (annihilation = dissipation / regeneration = pumping; steady state = detailed balance \ensuremath{\implies} Boltzmann). Falsifiable: each layer gives predictions for collision rates / relaxation timescales / steady-state distributions.

\textbf{Connection with LHC}: the single-event energy scale E\ensuremath{\sim}E\down{P}\ensuremath{\approx}10\up{19} GeV \ensuremath{\gg} LHC \ensuremath{\sim}10\up{4} GeV (15 orders of magnitude) \ensuremath{\implies} the event flow manifests only through collective modes, \textbf{predicting no new physics at LHC} (consistent with null results); isotropy + Poisson sprinkling \ensuremath{\implies} no Lorentz/CPT violation; leakage \ensuremath{\varepsilon}\ensuremath{\sim}10\up{-50} produces no LHC signal; the framework's discriminating predictions are concentrated on the gravity side (S-E1/E5/S4). \textbf{Honest caveat}: the event flow remains Axiom 1 (an input, like the Standard Model Lagrangian); the microscopic structure of H\down{quantum} and an analytic proof of ergodicity remain open (\S{}10.2 \#4).

\subsection{4.3 Time component and PPN \ensuremath{\gamma}=1}

\mtag{Derivation} From the lemma g\down{00} = \ensuremath{-}(\ensuremath{\rho}/\ensuremath{\rho}\ensuremath{\infty})\up{2}:

\[g_{00} = -e^{-2\Phi/c^2} = -\big(1 - 2x + 2x^2 - \tfrac43 x^3 + \cdots\big)\]

\subsection{4.4 Micro-variation weight conservation proposition (spatial components)}

\textbf{\mtag{Proposition (micro-variation weight conservation \ensuremath{\Rightarrow} \ensuremath{\gamma}=1)}}: If the annihilation event flow satisfies conservation (Axiom 1\ensuremath{^\prime}) and local isotropy (Axiom 1), then in a static spherically symmetric weak field g\down{rr} = \ensuremath{-}1/g\down{00} = (1\ensuremath{-}2GM/c\up{2}r)\up{-1}, hence \textbf{PPN \ensuremath{\gamma} = 1}.

\textbf{Proof} (three lines): (1) local inertial-frame isotropy \ensuremath{\Rightarrow} the temporal weight \ensuremath{\sqrt{-g\down{00}}} and radial weight \ensuremath{\sqrt{g\down{rr}}} are complementary (product = 1); (2) lemma: \ensuremath{\sqrt{-g\down{00}}} = f(r)/f\down{0} = e\up{-x}; (3) substitution \ensuremath{\Rightarrow} \ensuremath{\gamma}=1. \ensuremath{\blacksquare}

\subsection{4.5 PPN \ensuremath{\beta}=1 (second-order term)}

\mtag{Derivation} The second-order coefficient of the Boltzmann distribution is a = +1/2:

\[g_{00} = -e^{-2x} = -\big(1 - 2x + 2\beta x^2 + \cdots\big) \quad\Longrightarrow\quad \beta_{PPN} = 1\]

\textbf{Numerical verification}: PPN expansion comparison (Will convention g\down{00} = \ensuremath{-}(1\ensuremath{-}2x+2\ensuremath{\beta}x\up{2})) \ensuremath{\implies} \ensuremath{\beta} = a + 1/2 = 1 \checkmark.

\subsection{4.5b Complete metric = isotropic exponential metric + complete PPN parameter set (v1.21)}

\textbf{\mtag{Derivation + Numerical} Three independent derivations converge to a unique coordinate form} --- Boltzmann g\down{00}=\ensuremath{-}e\up{-2x} (\S{}4.2) + micro-variation weight conservation g\down{rr}=e\up{2x} (\S{}4.4) + Axiom-1 isotropy (g\down{ij}=\ensuremath{\delta}\down{ij}\ensuremath{\cdot}e\up{2x}):

\[g_{\mu\nu} = \mathrm{diag}\!\big(-e^{-2x},\ e^{2x},\ r^2 e^{2x},\ r^2\sin^2\theta\, e^{2x}\big),\qquad x = \frac{GM}{c^2 r}\]

Consistency check: micro-variation weight conservation g\down{rr}\ensuremath{\cdot}(\ensuremath{-}g\down{00}) = e\up{2x}\ensuremath{\cdot}e\up{-2x} = \textbf{1 exactly} --- the three derivations are coordinate-consistent. \mtag{Citation} This metric is the exponential metric known in the literature (Edery 2005 et al.) --- the framework derives it independently from axioms and connects directly to the mainstream literature.

\textbf{Complete PPN parameter set}: static isotropic purely diagonal metric \ensuremath{\implies}

\begin{table}[htbp]
\centering
\small
\begin{tabularx}{\textwidth}{@{}lXXX@{}}
\toprule
Parameter & \ensuremath{\gamma} & \ensuremath{\beta} & \ensuremath{\xi}, \ensuremath{\alpha}\down{1}, \ensuremath{\alpha}\down{2}, \ensuremath{\alpha}\down{3}, \ensuremath{\zeta}\down{1}--\ensuremath{\zeta}\down{4}, \ensuremath{\zeta}\down{W} \\
\midrule
Framework value & \textbf{1} & \textbf{1} & \textbf{0 (automatic GR values)} \\
\bottomrule
\end{tabularx}
\end{table}

\textbf{Indistinguishable from GR at 1PN}; bifurcation points precisely localized (Taylor expansion comparison):

\begin{table}[htbp]
\centering
\footnotesize
\begin{tabularx}{\textwidth}{@{}lXXXXX@{}}
\toprule
Order & Framework g\down{00} & GR g\down{00} & Framework g\down{rr} & GR g\down{rr} & Bifurcation \\
\midrule
x\up{1} & \ensuremath{-}2 & \ensuremath{-}2 & +2 & +2 & none (1PN) \\
x\up{2} & +2 & +2 & \textbf{+2} & \textbf{+1.5} & g\down{rr} 33\% (2PN) \\
x\up{3} & \textbf{\ensuremath{-}4/3} & \textbf{\ensuremath{-}3/2} & +4/3 & +0.5 & g\down{00} 11\% (3PN) \\
\bottomrule
\end{tabularx}
\end{table}

\textbf{Numerical verification (geodesic RK4)}: perihelion precession is 0.57\%/1.02\%/1.55\% larger than GR (relative difference \ensuremath{\delta}=0.30\ensuremath{\cdot}x\down{p} linear, quantifying the 2PN bifurcation); photon deflection larger than GR (5\% difference at b=10). 2PN observable effects \ensuremath{\propto}U\up{2} are currently unresolvable (binary pulsars); 3PN requires future LISA EMRI / pulsar-timing-array tests. \textbf{EMRI verdict program (v1.26--v1.35)}\mtag{Numerical}: EMRI (q\ensuremath{\sim}10\up{4}--10\up{7}) orbital cycles N\ensuremath{\sim}10\up{4}--10\up{5} amplify the precession difference into a measurable phase shift --- \textbf{\ensuremath{\Delta}\ensuremath{\Phi}\down{prec} = N\ensuremath{\cdot}\ensuremath{\Delta}\ensuremath{\omega}\down{GR}\ensuremath{\cdot}\ensuremath{\delta}} (\ensuremath{\delta}=0.30\ensuremath{\cdot}x\down{p}, x\down{p}=GM/c\up{2}p). Full waveform template + Fisher precision (v1.27): Peters--Mathews adiabatic evolution + dominant harmonic h(t)=h\down{0}\ensuremath{\cdot}cos[2\ensuremath{\lambda}+2(1+\ensuremath{\delta})\ensuremath{\omega}\down{peri}+\ensuremath{\varphi}\down{0}] + LISA noise (Barack--Cutler 2004) + \textbf{7-parameter Fisher matrix} {\ensuremath{\delta},lnM,ln\ensuremath{\mu},a\down{0},e\down{0},\ensuremath{\varphi}\down{0},lnD}: representative source (M=10\up{6} M\ensuremath{\odot}, \ensuremath{\mu}=10 M\ensuremath{\odot}, e\down{0}=0.3, p\down{0}=15M, D=3 Gpc) \ensuremath{\delta}\down{true}=2.00\% \ensuremath{\implies} \ensuremath{\sigma}\down{\delta{}}=8.3\ensuremath{\times}10\up{-4} (24\ensuremath{\sigma}, 2yr) / 1.2\ensuremath{\times}10\up{-4} (167\ensuremath{\sigma}, 4yr) --- values later corrected (see v1.35). Kerr + Bayesian analysis (v1.32): \ensuremath{\delta}\down{hat}=0.02000 recovered exactly, \ensuremath{\sigma}\down{\delta{}}=9.7\ensuremath{\times}10\up{-4}, \ensuremath{\Delta}lnL=338.4 \ensuremath{\implies} GR excluded \ensuremath{\sim}26\ensuremath{\sigma} (robust for Kerr), GR control shows no false positive. \textbf{v1.35 a-\ensuremath{\delta} joint analysis (three major numerical repairs)}: (1) \textbf{undersampling} (v1.32 dt/carrier-period=34.6 \ensuremath{\to} aliased \ensuremath{\delta}-likelihood comb) fixed by analytic inspiral r(t)=(r\down{0}\up{4}\ensuremath{-}kt)\up{1/4} + Nyquist sampling dt=400s; (2) \textbf{orbital-phase comb} (a enters fast phase 2\ensuremath{\lambda}, 4.3\ensuremath{\times}10\up{5} rad accumulated in 2yr) fixed by \ensuremath{\varphi}0-profiled likelihood |\ensuremath{\Sigma}d\ensuremath{\cdot}e\up{-i\Phi{}}|; (3) \textbf{Lense--Thirring precession added} (\ensuremath{\omega}\down{peri}=(1+\ensuremath{\delta})\ensuremath{\omega}\down{GR}+\ensuremath{\omega}\down{LT}(a); LT\ensuremath{\propto}x\up{3/2} vs 1PN\ensuremath{\propto}x breaks the \ensuremath{\delta}-a degeneracy). \textbf{Core results}: (i) \textbf{\ensuremath{\sigma}\down{\delta{}} corrected by 600\ensuremath{\times}}: \ensuremath{\sigma}\down{\delta{}}(1D, \ensuremath{\varphi}0-profile, a fixed)=1.63\ensuremath{\times}10\up{-6} (2yr) / 7.37\ensuremath{\times}10\up{-7} (4yr); (ii) degeneracy small with a free: \ensuremath{\sigma}\down{a}=1.5\ensuremath{\times}10\up{-5} (LT phase pins a), \ensuremath{\sigma}\down{\delta{}}(3D marginal)\ensuremath{\approx}3.8\ensuremath{\times}10\up{-6} (2.4\ensuremath{\times} inflation); (iii) \textbf{\ensuremath{\delta}=2\% detection robust: \ensuremath{\Delta}lnL(GR)=357 (2yr) / 396 (4yr) \ensuremath{\implies} 27\ensuremath{\sigma}/28\ensuremath{\sigma} (a-\ensuremath{\varphi}0 fully free)} --- the earlier "167\ensuremath{\sigma}" was aliased-likelihood inflation, true value 28\ensuremath{\sigma}; (iv) \textbf{main degeneracy identified: \ensuremath{\delta}-\ensuremath{\varphi}0 (\ensuremath{\rho}\ensuremath{\approx}\ensuremath{-}0.85, orbital phase), not \ensuremath{\delta}-a}; (v) GR control \ensuremath{\delta}\down{hat}\ensuremath{\approx}0, \ensuremath{\Delta}lnL\ensuremath{\approx}0 (no false positive). Honest caveat: 3D MCMC mixing difficulty (very narrow \ensuremath{\sigma}\down{a}), \ensuremath{\sigma}\down{\delta{}} is local curvature (robust criterion is \ensuremath{\Delta}lnL), circular orbit / quadrupole / non-Teukolsky --- \textbf{the S-E1 prediction remains measurable at 27-28\ensuremath{\sigma} with spin and initial phase fully free}.

\textbf{\mtag{v2.4 supplement: comparison with the exponential-metric literature (referee M4)}}\mtag{Literature comparison}: the framework metric g\down{\mu{}\nu{}}=diag(\ensuremath{-}e\up{-2x}, e\up{2x}, r\up{2}e\up{2x}, r\up{2}sin\up{2}\ensuremath{\theta}e\up{2x}) is \textbf{the Yilmaz--Rosen exponential metric} (first proposed by Papapetrou 1954; Yilmaz 1958/1971; Rosen 1962; also Vaidya 1954 / Coleman 1971 / Harvey 1972; recent: Makukov \& Mychelkin 2018, Boonserm et al. 2021). Literature-known properties vs framework: (i) \textbf{PPN \ensuremath{\alpha}=\ensuremath{\beta}=\ensuremath{\gamma}=1} (indistinguishable from GR at 1PN) --- established in the literature (Yilmaz tradition); the framework provides a statistical-mechanical source for its emergence (\S{}4.2 flux Boltzmann) rather than assuming it; (ii) \textbf{2PN g\down{rr} bifurcation 33\% / 3PN g\down{00} bifurcation 11\% = literature-known higher-order PPN properties} (Yilmaz \ensuremath{\delta}=1 vs GR \ensuremath{\delta}=3/4; analytic expansion verified: g\down{rr}=e\up{2x}=1+2x+2x\up{2}+\ensuremath{\cdots} vs GR (1+x/2)\up{4}=1+2x+1.5x\up{2}+\ensuremath{\cdots} \ensuremath{\implies} 33\%; g\down{00}=\ensuremath{-}e\up{-2x} vs GR \ensuremath{-}[(1\ensuremath{-}x/2)/(1+x/2)]\up{2} \ensuremath{\implies} 11\% at 3PN) --- \textbf{numerically identical, but the framework's contribution is the emergent derivation of the mechanism and the discrimination scheme} (EMRI \ensuremath{\delta}=0.30\ensuremath{\cdot}x\down{p} scaling law, S-E1--E4); (iii) \textbf{no event horizon} (g\down{00} never vanishes) --- literature-known, consistent with the framework's v1.13 withdrawal of the geometric horizon and the horizon = storm-overload equipotential surface; (iv) \textbf{the exponential metric is not an EFE vacuum solution} (R\down{\mu{}\nu{}}=\ensuremath{-}\ensuremath{\tfrac12}\ensuremath{\nabla}\ensuremath{\varphi}\ensuremath{\nabla}\ensuremath{\varphi}, requiring a negative-kinetic-energy ghost scalar field) --- the framework's "mass-shielded flux producing negative pressure" (\S{}4.1) is proposed as the physical counterpart of this source term, to be rigorized in future work; (v) \textbf{traversable wormhole} (Boonserm et al. 2021: YEM = traversable wormhole) --- corresponds to the framework's "cross-branch ER channel / quantum line" (\S{}4.6). \textbf{Positioning among emergent-gravity literature}: \S{}4.7 horizon-thermodynamics path = Jacobson 1995 method (Clausius \ensuremath{\implies} Einstein); \S{}4.2 flux-Boltzmann mechanism differs from Verlinde 2011 entropic gravity and Padmanabhan 2010 equipartition (different mechanism, compatible conclusions); \S{}3 causal sets and \S{}5.8 string-nets are standard tools (Bombelli 1987 / Malament 1977 / Wen 2003 / Kitaev 2003). \mtag{Innovation re-assessment} the metric and bifurcations are literature-known; \textbf{the framework's additions}: Boltzmann-mechanism emergence of the metric, unified flux-occlusion language, EMRI discrimination scaling law, pinned equation of state, packing theorem package, string-net rigorization series.

\subsection{4.6 Strong-field mechanism and black holes (revised)}

\textbf{\mtag{v1.1 important revision}}: An early version (v1.0) extrapolated the weak-field Boltzmann metric g\down{00}=\ensuremath{-}e\up{-2GM/c^{2}r} to strong fields, predicting a "photon sphere 1GM vs 3GM" conflict with EHT. On review, \textbf{that extrapolation is outside the framework's domain of validity} (the thermal-equilibrium proof already established that a pure gravitational potential has no global thermal equilibrium; e\up{-\Phi{}/c^{2}} is valid only in the weak field / locally / with cutoff) --- \textbf{formally withdrawn}.

\textbf{\mtag{v1.14 mathematical exclusion of strong-field metrics}}: Attempts to superimpose flux-saturation cutoffs (soft-core / linear / logarithmic) on the Boltzmann weak-field form (\ensuremath{\beta}=1 surviving) and test against GR strong-field facts (horizon 2GM, photon sphere 3GM, shadow 5.196GM) --- \textbf{all failed}: (1) the pure-Boltzmann photon sphere is analytically proven at 1GM (V=e\up{-2x}x\up{2}, dV/dx=0 \ensuremath{\implies} x=1); (2) saturated corrections give photon spheres 2.16--4.56GM (deviations 9--52\%); (3) the saturation factor leaks into the weak field (x=0.01), destroying \ensuremath{\beta}=1. \textbf{\ensuremath{\implies} All Boltzmann-based saturated strong-field metrics are mathematically excluded; the framework has no strong-field metric formula.} Strong field = storm-overload physical mechanism (horizon = pressure-critical equipotential surface, not a g\down{00}=0 geometric singularity); no conflict with EHT (no strong-field metric commitment).

\textbf{\mtag{v1.15 storm-overload observational effects (one-parameter family \ensuremath{\alpha})}}: The framework's strong-field observational content is encoded in a minimal-commitment one-parameter family --- flux-vanishing radius r\down{h} = \ensuremath{\alpha}\ensuremath{\cdot}2GM/c\up{2} (\ensuremath{\alpha}=1 is the GR limit), predicting: shadow b=5.196\ensuremath{\alpha}\ensuremath{\cdot}GM/c\up{2}, QNM f=f\down{GR}/\ensuremath{\alpha}, firewall spread angle \ensuremath{\delta}/R=\ensuremath{\sqrt{2l\down{P}/R}}. \textbf{EHT constrains \ensuremath{\alpha} = 1.00\ensuremath{\pm}0.06 (M87*/Sgr A* shadow 5.2 GM) --- the framework is fully compatible with black-hole shadow observations}; future ET/CE QNM precision \ensuremath{\sim}1\% \ensuremath{\implies} \ensuremath{\alpha} constrained to \ensuremath{\pm}0.01 (first precise verdict); \textbf{framework-unique signature: firewall spread \ensuremath{\delta}/R (primordial black holes \ensuremath{\sim}10\up{-9} measurable; GR shadows are sharp)}; overload region without singularity (storm suppression, Assumption H-7) vs GR singularity. Prediction list S1--S5 in \S{}9.

\textbf{\mtag{v1.16 micro-origin of \ensuremath{\alpha} = entropy consistency (theoretical necessity, not fitting)}}: \ensuremath{\alpha} is forced to 1 by the framework's thermodynamic self-consistency --- (1) the framework's Kakeya microstate counting S=A/(4l\down{P}\up{2}) has the same 1/4 coefficient as Bekenstein--Hawking \ensuremath{\implies} entropy consistency \ensuremath{\implies} area consistency A\down{overload}=A\down{BH} \ensuremath{\implies} \ensuremath{\alpha}\up{2}=1 \ensuremath{\implies} \textbf{\ensuremath{\alpha}=1 (unique solution)}; (2) Jacobson-path cross-check: \ensuremath{\alpha}\ensuremath{\neq}1 \ensuremath{\implies} G\down{eff}=G/\ensuremath{\alpha}\up{2}\ensuremath{\neq}G, conflicting with PPN/gravitational-wave observations (1\ensuremath{\pm}10\up{-5}); (3) numerically, at \ensuremath{\alpha}=1 the solar-mass black hole S=1.05\ensuremath{\times}10\up{77} nats exactly equals standard BH entropy. \textbf{\ensuremath{\implies} EHT observation \ensuremath{\alpha}=1.00\ensuremath{\pm}0.06 is a direct verification of the framework's self-consistency} --- the overload-region radius is locked to the Schwarzschild radius by entropy consistency.

\textbf{\mtag{v1.17 pressure-critical area locking (M\up{2}/A Planck criticality)}}: "A black hole depends only on M/A" in GR gives M/A=c\up{4}/(16\ensuremath{\pi}GM) (M-dependent); \textbf{the squared invariant M\up{2}/A = c\up{4}/(16\ensuremath{\pi}G\up{2}) \ensuremath{\approx} 3.6\ensuremath{\times}10\up{52} kg\up{2}/m\up{2} (Planck constant, universal criticality)} \ensuremath{\implies} A=16\ensuremath{\pi}G\up{2}M\up{2}/c\up{4} \ensuremath{\implies} r\down{h}=2GM/c\up{2} (exact; verified for Sun/M87*/Sgr A*) --- storm-suppression strength \ensuremath{\propto} M\up{2} (gravitational self-binding), surface resistance \ensuremath{\propto} A; the suppression/resistance ratio reaching the Planck constant = the Schwarzschild surface. Candidate alternatives excluded: gravitational self-pressure GM\up{2}/(4\ensuremath{\pi}r\up{4}) (r\down{h} off by 98\%), (M/A)c\up{2}/r (r\down{h}\ensuremath{\propto}M\up{1/3}). \textbf{\ensuremath{\alpha}=1 acquires a dual origin: entropy consistency (v1.16 thermodynamic path) + M\up{2}/A Planck criticality (mechanical path) --- area locking upgraded from a thermodynamic requirement to a physical equilibrium surface of storm pressure.}

\textbf{\mtag{v1.18 + v1.24 overload-region fluid model (positive-matter "paste" equation of state)}}: \mtag{Assumption H-7} "the maximum-entropy state of positive-matter paste\ldots the storm pressure is finite (Planck-scale limited)" \ensuremath{\implies} equation of state \textbf{P(\ensuremath{\rho}) = P\down{P}\ensuremath{\cdot}tanh(\ensuremath{\rho}/\ensuremath{\rho}\down{P})} (Planck pinning: linear region recovers P\ensuremath{\approx}c\up{2}\ensuremath{\rho}, saturation region pins pressure at P\down{P}\ensuremath{\approx}4.6\ensuremath{\times}10\up{113} Pa --- \textbf{the quantitative expression of singularity avoidance}). \textbf{v1.24 micro-derivation (from phenomenological assumption to emergence)} --- three mechanisms, all from existing framework premises, no new assumptions:
\begin{itemize}
\item \textbf{Mechanism (1) Planck filling ratio}: \ensuremath{\rho}/\ensuremath{\rho}\down{P} with \ensuremath{\rho}\down{P} = m\down{P}/l\down{P}\up{3} (one debris mass per Planck volume) --- geometric scaling, not a fitting parameter;
\item \textbf{Mechanism (2) tanh emergence}: overload-region pressure = debris--storm bidirectional net momentum flow (reflection +J\down{max} / transmission \ensuremath{-}J\down{max}); the net flow obeys logistic saturation dJ/d\ensuremath{\rho}=(1/\ensuremath{\rho}\down{P} J\down{max})(J\down{max}\up{2}\ensuremath{-}J\up{2}) \ensuremath{\implies} \textbf{J=J\down{max}\ensuremath{\cdot}tanh(\ensuremath{\rho}/\ensuremath{\rho}\down{P})} (numerical integration vs analytic deviation 8\ensuremath{\times}10\up{-4}) --- two-channel competition \ensuremath{\implies} hyperbolic tangent;
\item \textbf{Mechanism (3) Planck pressure}: P\down{P} = E\down{P}/l\down{P}\up{3} = single-collision momentum p\down{P}=\ensuremath{\hbar}/l\down{P} \ensuremath{\times} maximum collision frequency \ensuremath{\nu}\down{P}=c/l\down{P} (Axiom-1 event-flow upper bound; \mtag{Axiom-1 corollary} "annihilation has a maximum rate") / unit area l\down{P}\up{2};
\item \textbf{By-products}: merged mass m\down{eff} = m\down{P}\ensuremath{\cdot}(\ensuremath{\rho}/\ensuremath{\rho}\down{P})/tanh(\ensuremath{\rho}/\ensuremath{\rho}\down{P}) (independent debris \ensuremath{\to} storm-pressed clumps, micro-picture of "paste"); \textbf{new prediction: sound-speed profile c\down{s}\up{2} = \ensuremath{\partial}P/\ensuremath{\partial}\ensuremath{\rho} = c\up{2}\ensuremath{\cdot}sech\up{2}(\ensuremath{\rho}/\ensuremath{\rho}\down{P})} --- c\down{s}=c for \ensuremath{\rho}\ensuremath{\ll}\ensuremath{\rho}\down{P} (relativistic maximal causality), c\down{s}\ensuremath{\to}0 for \ensuremath{\rho}\ensuremath{\gg}\ensuremath{\rho}\down{P} (\textbf{no elastic waves in the pinned region = dynamical expression of structureless paste}, linking Kakeya-set information \ensuremath{\approx} 0 and "no characteristic radiation from the overload region"); \textbf{energy condition}: P \ensuremath{\leq} \ensuremath{\rho}c\up{2} everywhere (deeply sub-causal P/\ensuremath{\rho}c\up{2}\ensuremath{\to}\ensuremath{\rho}\down{P}/\ensuremath{\rho}\ensuremath{\to}0 for \ensuremath{\rho}\ensuremath{\gg}\ensuremath{\rho}\down{P}), \ensuremath{\partial}P/\ensuremath{\partial}\ensuremath{\rho}\ensuremath{\geq}0 \ensuremath{\implies} no thermodynamic instability \ensuremath{\implies} no collapse (contrast GR singularity);
\end{itemize}
Numerics: macroscopic black-hole overload-region average density far below the Planck density (Sun \ensuremath{\rho}\down{avg}=1.8\ensuremath{\times}10\up{19} vs \ensuremath{\rho}\down{P}=5.2\ensuremath{\times}10\up{96} kg/m\up{3}, 78 orders of magnitude apart; P\ensuremath{\approx}10\up{-6} P\down{P}, c\down{s}\ensuremath{\approx}c); only Planck-mass black holes approach \ensuremath{\rho}\down{P} (\ensuremath{\rho}\down{avg}=0.03\ensuremath{\rho}\down{P}, P=0.030P\down{P}, c\down{s}=0.9996c) --- \textbf{singularity naturally eliminated}; accretion behavior: matter is crushed by the storm into positive-matter debris (not GR free-fall \ensuremath{\to} singularity); radiation: the overload region = maximum-entropy structureless \ensuremath{\implies} no characteristic radiation (radiation comes from horizon-surface saturated modes + Hawking radiation); prediction: a subset of Hawking radiation is entangled with infalling matter (information-flip signature, S4 observational window).

\textbf{\mtag{v1.25 overload-region radial structure and radiative transfer (items (2)(3) closed)}}: \textbf{Isothermal-sphere exclusion (derivation + numerics)} --- assuming a Newtonian gravitational gradient inside the overload region (linear-region P\ensuremath{\approx}c\up{2}\ensuremath{\rho} hydrostatic equilibrium) \ensuremath{\implies} singular isothermal sphere \ensuremath{\rho}(r)=c\up{2}/(2\ensuremath{\pi}Gr\up{2}), M(r)=2c\up{2}r/G \ensuremath{\implies} \textbf{M(r\down{h})=4M constant overweight} (verified 4.0\ensuremath{\times}M for six masses: Sun/Sgr A*/M87*/primordial 10\up{11}/10\up{13} kg/Planck mass; the singular isothermal sphere has no scale-free self-similar solution) \ensuremath{\implies} contradicts total-mass conservation; \textbf{framework-mechanism confirmation}: horizon = flux=0 equipotential surface + storm-suppression region (seven-loop ring 1) \ensuremath{\implies} no net flux gradient inside \ensuremath{\implies} no gravity (gravity = flux-occlusion deficit vanishing) \ensuremath{\implies} dP/dr=0 \ensuremath{\implies} \textbf{uniform paste} \ensuremath{\bar{\rho}} = 3c\up{6}/(32\ensuremath{\pi}G\up{3}M\up{2}) (locked by total mass + volume; exactly consistent with v1.18 \ensuremath{\rho}\down{avg}: Sun 1.842\ensuremath{\times}10\up{19} kg/m\up{3} \checkmark, Sgr A* 3.94\ensuremath{\times}10\up{6} \checkmark, M87* 1.73 \checkmark); P\ensuremath{\approx}c\up{2}\ensuremath{\bar{\rho}} (linear region), c\down{s}\ensuremath{\approx}c --- \textbf{"uniform = maximum-entropy structureless" as dynamical expression} (no internal gradient \ensuremath{\implies} no internal structure \ensuremath{\implies} no characteristic radiation). \textbf{Radiative transfer (quantitative)}: debris cross-section \ensuremath{\sigma}=l\down{P}\up{2} (minimum Planck-cell cross-section assumption) \ensuremath{\implies} optical depth \ensuremath{\tau}=(\ensuremath{\bar{\rho}}/m\down{P})l\down{P}\up{2}r\down{h}=3c\up{4}l\down{P}\up{2}/(16\ensuremath{\pi}G\up{2}m\down{P} M) --- \textbf{macroscopic black holes \ensuremath{\tau}\ensuremath{\sim}10\up{-40}\ensuremath{\sim}10\up{-49} (optically transparent: the quantitative reason the overload region neither absorbs nor emits; infalling radiation passes straight through)}, primordial 10\up{11} kg \ensuremath{\tau}=1.3\ensuremath{\times}10\up{-20}, \textbf{Planck-mass black hole \ensuremath{\tau}=0.0597 (T=0.942, first "visible" critical window of the overload region)}.

\textbf{\mtag{v1.31 accretion-disk radiative transfer (item (3) fully closed)}}\mtag{Citation + Numerical}: the standard thin disk (Shakura--Sunyaev multi-temperature blackbody T(r)=[3GM\ensuremath{\dot{M}}/(8\ensuremath{\pi}\ensuremath{\sigma}r\up{3})(1\ensuremath{-}\ensuremath{\sqrt{r\down{in}/r}})]\up{1/4}, r\down{in}=6GM/c\up{2}) is the framework's accretion-disk \textbf{bulk --- completely consistent with GR} (outer-region geometry indistinguishable at 1PN, \S{}4.5b; compatible with EHT): for M=10\up{6} M\ensuremath{\odot}, \ensuremath{\dot{M}}=0.01\ensuremath{\dot{M}}\down{Edd}, T(2r\down{in})=\textbf{1.2\ensuremath{\times}10\up{5} K}, spectral peak \ensuremath{\nu}\ensuremath{\approx}6.7\ensuremath{\times}10\up{15} Hz (X-rays), L\down{acc}=1.05\ensuremath{\times}10\up{35} W\ensuremath{\approx}0.01L\down{Edd}; \textbf{overload-region radiative transfer}: \ensuremath{\tau}\ensuremath{\sim}10\up{-40} (Sun/Sgr A*/M87*), primordial 10\up{11} kg \ensuremath{\tau}=1.3\ensuremath{\times}10\up{-20} --- \textbf{infalling disk radiation passes through the overload region without secondary absorption} (quantitative continuation of v1.25; the overload region does not alter disk radiation); \textbf{framework-unique differences at the horizon surface}: firewall emission layer (F4) non-thermal correction \ensuremath{\sim}\ensuremath{\sqrt{l\down{P}/R}} (Sun 7.4\ensuremath{\times}10\up{-20}, \textbf{primordial 10\up{11} kg 3.3\ensuremath{\times}10\up{-10}}, Planck 0.71) + \textbf{Hawking-radiation--infalling-matter entanglement}\mtag{Framework prediction S4} --- the final-stage \ensuremath{\gamma}-ray window of primordial black holes (10\up{11}--10\up{13} kg, T\down{H}=10\up{10}--10\up{13} K) contains framework-unique signatures. \mtag{Honest caveat} thin disk is a Newtonian approximation (no GR redshift/light bending), in-firewall transport not modeled, Kerr disk not treated.

\textbf{\mtag{v1.21 shadow of the isotropic exponential metric compatible with EHT (correcting v1.14)}}: v1.14's "Boltzmann photon sphere 1GM conflicts with EHT" arose from the \textbf{standard coordinate convention} (g\down{\theta{}\theta{}}=r\up{2}, no angular bending; its shadow b\down{min}=e=2.718 GM/c\up{2} indeed conflicts); but the \textbf{isotropic exponential metric} identified in \S{}4.5b (the natural choice of Axiom-1 isotropy) gives the coordinate-invariant impact parameter b\down{min}=\ensuremath{\sqrt{g\down{\varphi{}\varphi{}}/|g\down{00}|}}:

\begin{table}[htbp]
\centering
\small
\begin{tabularx}{\textwidth}{@{}lXXX@{}}
\toprule
Metric & Photon sphere r\down{ph} & Shadow b\down{min} & vs EHT (5.2\ensuremath{\pm}0.3 GM/c\up{2}) \\
\midrule
\textbf{Framework isotropic exponential metric} & \textbf{2.00 GM/c\up{2}} & \textbf{5.436 GM/c\up{2} (=2e)} & \textbf{\checkmark compatible (+4.55\% < 5.8\% error)} \\
GR Schwarzschild & 3.00 GM/c\up{2} & 5.196 GM/c\up{2} (=3\ensuremath{\sqrt{3}}) & baseline \\
{[contrast] standard coordinates g\down{\theta{}\theta{}}=r\up{2}} & 1.00 GM/c\up{2} & 2.718 GM/c\up{2} (=e) & \ensuremath{\times} conflict (v1.14) \\
\bottomrule
\end{tabularx}
\end{table}

\textbf{The framework's strong field is upgraded from "no metric commitment (v1.14)" to "geometric predictions compatible with observation (v1.21)"} --- compatible with EHT without saturation corrections (v1.14 excluded "standard-coordinate Boltzmann + saturation corrections"); incidental consistency: the exponential-metric photon sphere 2GM/c\up{2} exactly coincides with the framework horizon (entropy-consistency \ensuremath{\alpha}=1, v1.16). The storm overload remains the internal mechanism narrative; geometric and mechanistic sides are complementary. Honest caveat: the shadow computation is an analytic continuation of the weak-field metric to strong field (x\ensuremath{\sim}0.5 beyond the thermal-equilibrium cutoff r\down{c}=GM/c\up{2}, B2).

The framework's strong-field content is given by the seven-loop mechanism chain (ontological level, not metric-equation modification):

\begin{verbatim}
Black hole = storm-suppression region -> horizon = pressure-critical equipotential surface (flux=0)
  -> interior = 3D Kakeya set (overload region, information measure ~ 0)
  -> quantum line = Kakeya segment (ER=EPR)
  -> time-arrow flip (information flips branch, not lost)
  -> freeze = ontology (positive-branch projection source halts; gentle firewall preserves equivalence principle)
  -> residue shedding = Hawking radiation (Page curve = residue exhaustion)
\end{verbatim}

\begin{itemize}
\item \textbf{Strong-field metric form not given} (interior is an overload region, not a geometric region) --- the framework makes no strong-field metric commitment, no conflict with EHT, no exclusion risk;
\item Strong-field predictions in \S{}9, table B1--B5 (extreme-Kerr angular distribution, primordial black-hole volume correction, Hawking-radiation--infalling-matter entanglement, ER-bridge packing bound, non-thermal correction 1/S).
\end{itemize}

\textbf{Firewall thickness verdict (F1--F4)}\mtag{Framework assumption: horizon-surface saturated modes, thickness \ensuremath{\alpha}=0 or 1/2 undecided} --- \textbf{verdict \ensuremath{\delta} \ensuremath{\sim} \ensuremath{\sqrt{R\cdot{}l\down{P}}} (\ensuremath{\alpha}=1/2, horizon fleece)}:
\begin{itemize}
\item \textbf{F1 (diffusion mechanism, strongest)}: Axiom 1 (isotropic random event flow) \ensuremath{\implies} flux recovers from the horizon (=0) by random diffusion \ensuremath{\implies} front width \ensuremath{\delta} \ensuremath{\sim} \ensuremath{\sqrt{N}}\ensuremath{\cdot}l\down{P} (N\ensuremath{\sim}R/l\down{P}) \ensuremath{\implies} \ensuremath{\delta} \ensuremath{\sim} \ensuremath{\sqrt{R\cdot{}l\down{P}}} (numerically \ensuremath{\sigma}/\ensuremath{\sqrt{N}} \ensuremath{\approx} 1.0);
\item \textbf{F1\ensuremath{^\prime} (exact diffusion-equation solution)}: random-walk continuum limit \ensuremath{\implies} diffusion equation \ensuremath{\partial}\ensuremath{\rho}/\ensuremath{\partial}t = D\ensuremath{\partial}\up{2}\ensuremath{\rho}/\ensuremath{\partial}x\up{2} (D = l\down{P}\ensuremath{\cdot}c/2) \ensuremath{\implies} horizon boundary \ensuremath{\rho}(0,t)=0 \ensuremath{\implies} erf solution \ensuremath{\rho}=\ensuremath{\rho}\ensuremath{\infty}\ensuremath{\cdot}erf(x/\ensuremath{\sqrt{4Dt}}) \ensuremath{\implies} recovery time t\down{rec}=R/c \ensuremath{\implies} \textbf{\ensuremath{\delta} = \ensuremath{\sqrt{2}}\ensuremath{\cdot}\ensuremath{\sqrt{R\cdot{}l\down{P}}}} (exact 1D coefficient; zero numerical deviation);
\item \textbf{F1\ensuremath{^{\prime\prime}} (3D spherically symmetric exact solution)}: the u = r\ensuremath{\rho} transformation exactly absorbs the 2/r term \ensuremath{\implies} 3D spherical symmetry reduces to 1D diffusion \ensuremath{\implies} mirror-method Green's function \ensuremath{\implies} \textbf{closed-form solution \ensuremath{\rho} = \ensuremath{\rho}\ensuremath{\infty}\ensuremath{\cdot}[R\ensuremath{\cdot}erf((r\ensuremath{-}R)/\ensuremath{\sqrt{4Dt}}) + (r\ensuremath{-}R)]/r} (horizon \ensuremath{\rho}(R)=0, far field \ensuremath{\rho}\ensuremath{\to}\ensuremath{\rho}\ensuremath{\infty}, x\ensuremath{\ll}R reduces to 1D --- all three limits pass) \ensuremath{\implies} characteristic width \textbf{\ensuremath{\delta} = \ensuremath{\sqrt{2}}\ensuremath{\cdot}\ensuremath{\sqrt{R\cdot{}l\down{P}}}}, half-value thickness \ensuremath{\delta}\down{1}/\down{2} = 0.6745\ensuremath{\cdot}\ensuremath{\sqrt{R\cdot{}l\down{P}}}; \textbf{the validity of the 1D approximation is rigorously proven} (geometric correction O(\ensuremath{\varepsilon})=\ensuremath{\sqrt{2l\down{P}/R}} \ensuremath{\sim} 10\up{-19}); dimensionless diffusion constant \textbf{D\down{eff} = 1/4 analytic identity}; numerical verification: finite-difference deviation 0.0024\%, half-value-point deviation 6\ensuremath{\times}10\up{-6};
\item \textbf{F1\ensuremath{^{\prime\prime\prime}} (finite-R exact solution)}: introducing an outer boundary L (cosmological horizon / nearby body) \ensuremath{\implies} separation of variables, spherical-Bessel series u = \ensuremath{\rho}\ensuremath{\infty}(r\ensuremath{-}R)L/(L\ensuremath{-}R) + \ensuremath{\Sigma}c\down{n} sin(k\down{n}(r\ensuremath{-}R))e\up{-Dk\down{n}^{2}t} (k\down{n}=n\ensuremath{\pi}/(L\ensuremath{-}R)) \ensuremath{\implies} all numerical checks pass (convergence N\ensuremath{\geq}200, initial/steady limits exact) \ensuremath{\implies} \textbf{half-infinite approximation error quantified: deviation <0.01\% for L/R\ensuremath{\geq}1.1} (only failure = extremely close binary merger); \textbf{important physical finding: the steady-state solution \ensuremath{\rho}\down{ss}(r) = \ensuremath{\rho}\ensuremath{\infty}(1\ensuremath{-}R/r) --- the flux deficit \ensuremath{\propto} 1/r, of the same form as the Newtonian potential} (P0 derived the 1/r\up{2} force from solid-angle integration; this batch independently derives the 1/r potential from the diffusion steady state --- two paths converge, independently supporting "gravity = flux occlusion");
\item \textbf{F2 (mode energy scale)}: \ensuremath{\delta}=l\down{P} \ensuremath{\implies} scale \ensuremath{\sim} E\down{P} (AMPS risk); \ensuremath{\delta}=\ensuremath{\sqrt{R\cdot{}l\down{P}}} \ensuremath{\implies} scale \ensuremath{\sim} E\down{P}\ensuremath{\cdot}\ensuremath{\sqrt{l\down{P}/R}} (solar black hole 7.4\ensuremath{\times}10\up{-20} E\down{P}, genuinely gentle);
\item \textbf{F3 (observational compatibility)}: tidal forces of both thicknesses are imperceptible (equivalence principle preserved); QNM correction \ensuremath{\delta}/R \ensuremath{\sim} 10\up{-19} unmeasurable --- the verdict must come from theoretical mechanism;
\item \textbf{F4 (emission-layer prediction)}: firewall = Hawking-radiation emission layer\mtag{Framework assumption}\ensuremath{\implies} non-thermal correction \ensuremath{\sim} \ensuremath{\sqrt{l\down{P}/R}} --- primordial black holes \ensuremath{\sim}10\up{-9}\ensuremath{\sim}10\up{-4} measurable window (complementary to B5).
\item \mtag{Honest caveat} the early record "primordial black hole R\ensuremath{\sim}10\up{8} l\down{P}" differs from the standard formula R\down{s}(10\up{11}kg)\ensuremath{\approx}9\ensuremath{\times}10\up{18} l\down{P} by \ensuremath{\sim}10\up{11} (likely a typographical error; the standard formula is adopted).
\end{itemize}

\subsection{4.7 Emergence of the Einstein equation (Jacobson path)}

\mtag{Citation + Derivation} Rindler horizon + entropy area law (\S{}4.8) + Clausius relation \ensuremath{\delta}Q = T dS \ensuremath{\implies} \textbf{the Einstein equation emerges as a horizon-thermodynamics identity}:

\[G_{\mu\nu} = 8\pi G\,T_{\mu\nu}\]

Framework connection: entropy area law (\S{}4.8), temperature = flux gradient (Unruh), energy flow = flux occlusion, cutoff = Planck (naturally from Axiom 1). \mtag{To be verified} full closure still requires a rigorous dynamical derivation "event flow \ensuremath{\to} entropy area law".

\subsection{4.8 Black-hole entropy S = A/4: flux-tube coherent-block microstate counting}

\textbf{\mtag{Microstate-counting proposition}}: the horizon = flux-saturated surface = dense packing of flux-tube endpoints. The surface has N = A/l\down{P}\up{2} plaquettes; every 4 plaquettes form 1 coherent unit (Clifford 4 components: two time directions \ensuremath{\times} two chirality directions), 1 bit per unit:

\[S = \frac{A}{4l_P^2}\ln 2 \quad\text{(bits)} = \frac{A}{4l_P^2}\quad\text{(nats)},\qquad k = 2^{1/4} \approx 1.189\]

\textbf{Numerical verification}: c\ensuremath{\cdot}t\down{P} = l\down{P} (1 event per plaquette per Planck time); solar-mass black hole S = 1.05\ensuremath{\times}10\up{77} \checkmark.

\textbf{Mechanism of appearance of G, \ensuremath{\hbar}}\mtag{Derivation}: S = A\ensuremath{\cdot}c\up{3}/(4G\ensuremath{\hbar}) --- G, \ensuremath{\hbar} appear as composites of the Planck area l\down{P}\up{2} = G\ensuremath{\hbar}/c\up{3} (not fundamental constants; derived quantities of the event-flow minimal scale).

\textbf{Candidate comparison}: Boltzmann factor (\ensuremath{\times}), LQG area spectrum (different origin, but same ln 2), Bekenstein 1-photon argument (weak bound 8\ensuremath{\pi}, \ensuremath{\times}) --- the flux-tube coherent-block model is the unique exact match. The \textbf{geometric origin of the 1/4 coefficient} (Kakeya packing, k\ensuremath{\geq}4 lower-bound theorem) is in \S{}4.10. \mtag{To be verified} the dynamical origin of k=2\up{1/4} still requires "flux--coupling-mode counting".

\textbf{\mtag{v2.8 supplement: independent derivation of the k=2\up{1/4} factor (referee M7 confirmation-bias objection)}}:
(1) \textbf{identity positioning of k} (`\down{m7}\down{factor}.py`): k=2\up{1/4} is an \textbf{equivalent notation} for "1 bit per 4 plaquettes", not an independent fitting parameter --- \ensuremath{\Omega}=2\up{N/4} \ensuremath{\implies} S\down{bit}=N/4=A/(4l\down{P}\up{2}), k=\ensuremath{\Omega}\up{1/N}=2\up{1/4} (numerically verified) --- the real question is the origin of 1/4;
(2) \textbf{sampling-theorem independent argument}: direction cluster k\ensuremath{\geq}4 (Theorem 1) + Tammes regular tetrahedron \ensuremath{\implies} 4 independent direction modes per plaquette; Nyquist sampling theorem: 4 modes require \ensuremath{\geq}4 sampling points (numerical rank check: 4 modes \ensuremath{\times} 2/3 points rank 2/3 undersampled, \ensuremath{\times}4 points rank 4 full) \ensuremath{\implies} \textbf{coherent block \ensuremath{\geq} 4l\down{P}\up{2} determined independently by general information theory} (not contrived for S=A/4); maximal compression of the 4 internal degrees of freedom \ensuremath{\implies} 1 bit/block \ensuremath{\implies} entropy density 1/(4l\down{P}\up{2});
(3) \textbf{three independent paths cross-validate 1/4}: packing geometry (Proposition 5, mathematically rigorous) / microstate counting (Clifford 4 components) / \textbf{entropy consistency \ensuremath{\alpha}=1} (solar-mass black hole S=1.049\ensuremath{\times}10\up{77} nats exactly equal to Bekenstein--Hawking, 0.00\% deviation) --- if it were "contrived", the probability that three independent mechanisms simultaneously give exactly 1/4 is extremely low;
(4) \textbf{intermediate quantities falsifiable}: coherent-block size 4l\down{P}\up{2}, entropy density 1/4 bit/plaquette, direction cluster k=4 all independently measurable (quantum simulation / topological entanglement entropy / packing mathematics) --- intermediate quantities falsifiable \ensuremath{\implies} confirmation-bias risk substantially mitigated. Honest caveat: the "dynamical uniqueness of 4" (why the interior space has 4 components) lies at the framework-assumption level (same level as the vortex winding-mode number).

\subsection{4.9 Holographic concentration mechanism (why the area law)}

\mtag{Derivation} In 3+3-dimensional space--time, the observer = local soliton on the present plane\mtag{Framework assumption: observer = present-plane soliton}:
\begin{itemize}
\item Lemma 1: accessible information = in-plane 4-direction connections (2 out-of-plane directions curled)\mtag{Framework assumption};
\item Lemma 2: bulk information = surface projection (bulk redundant, surface primary)\mtag{Framework assumption};
\item \textbf{Proposition (holographic concentration)}: S\down{max}(R) = (A/4l\down{P}\up{2})\ensuremath{\cdot}ln 2 --- observer limitation is the origin of information concentration (not volume compression, but bulk information inaccessible).
\end{itemize}

\textbf{Observer--information proposition chain (O1--O4)}: rigorous version of Proposition 6 Lemma A-1 + mutual-information framework of holographic concentration\mtag{Framework assumption}:
\begin{itemize}
\item \textbf{O1 (observer structure)}: observer = present-plane soliton; readable degrees of freedom = in-plane 4-direction connections + interior-space Clifford 4 components\mtag{Framework assumption};
\item \textbf{O2 (accessible information)}: S\down{access}(R) = I(R;\ensuremath{\Sigma}) (mutual information); bulk information redundant, surface information primary\mtag{Framework corollary};
\item \textbf{O3 (uniform prior derived)}: observer geometry (no information channel along out-of-plane directions) \ensuremath{\implies} uniform prior over external spatial directions --- \textbf{Proposition 6 Lemma A-1 upgraded from assumption to corollary} (numerically: centered-template signal \ensuremath{\pm}0.0001 \ensuremath{\approx} 0);
\item \textbf{O4 (information capacity)}: 4 nearest neighbors \ensuremath{\implies} 1 bit per 4 plaquettes \ensuremath{\implies} S\down{max} = (A/4l\down{P}\up{2})\ensuremath{\cdot}ln2\mtag{Framework corollary} --- holography = geometric necessity of observer limitation.
\end{itemize}

\textbf{Present-plane selection propositions (N1--N4)}: dynamical basis of O1 --- why the observer aligns with exactly one temporal plane\mtag{Framework assumption}:
\begin{itemize}
\item \textbf{N1 (existence)}: a steady vortex requires phase accumulation of integer multiples of 4\ensuremath{\pi}\mtag{Framework assumption}\ensuremath{\implies} must lock a rotation axis \ensuremath{\implies} at least one plane;
\item \textbf{N2 (uniqueness)}: 0 planes degenerate (no vortex), 2 planes conflict (shared-axis beating), 3 planes isotropic --- exactly one (numerically: stabilities 0.00/1.00/0.30/0.05);
\item \textbf{N3 ("present" definitional)}: three planes fully equivalent (numerically: equal probabilities 0.332/0.334/0.334 \ensuremath{\approx} 1/3) \ensuremath{\implies} "present" = the selected plane (definition), no a-priori special status --- dissolving the "why now" circularity;
\item \textbf{N4 (dynamical mechanism)}: initial perturbation\mtag{Framework assumption: "an accidental surplus of positive matter, an initial perturbation direction"}breaks the SO(3) degeneracy (spontaneous symmetry breaking) \ensuremath{\implies} locks a single plane; connects to Axiom 3 (time-arrow bifurcation).
\end{itemize}

\textbf{Initial-perturbation origin propositions (P1--P4)}: the seed of N4 --- why the "accidental surplus positive matter" is accidental\mtag{Framework assumption}:
\begin{itemize}
\item \textbf{P1 (origin = \ensuremath{\varepsilon} leakage)}: the "accidental surplus positive matter" = positive matter produced by leakage --- "the anti fraction kicked to the mirror branch, the positive fraction stays in our branch";
\item \textbf{P2 (necessity)}: vacuum ground-state instability ("absolute uniformity is unstable; fluctuations are unavoidable") + hairy-ball-theorem topological forcing + pre-electroweak expected leakage \ensuremath{\sim}10\up{22} \ensuremath{\implies} the initial perturbation necessarily exists;
\item \textbf{P3 (contingency)}: leakage time Poisson-random + direction isotropically random \ensuremath{\implies} selected plane equiprobable 1/3 (consistent with N3);
\item \textbf{P4 (quantitative: F26 unification)}: \ensuremath{\varepsilon}\down{0} = \ensuremath{\varepsilon}(1TeV)/(M\down{Pl}/1TeV)\up{0.615} = \textbf{7.7\ensuremath{\times}10\up{-20}} (P0 cross-checked); verification \ensuremath{\varepsilon}(1meV) = 1.01 \checkmark --- \textbf{initial perturbation (\ensuremath{\varepsilon}\down{0}), baryon asymmetry (6\ensuremath{\times}10\up{-10} @1TeV), dark energy ((meV)\up{4}): three sources, one origin}; the observer model and cosmology side (\S{}8) formally share one origin.
\end{itemize}

\textbf{Two-layer structure of "necessary contingency"}\mtag{Derivation}: necessity (topology + instability) guarantees "an initial perturbation will exist"; contingency (Poisson + isotropy) determines "where, when, and in which direction" --- both the "accidental" and "necessary" aspects of the phrase "accidental surplus positive matter" are simultaneously explained.

\textbf{Six occurrences of "4" share one source}: surface plaquette 4 neighbors (observer geometry, framework corollary) = Clifford 4 components (interior space, framework assumption) = flux-tube coherent cross-section 4l\down{P}\up{2} (k model) = 1 bit per 4 plaquettes (framework assumption) = Kakeya packing 1/4 (framework corollary) = regular-tetrahedron 4-direction cluster (Tammes n=4) --- \textbf{the information structure encodable by the observer = interior-space structure = the inverse of surface information density}; holographic concentration and microstate counting connect seamlessly.

\textbf{\mtag{v1.21 direction discretization and the Bekenstein bound (B4) + observer mutual-information area law (B5)}}:

\textbf{Direction discretization = dynamical corollary of the Planck tube diameter (not an assumption)}\mtag{Derivation + Numerical} --- the framework assumes "the finest scale of a direction tube is l\down{P} (any finer and the storm's discreteness emerges, the tube disintegrates)" \ensuremath{\implies} direction space discretizes into a direction cluster (no continuous rotation symmetry, framework assumption) \ensuremath{\implies} Wang Hong H-dim=3 \ensuremath{\implies} k=4 optimal packing (framework corollary: Tammes regular tetrahedron rigorous) \ensuremath{\implies} 4 basis directions whose convex cones cover all direction space (numerically: 20000 random directions all reconstructable with positive combinations, error 10\up{-6}\ensuremath{^{\circ}}) \ensuremath{\implies} macroscopic continuous directions = statistical-average emergence of 4 basis directions (central limit, angular error \ensuremath{\sim}1/\ensuremath{\sqrt{N}}). \textbf{\ensuremath{\implies} Bekenstein bound recovered from direction discretization}: each Planck plaquette can "tube" only a finite direction cone (k=4 \ensuremath{\implies} 1 bit per 4 plaquettes) \ensuremath{\implies} \textbf{S \ensuremath{\leq} A/(4l\down{P}\up{2})} --- doubly closed with the k=2\up{1/4} microstate counting ("=" directions) (upper and lower bounds given by two independent mechanisms).

\textbf{Observer mutual-information area law = corollary of bulk = surface projection}\mtag{Derivation + Numerical} --- observer (present-plane soliton) reads only the surface\mtag{Framework assumption: "connection states in four directions"}+ bulk = surface projection\mtag{Framework assumption}\ensuremath{\implies} H(R|\ensuremath{\Sigma})=0 \ensuremath{\implies} \textbf{I(R;\ensuremath{\Sigma})=H(\ensuremath{\Sigma})=M bits \ensuremath{\propto} A (area law)} (toy-model numerics: M=4/8/16/32 \ensuremath{\to} I=M \checkmark); \textbf{16-fold compression}: observer's raw 4 bits/plaquette \ensuremath{\times} coherent block (4 plaquettes \ensuremath{\times} 4 directions = 16 degrees of freedom \ensuremath{\to} 1 bit, k model) \ensuremath{\implies} entropy density 1/4 bit/plaquette = 1/(4l\down{P}\up{2})\mtag{Framework corollary: "1 bit per 4 plaquettes"} --- unified observer information-capacity formula S\down{max}=(A/l\down{P}\up{2})\ensuremath{\times}(4\ensuremath{\div}16)=A/(4l\down{P}\up{2}).

\subsection{4.10 Black-hole interior as a 3D Kakeya set: packing proposition package}

\begin{quote}
The full chapter draft is in the standalone Chinese file (Paper\down{Sec4}.10\down{packing}\down{theorem}\down{package}.md); this section is the embedded concise version.
\end{quote}

\textbf{\mtag{Framework assumption}Physical setting}: black-hole interior = storm-overload region = 3D Kakeya set (contains all inward directions but information measure \ensuremath{\mu}\down{info}(V)\ensuremath{\approx}0; information lives only in horizon-surface saturated modes); quantum line = Kakeya segment (cross-branch ER channel)\mtag{Framework assumption}; horizon = time-arrow-flip surface\mtag{Framework assumption}.

\textbf{Theorem 1 (direction-cluster-number lower bound k \ensuremath{\geq} 4)\mtag{New proof, rigorous}}: Let V\ensuremath{\subset}S\up{2} be the direction cluster of quantum-line endpoints. If the covering radius is strictly < 90\ensuremath{^{\circ}} (Kakeya "covers all directions" requirement, no zero-information-residual direction), then k \ensuremath{\geq} 4.

\textit{Proof}: Suppose 3 open hemispheres cover S\up{2} \ensuremath{\implies} 0 \ensuremath{\in} int(conv{v\down{1},v\down{2},v\down{3}}) (separating-hyperplane theorem) \ensuremath{\implies} v\down{i} coplanar encircling \ensuremath{\implies} take plane normal p: p\ensuremath{\cdot}v\down{i} = 0 \ensuremath{\forall}i, covering radius exactly \ensuremath{\geq} 90\ensuremath{^{\circ}}, contradiction. \ensuremath{\blacksquare} (Numerically: k=3 optimal covering radius = 90.05\ensuremath{^{\circ}}, k=4 = 70.291\ensuremath{^{\circ}})

\textbf{Theorem 2 (k=4 configuration)\mtag{Theorem}}: Tammes n=4 --- the unique optimal configuration of 4 points = regular tetrahedron, minimal angle arccos(\ensuremath{-}1/3) \ensuremath{\approx} 109.47\ensuremath{^{\circ}} (random 2000 trials optimal only 96.9\ensuremath{^{\circ}}).

\textbf{Theorem 3 (exact Voronoi partition)\mtag{Theorem, rigorous}}: the regular tetrahedron's 4 directions partition S\up{2} into 4 congruent cells, each of area 4\ensuremath{\pi}/4 = \ensuremath{\pi} (Monte Carlo: 3.138--3.148 \ensuremath{\approx} \ensuremath{\pi}).

\textbf{Argument 4 ("exactly 4" = Clifford interior space)\mtag{Structural argument}}: regular-tetrahedron 4 directions = interior-space component structure (Clifford Cl(1,3) representation dimension 4 = two time directions \ensuremath{\times} two chirality directions); octahedron 6 directions = exterior-space axis structure (3 axes \ensuremath{\pm}). Entropy encoding uses interior-space components \ensuremath{\implies} direction-cluster number 4. The three candidates (four quadrants / Hopf fibration / SU(2) double cover) unify in the same "4".

\textbf{Proposition 5 (1/4 coefficient)\mtag{Derivation}}: Theorem 1 (k\ensuremath{\geq}4) \ensuremath{\to} Theorem 2 (regular tetrahedron) \ensuremath{\to} Theorem 3 (Voronoi \ensuremath{\pi}/cell) \ensuremath{\to} Argument 4 (4=Clifford) \ensuremath{\to} Nyquist (coherent cross-section 4l\down{P}\up{2}) \ensuremath{\to} 1 bit per unit \ensuremath{\implies} \textbf{S = (A/4l\down{P}\up{2})\ensuremath{\cdot}ln2}, cross-validated with the \S{}4.8 microstate counting as \textbf{two paths}.

\textbf{Proposition 6 (interior-space encoding priority)\mtag{Derivation, framework-internal deduction}}: the final step of Argument 4 --- why entropy encoding uses the interior-space 4 components rather than the exterior-space 3 axes (6-direction octahedron). Three-lemma chain:
\begin{itemize}
\item Lemma A (exterior-space entropy = 0): isotropy (Axiom 1) \ensuremath{\implies} no information in directional preference (uniform prior, numerically \ensuremath{\|}\ensuremath{\langle}n\ensuremath{\rangle}\ensuremath{\|}\ensuremath{\approx}0); no-hair\mtag{Framework assumption: no-hair}\ensuremath{\implies} exterior structure erased; \textbf{event point-like-ness\mtag{Assumption H-8: event point-like-ness} "two vortices dragged in different directions cannot coexist in the same micro-variation event" \ensuremath{\implies} directions are not intra-event attributes};
\item Lemma B (interior-space non-zero): quantum-line endpoints carry Clifford 4-component states latched across branches\mtag{Framework assumption}; Hawking radiation = leakage of interior-space states \ensuremath{\implies} the interior space is the information carrier;
\item Lemma C (4 = minimal complete spinor): Cl(1,3) spinor dimension 4 = two time directions \ensuremath{\times} two chirality directions (1D trivial, 2D Weyl missing half, 4D Dirac complete --- representation theory);
\item Proposition 6: entropy encoding = interior-space 4 components \ensuremath{\implies} direction-cluster number = 4. \textbf{The three premises A1/A2/A3 are all existing framework assumptions; no new assumptions introduced}.
\end{itemize}

\textbf{Testable predictions (corollaries)}: B1 extreme-Kerr Hawking-radiation angular distribution deviates from axisymmetry (\ensuremath{\sim}10\up{-8}); B2 primordial black-hole entropy volume correction \ensuremath{\Delta}S/S \ensuremath{\sim} l\down{P}/R (\ensuremath{\sim}10\up{-8}); B3 Hawking-radiation--infalling-matter entanglement (information-flip verdict signature); B4 ER-bridge packing-density bound (BBH angular-mode distortion); B5 non-thermal correction \ensuremath{\propto}1/S\down{BH} (\ensuremath{\sim}10\up{-20} window).

\textbf{Honest notes}: Theorems 1, 3 rigorous; Theorem 2 established; Argument 4 structural; Proposition 6 framework-internal deduction (three premises direct axioms); \mtag{Assumption} direction discretization (requires dynamical proof, \S{}10.2.1).


\section{5 Gauge-Group Structure}

\subsection{5.1 "Pairing order" and the three gauge groups}

\mtag{Derivation} Define the \textbf{pairing order} n = the smallest positive integer such that the representation tensor product \ensuremath{\otimes}\up{n} contains the trivial representation:

\begin{table}[htbp]
\centering
\small
\begin{tabularx}{\textwidth}{@{}lXXX@{}}
\toprule
Group & n & Representation theory & Physics \\
\midrule
U(1) & 1 & all one-dimensional, q\ensuremath{\otimes}(\ensuremath{-}q)=1 & long-range, massless photon \\
SU(2)\down{L} & 2 & 2 pseudo-real (2\ensuremath{\cong}\ensuremath{\bar{2}}, \ensuremath{\varepsilon}\up{ij} sign-preserving) & short-range (massive W/Z), chiral \\
SU(3)\down{c} & 3 & 3 complex (3\ensuremath{\ncong}\ensuremath{\bar{3}}), \ensuremath{\varepsilon}\down{ijk} third-order & color confinement (qqq cube has no solution) \\
\bottomrule
\end{tabularx}
\end{table}

\textbf{Representation-theoretic root of color confinement}: colorless combinations of baryons qqq require the third-order antisymmetric invariant tensor \ensuremath{\varepsilon}\down{ijk} --- three quarks form a singlet via the third-order tensor. Mesons q\ensuremath{\bar{q}} use \ensuremath{\delta}\down{i}\up{j} (n=2 structure).

\textbf{\mtag{v2.5 definition revision (referee M6)}}: the pairing order is precisely defined as \textbf{n \ensuremath{\equiv} d\down{min}(G)} (smallest non-trivial irreducible representation dimension = number of color components on string-net edges), supplemented by the \textbf{same-sign pairing order n\ensuremath{^\prime}} (smallest n such that \ensuremath{\rho}\ensuremath{\otimes}\up{n} without conjugation contains the trivial representation): U(1) n\ensuremath{^\prime}=\ensuremath{\infty} (same-sign never contains trivial \ensuremath{\implies} \textbf{must pair positive--anti = representation-theoretic root of charge conjugation C}), SU(2) n\ensuremath{^\prime}=2 (2\ensuremath{\otimes}2=3\ensuremath{\oplus}1, pseudo-real), SU(3) n\ensuremath{^\prime}=3 (3\ensuremath{\otimes}3\ensuremath{\otimes}3=10\ensuremath{\oplus}8\ensuremath{\oplus}8\ensuremath{\oplus}1, \ensuremath{\varepsilon}\down{ijk}). \textbf{Structural necessity: non-Abelian-ness \ensuremath{\iff} n\ensuremath{^\prime} finite \ensuremath{\iff} same-sign tensor products can form singlets} (see \S{}5.2 screening argument).

\subsection{5.2 Screening argument (why SU(3)\ensuremath{\times}SU(2)\ensuremath{\times}U(1) --- structural screening)}

\textbf{\mtag{v2.5 structural screening (referee M6)}}: Given the pairing order n, the candidate group pool (U(1)/Z\down{N}/SU(2)/Q\down{8}/S\down{3}/SO(3)/SU(3)/G\down{2}/SU(4)/Sp(2), 10 total) is eliminated one by one by three criteria, with a \textbf{unique survivor at each step} (numerical verification in the M6 report, dimension conservation all \checkmark):

\begin{itemize}
\item \textbf{S1 dimension matching}: d\down{min}(G) = n (n = vortex winding-mode number, \S{}5.3 input)
\item \textbf{S2 same-sign pairing}: n\ensuremath{^\prime}(G) = n (non-Abelian-ness \ensuremath{\iff} n\ensuremath{^\prime} finite)
\item \textbf{S3 third-order antisymmetry}: n=3 requires \ensuremath{\varepsilon}\down{ijk} (det=1 \ensuremath{\implies} unique to the SU(N) sequence)
\end{itemize}

\begin{table}[htbp]
\centering
\small
\begin{tabularx}{\textwidth}{@{}lXXX@{}}
\toprule
n & Candidates (d\down{min} matching) & Reason excluded & Unique survivor \\
\midrule
1 & U(1), Z\down{N} & Z\down{N} discrete, Villain limit = U(1) (v1.30 Theorem A generalized) & \textbf{U(1)} (unique 1D connected compact Abelian Lie group) \\
2 & SU(2), Q\down{8}, Sp(2), S\down{3} & Q\down{8} \ensuremath{\subset} SU(2) discrete; Sp(2) symplectic pseudo-real (chirality picks SU(2)); S\down{3} same class (d\down{min}=2) & \textbf{SU(2)} (2 pseudo-real \ensuremath{\to} chiral) \\
3 & SU(3), SO(3), G\down{2}, SU(4), A\down{4}/S\down{4}/\ensuremath{\Delta}(27) & \textbf{SO(3) n\ensuremath{^\prime}=2 (3 real self-dual \ensuremath{\to} 3\ensuremath{\otimes}3 contains trivial \ensuremath{\to} no confinement forcing)}; G\down{2} d\down{min}=7; SU(4) n=4; A\down{4} etc. are SU(3) discrete subgroups & \textbf{SU(3)} (3 complex 3\ensuremath{\ncong}\ensuremath{\bar{3}} + \ensuremath{\varepsilon}\down{ijk} \ensuremath{\to} qqq third-order singlet) \\
\bottomrule
\end{tabularx}
\end{table}

\textbf{Decisive exclusion (SO(3))}: although d\down{min}(SO(3))=3, its 3-dimensional representation is real self-dual (3\ensuremath{\cong}\ensuremath{\bar{3}}) \ensuremath{\to} same-sign pairing order n\ensuremath{^\prime}=2 \ensuremath{\neq} 3 (3\ensuremath{\otimes}3=5\ensuremath{\oplus}3\ensuremath{\oplus}1 already contains trivial) \ensuremath{\to} \textbf{structurally incapable of third-order confinement} --- a direct answer to "why SU(3) rather than SO(3)".

\textbf{n\ensuremath{\geq}4 exclusion}: three spatial dimensions (locked by the Kakeya dimension) \ensuremath{\implies} the fourth-order antisymmetric \ensuremath{\varepsilon}\down{ijkl} requires 4-dimensional space; n=4 Fermat has no solution (original argument retained); SU(2)\ensuremath{\times}U(1) cannot merge (chirality, original argument retained).

\textbf{Discrete--continuum bridge (correcting the crude A2 mapping)}: the framework's numerical layer uses quantum-double D(G) discrete groups --- Z\down{N}\ensuremath{\subset}U(1), \textbf{Q\down{8}\ensuremath{\subset}SU(2)} (i\up{2}=j\up{2}=k\up{2}=\ensuremath{-}I), \textbf{S\down{3}\ensuremath{\subset}SU(3)} (even-permutation matrices, det=1) --- with unique continuous parent groups; \textbf{S\down{3} itself has d\down{min}=2 and belongs to the n=2 class; it is used as an SU(3) permutation subgroup for non-Abelian numerical verification (v1.33)}.

\textbf{Residual contingency localized}: after screening uniqueness, the only remaining input = the micro-origin of the vortex winding-mode numbers n=1,2,3 (\S{}5.3 assumption, not derivable from representation theory --- honestly retained).

\subsection{5.3 Three fermion generations = vortex winding modes}

\mtag{Framework assumption} Each fermion generation = excitation of the same vortex winding mode at different energy scales: first generation n=1, second n=2, third n=3. \textbf{Fourth-generation prediction}: n=4 vortices decay unstably at z>1000, extremely heavy and short-lived above 1 TeV.

\subsection{5.4 Kaluza--Klein equivalence and terminology resolution}

\mtag{Derivation} Abandon the phrasing "fifth dimension = flux-field dimension"; instead: \textbf{"the flux field \ensuremath{\Phi} = low-energy projection of the KK scalar mode field (g\down{55}), whose VEV determines coupling constants and masses"}. The mathematical identity of D[\ensuremath{\Phi}]=S is the \textbf{spectral action} (S = Tr f(D\up{2}/\ensuremath{\Lambda}\up{2}), heat-kernel expansion a\down{1}\ensuremath{\to}R, a\down{2}\ensuremath{\to}F\up{2}).

\subsection{5.5 Multi-representation unified field: four forces = partial-wave decomposition}

\mtag{Derivation} Decompose flux-field fluctuations \ensuremath{\delta}\ensuremath{\Phi}(x) into Lorentz-group irreducibles:

\begin{table}[htbp]
\centering
\small
\begin{tabularx}{\textwidth}{@{}lXXX@{}}
\toprule
Component & Spin & Fluctuation type & Corresponding force \\
\midrule
\ensuremath{\Phi} scalar & 0 & isotropic & Higgs/modulus field \\
A\down{\mu{}} vector & 1 & directional & electromagnetism U(1) \\
h\down{\mu{}\nu{}} tensor & 2 & second-order tensor & gravity \\
\bottomrule
\end{tabularx}
\end{table}

\textbf{New statement of unification}: not "different limits of one operator" but "\textbf{the fluctuation spectrum of one flux field}" --- the differentiation of forces = partial-wave decomposition of fluctuation modes. \mtag{To be verified} complete emergence of the non-Abelian components (SU(2)\ensuremath{\times}SU(3)), mass gap, chiral dynamics.

\subsection{5.6 Quantum-double algebra: the framework's gauge structure (v1.19)}

\mtag{Derivation + Proposition} The framework's three event-flow operations from Axiom 1 (encounter/annihilation/regeneration) + P2-A non-commutativization \ensuremath{\implies} Hopf-algebra structure, with the most natural discrete realization = \textbf{quantum double D(G) = C[G]\ensuremath{\ltimes}F(G)} (Drinfeld double). Proposition chain (stage-A proposition package A1--A6):

\begin{table}[htbp]
\centering
\small
\begin{tabularx}{\textwidth}{@{}lXX@{}}
\toprule
Proposition & Content & Nature \\
\midrule
A1 & Three event-flow operations \ensuremath{\leftrightarrow} three Hopf structures (encounter=multiplication, annihilation=involution, regeneration=comultiplication) & framework\ensuremath{\to}structure \\
A2 & Pairing order \ensuremath{\to} group selection (v2.5 revision): discrete layer n=1\ensuremath{\to}Z\down{N} (\ensuremath{\subset}U(1)) / n=2\ensuremath{\to}Q\down{8} (\ensuremath{\subset}SU(2), pseudo-real) / n=3\ensuremath{\to}S\down{3} (\ensuremath{\subset}SU(3), minimal non-Abelian); continuous parent groups \textbf{U(1)/SU(2)/SU(3) unique} (\S{}5.2 three-criterion screening, structural not fitting) & \mtag{Proposition}+ numerical \\
A3 & annihilation--regeneration \ensuremath{\to} gauge operators A\down{v} (vertex = generator of gauge transformations) / B\down{p} (face = flux closure) & \mtag{Derivation+Numerical}\checkmark \\
A4 & flux tube = confining string \ensuremath{\to} area law (2D YM \ensuremath{\sigma}=g\up{2}C\down{2}/2 + strong-coupling Wilson) & \mtag{Proposition} \\
A5 & \textbf{main theorem: emergence of gauge symmetry} (Kitaev 2003: H=\ensuremath{-}\ensuremath{\Sigma}A\down{v}\ensuremath{-}\ensuremath{\Sigma}B\down{p} local symmetry = gauge symmetry) & \mtag{Theorem citation} \\
A6 & why n=3 (Clifford interior-space 4 components / number-line classification / Fermat-solvability candidates) & \mtag{Open} \\
\bottomrule
\end{tabularx}
\end{table}

\textbf{Numerical verification (Kitaev single-face lattice)}: Z\down{2}(16 states)/Z\down{3}(81 states)/S\down{3}(216 states) all pass --- A\down{v}\up{2}=A\down{v}, B\down{p}\up{2}=B\down{p}, [A\down{v},B\down{p}]=0, [A\down{v},A\down{v}\ensuremath{^\prime}]=0, ground-state dimension=1; D(S\down{3}) has 8 anyons (6 non-Abelian, quantum dimensions 2,3), D\up{2}=36=|G|\up{2}, fusion-dimension identity \ensuremath{\tau}\ensuremath{\otimes}\ensuremath{\tau}=A\down{1}\ensuremath{\oplus}A\down{2}\ensuremath{\oplus}\ensuremath{\tau} \checkmark.

\textbf{Core significance}: \textbf{non-Abelian gauge structure can emerge exactly from discrete algebra}\mtag{Theorem citation, Kitaev 2003} --- the framework's "2 vs 3 \ensuremath{\to} confinement" has rigorous mathematical existence at the discrete level (honest note: S\down{3}\ensuremath{\neq}SU(3); structural correspondence, not group isomorphism).

\subsection{5.7 2+1D framework Hamiltonian and emergent gauge symmetry (v1.19)}

\mtag{Derivation + Numerical} Explicit mapping from framework axioms to Hamiltonian (stage-B proposition package B1--B5):

\[H_F = -J\sum_v A_v - K\sum_p B_p\]

\begin{table}[htbp]
\centering
\small
\begin{tabularx}{\textwidth}{@{}lXX@{}}
\toprule
Framework axiom & Hamiltonian term & Physical meaning \\
\midrule
Axiom 1: encounter (pairing symmetry) & \ensuremath{-}J\ensuremath{\cdot}A\down{v} & kinetic energy of gauge-transformation generators \\
Axiom 1: annihilation--regeneration (circulation conservation) & \ensuremath{-}K\ensuremath{\cdot}B\down{p} & flux closure (flux tube) \\
vacuum = 50:50 balance & \textbf{ground state} (A\down{v}=B\down{p}=1) & perfect annihilation balance \\
\ensuremath{\varepsilon} leakage & vertex violation (charge anyon) & seed of baryon asymmetry \\
\bottomrule
\end{tabularx}
\end{table}

\textbf{Proposition chain B1--B5}:
\begin{itemize}
\item \textbf{B1 emergent gauge symmetry}: [H\down{F}, A\down{v}]=[H\down{F}, B\down{p}]=0 (A\down{v}/B\down{p} projectors + commutation \ensuremath{\implies} linear combinations commute) --- \textbf{gauge symmetry is a dynamical symmetry of the Hamiltonian, not an externally imposed constraint};
\item \textbf{B2 vacuum = perfect annihilation balance}: ground state satisfies all A\down{v}=B\down{p}=1, E\down{0}=\ensuremath{-}(N\down{v}+N\down{p});
\item \textbf{B3 energy gap \ensuremath{\Delta}=1>0}: rigorous example of a gapped gauge theory in 2+1D (topological order deconfined)\mtag{Theorem, Kitaev 2003};
\item \textbf{B4 leakage excitations = anyons/flux tubes}: A\down{v}=0\ensuremath{\to}charge anyon e, B\down{p}=0\ensuremath{\to}flux m, e\ensuremath{\otimes}m composite --- \ensuremath{\varepsilon} leakage = vertex violation = seed of baryon asymmetry;
\item \textbf{B5 torus ground-state degeneracy = |G|\up{2}}: topological-order fingerprint (Z\down{2} torus numerics = 4 \checkmark; Z\down{3}=9, S\down{3}=36 cited).
\end{itemize}

\textbf{Numerical verification (Z\down{2} 2\ensuremath{\times}2 torus, 256 states, all pass)}: [A\down{w},A\down{v}]=0, [B\down{p},A\down{v}]=0, [B\down{p},B\down{q}]=0 \ensuremath{\implies} [H\down{F},A]=[H\down{F},B]=0 \checkmark; ground state E=\ensuremath{-}8.0000 \checkmark; torus degeneracy tr=4.000=|G|\up{2} \checkmark; gap \ensuremath{\Delta}=1 \checkmark.

\textbf{Numerical realization of confined string tension (v1.20)} --- the framework's "confinement = flux-tube linear potential"\mtag{Framework mechanism: confinement = flux-tube linear potential; citation Wilson 1974/Kogut--Susskind 1974} obtains \textbf{thermodynamic-limit-convergent quantitative numerics} in 2+1D lattice gauge theory (correct geometry + constrained-sector method; earlier numerics (pre-v1.19 \ensuremath{\varepsilon}-term/endpoint-potential reports) were found on review to contain geometric-convention and measurement-method defects and were superseded by the P0-1 review report):

\begin{table}[htbp]
\centering
\footnotesize
\begin{tabularx}{\textwidth}{@{}lXXXX@{}}
\toprule
Field version & Condensed object & Confined object & String tension (thermodynamic limit) & Phase \\
\midrule
\textbf{X field} \ensuremath{-}\ensuremath{\kappa}\ensuremath{\Sigma}X\down{e} & electric flux (\ & +\ensuremath{\rangle} condensation) & \textbf{charge e-e} & \textbf{\ensuremath{\sigma}\down{E} \ensuremath{\approx} 2\ensuremath{\kappa}} (0.969/1.984/3.992 at \ensuremath{\kappa}=0.5/1.0/2.0; three sizes 3\ensuremath{\times}3/3\ensuremath{\times}4/4\ensuremath{\times}4 converge <0.02\%; \ensuremath{\kappa}=0 strictly deconfined) \\
\textbf{Z field} \ensuremath{-}\ensuremath{\lambda}\ensuremath{\Sigma}Z\down{e} & magnetic flux & \textbf{flux m-m} & \textbf{\ensuremath{\sigma}\down{m} \ensuremath{\approx} 2\ensuremath{\lambda}} (0.21/1.09/3.04 at \ensuremath{\lambda}=0.5/1.0/2.0, 3\ensuremath{\times}4 constrained sector; \ensuremath{\lambda}=0 strictly deconfined) & magnetic condensation \\
\bottomrule
\end{tabularx}
\end{table}

\begin{itemize}
\item \textbf{Constrained-sector method\mtag{Derivation, Numerical}}: an A\down{v}-flipped vertex flips 4 edges \ensuremath{\to} each face flipped 0 or 2 times \ensuremath{\to} \textbf{the B\down{p} sector is an invariant subspace of H} --- find the ground state in the "faces A, B violated (B\down{p}=0)" subspace; E\down{pair}(d)=2E\down{m}+\ensuremath{\sigma}\down{m}\ensuremath{\cdot}d gives the \textbf{true flux-pair string potential} (the old "boundary-path-flip" method measured the flux-number effect, not the string tension; deprecated);
\item \textbf{\ensuremath{\kappa}=0/\ensuremath{\lambda}=0 strictly deconfined}\mtag{Derivation, Numerical}: dE independent of distance (=2E\down{m} constant) --- correct topological-order behavior, verified for both field versions;
\item \textbf{Phase-transition picture corrected (v1.20)}: on the torus, ground-state degeneracy for \ensuremath{\lambda}>0 drops \textbf{immediately} 4\ensuremath{\to}1 (degeneracy is an exact symmetry of \ensuremath{\lambda}=0; no finite-size transition point); string density decreases continuously from 0.5 (maximum entropy) to 0 (vacuum polarization); gap opens monotonically --- the true topological-order\ensuremath{\to}trivial transition (if any) requires \textbf{infinite-lattice finite-size scaling} to localize (the early "\ensuremath{\lambda}\down{c}\ensuremath{\approx}0.12--0.15" was a criterion-threshold artifact, withdrawn);
\item \textbf{Transition-point localization (v1.22 P0-3 open-boundary FSS \ensuremath{\to} v1.23 P0-3b torus refinement)} --- finite-size scaling of the topological-order\ensuremath{\to}confinement transition\mtag{Numerical}. Order parameter m(\ensuremath{\kappa},L)=\ensuremath{\Sigma}\down{e}\ensuremath{\langle}X\down{e}\ensuremath{\rangle}/E (m=0 for the toric-code ground state at \ensuremath{\kappa}=0; m\ensuremath{\to}1 as \ensuremath{\kappa}\ensuremath{\to}\ensuremath{\infty}), \ensuremath{\kappa}*(L) taken at the m=0.5 crossing:
\begin{itemize}
\item \textbf{Open-boundary FSS (v1.22, P0-3)}: \ensuremath{\kappa}* increases with L (0.0907\ensuremath{\to}0.1300), 1/L extrapolation \ensuremath{\kappa}\down{c}=0.2255 (0.2707 after dropping the smallest 2\ensuremath{\times}3 lattice) --- estimate \ensuremath{\kappa}\down{c}(\ensuremath{\infty})\ensuremath{\approx}0.23--0.27; the "\ensuremath{\kappa}* increasing with L" is mainly a boundary effect (confirmed by the torus, v1.23 below);
\item \textbf{Torus (periodic boundary, v1.23, P0-3b)}: \ensuremath{\kappa}* nearly \textbf{size-independent} across 4 sizes ---
\end{itemize}
\end{itemize}

    | Torus | Edges | States | L=\ensuremath{\sqrt{area}} | \textbf{\ensuremath{\kappa}\*(m=0.5)} | \ensuremath{\chi} peak |
    |---|---|---|---|---|---|
    | 2\ensuremath{\times}2 | 8 | 256 | 2.000 | \textbf{0.1942} | 0.17 |
    | 2\ensuremath{\times}3 | 12 | 4096 | 2.449 | \textbf{0.2017} | 0.19 |
    | 2\ensuremath{\times}4 | 16 | 65536 | 2.828 | \textbf{0.2044} | 0.20 |
    | 3\ensuremath{\times}3 | 18 | 262144 | 3.000 | \textbf{0.1994} | 0.19 |

    Ground-state verification all passes (E0=\ensuremath{-}(N\down{v}+N\down{p}) exact, m(0)=0, \textbf{4-fold ground-state degeneracy |G|\up{2}=4, \ensuremath{\Delta}(0)=0}, 2\ensuremath{\times}2 torus 256 states exact diagonalization); the symmetric convention (J\down{A}=J\down{B}=1, standard literature form) gives identical results at the crossing --- \textbf{the framework Hamiltonian \ensuremath{\equiv} standard toric code + transverse field, with \ensuremath{\kappa} directly corresponding to the literature h}. \ensuremath{\kappa}* 1/L extrapolation \textbf{0.217} (slope \ensuremath{-}0.044, nearly flat), \ensuremath{\chi}-peak extrapolation \textbf{0.247};
\begin{itemize}
\item \textbf{Criterion correction (v1.23, important)}: the m=0.5 crossing is a \textbf{half-polarization point}, not the transition point --- torus 3\ensuremath{\times}3 measures m(h=0.32)\ensuremath{\approx}0.84 (m already far above 0.5 at the literature transition point); the order parameter of the charge-condensation transition is the string/Wilson loop, not single-site magnetization (honest note). \textbf{True transition estimate \ensuremath{\kappa}\down{c} \ensuremath{\approx} 0.25--0.32} (lower bound: \ensuremath{\chi}-peak extrapolation 0.247; upper bound: literature infinite-size value);
\item \textbf{Literature-comparison correction (v1.23, important)}: h\down{c}=0.328474(3) is the critical value of the \textbf{\ensuremath{\sigma}\up{z} field (Z field / flux-condensation version)} (dual 2+1D transverse-field Ising, Trebst 2007/Bl\"{o}te--Deng 2002/Wu 2012); the framework's \textbf{X-field version (\ensuremath{\sigma}\up{x} field / charge condensation) should be compared with h\down{xc}\ensuremath{\approx}0.32} (Trebst 2007/Vidal 2009, verified via arXiv:1304.0780) --- the v1.22 citation "h\down{c}\ensuremath{\approx}0.328 same order" mixed the two versions, now corrected;
\item \textbf{Framework conclusions}: (1) "no transition point at finite size" is a finite-lattice manifestation; the true transition exists in the infinite-lattice limit; (2) \textbf{domain of validity of \ensuremath{\sigma}\down{E}\ensuremath{\approx}2\ensuremath{\kappa}}: \ensuremath{\kappa}<\ensuremath{\kappa}\down{c} deconfined (topological order), \ensuremath{\kappa}>\ensuremath{\kappa}\down{c} confined --- the string-tension numerics (\ensuremath{\kappa}=0.5/1.0/2.0 all > \ensuremath{\kappa}\down{c}\up{true\approx{}0.3}) lie in the confined phase \checkmark self-consistent; (3) confinement phase diagram complete (P0 stage closed);
\end{itemize}
\begin{itemize}
\item \textbf{Framework correspondence}: if \ensuremath{\varepsilon} leakage corresponds to flux (B\down{p} violation) \ensuremath{\to} \ensuremath{\sigma}\down{m} is its confinement mechanism; if to charge (A\down{v} violation) \ensuremath{\to} \ensuremath{\sigma}\down{E} is its confinement mechanism --- \textbf{two realizations of the flux-tube linear potential}.
\end{itemize}

\textbf{\mtag{v1.33 non-Abelian 3+1D numerical realization (action-list P0-2; cf. \mtag{confinement mechanism + BESIII anchor})}}\mtag{Numerical}: generalize the 2+1D Abelian framework to a \textbf{non-Abelian (S\down{3}) 3+1D face model} --- single cube (6 faces, face variables g\down{f} \ensuremath{\in} S\down{3}, 6\up{6}=46,656 states), H(\ensuremath{\lambda})=\ensuremath{-}A\down{c}\ensuremath{-}\ensuremath{\lambda}\ensuremath{\Sigma}\down{f} X\down{f}, \textbf{A\down{c}=\ensuremath{\delta}\down{e}(\ensuremath{\Pi}\down{f\in{}\partial{}c}g\down{f}) (non-Abelian product order)}, X\down{f}: g\down{f}\ensuremath{\to}(12)g\down{f}:
\begin{itemize}
\item \textbf{\ensuremath{\lambda}=0 string-net condensation (rigorous)}: E0=\ensuremath{-}N\down{c} (=\ensuremath{-}1), ground-state degeneracy \textbf{6\up{n-1}=7776} (constraint counting; group order generalizes Z\down{2}'s 2\up{n-1}), closed-loop weight 1.0000;
\item \textbf{Non-Abelian fingerprint\mtag{Key}}: excitations (A\down{c} violations) split into \textbf{two classes by conjugacy class} --- transpositions 23,328 (6\up{5}\ensuremath{\times}3), 3-cycles 15,552 (6\up{5}\ensuremath{\times}2) --- the Abelian Z\down{2} has only one class; non-Abelian "flux" carries group elements and decomposes irreducibly by conjugacy class (non-trivial fusion rules/braiding statistics);
\item \textbf{Gap spectrum}: opens monotonically with \ensuremath{\lambda} (\ensuremath{\lambda}=0.1 \ensuremath{\Delta}=0.037 \ensuremath{\to} \ensuremath{\lambda}=1.0 \ensuremath{\Delta}=1.546), closed-loop weight decreases (condensation-phase stability requires \ensuremath{\lambda}\down{c} analysis, C1);
\item \textbf{Glueball experimental anchor}[BESIII 2026]: BESIII announced at ICHEP 2026 that the main component of X(2370) is a \textbf{glueball} (0\up{-+}, \ensuremath{\sim}2395 MeV/c\up{2}, flavor singlet K(892)\ensuremath{\bar{K}}<2.7\ensuremath{\times}10\up{-6}, within the lattice-QCD predicted 2.3--3.0 GeV range) --- the framework's "glueball = bound state of the force/flux itself"[BESIII 2026] gains an experimental anchor; infrared slavery / quark confinement\mtag{Framework mechanism: infrared slavery = string tension}= string tension (splitting a flux tube grows new pairing) --- the non-Abelian counterpart of the framework's confinement mechanism.\mtag{Honest note} X(2370) is an experimental QCD fact; the framework's "flux self-binding" is a linguistic translation (framework-level interpretation, not experimental validation of the framework); S\down{3} is the minimal non-Abelian group (corresponding to an SU(2) discretization); complete 4D SU(3) emergence (mass gap/chirality) remains open (world-class).
\end{itemize}

\textbf{Significance}: the framework's "event flow \ensuremath{\implies} gauge structure" is upgraded from metaphor to \textbf{explicit Hamiltonian + dynamical-symmetry proof}; the 2+1D mass gap (topological order) is a low-dimensional rigorous realization of the 4D confinement mass gap (B6).

\subsection{5.8 Wen string-net 3+1D route (v1.19)}

\mtag{Framework + Citation, semi-rigorous} The framework's "quantum line = string" connects to Wen string-net condensation (stage C):

\begin{table}[htbp]
\centering
\small
\begin{tabularx}{\textwidth}{@{}lXX@{}}
\toprule
Framework concept & String-net language & Status \\
\midrule
quantum line & string-net edge (string segment) & framework + original \\
cosmic string network ("dynamic, continuously emerging network") & dynamically emergent string-net & original \\
matter--antimatter pair storm (Axiom 1) & \textbf{string-net condensation} (ground state = huge superposition of strings) & framework\ensuremath{\to}structure \\
pairing order n (A2) & string-net branching rules & framework\ensuremath{\to}structure \\
electromagnetism & U(1) string-net emergent photon (spin 1) & \mtag{Citation, semi-rigorous} \\
\textbf{gravity} & \textbf{string-net emergent spin-2 excitation (graviton)}\mtag{Citation, Wen 2006, semi-rigorous} & two-path cross-validation \\
\bottomrule
\end{tabularx}
\end{table}

\textbf{Numerical verification (3+1D Z\down{2} string-net)}: H(\ensuremath{\lambda})=\ensuremath{-}\ensuremath{\Sigma}A\down{cube}\ensuremath{-}\ensuremath{\lambda}\ensuremath{\Sigma}X\down{face} --- single cube (64 states): \ensuremath{\lambda}=0 closed-string weight 100\%, \ensuremath{\lambda}=0.5 gap \ensuremath{\Delta}=0.980, closure 58.2\% (>50\% random); double cube (2048 states): \ensuremath{\Delta}=0.963, double-closure 34\% (>25\% random). \textbf{Ground state favors closed endpoint-free string-nets (condensation) + string-endpoint excitations with positive gap (emergent particles)}.

\textbf{\mtag{v1.28 3+1D string-net rigorization (follow-up item (2), semi-rigorous \ensuremath{\to} partially rigorous)}}:
\begin{itemize}
\item \textbf{(1) quantum-line--string-net-edge Hopf bridge (B8 roadmap landed)}\mtag{Derivation}: the framework's three event-flow operations (Proposition A1) \ensuremath{\implies} Hopf structure \ensuremath{\implies} \textbf{operator-level correspondence} with string-net edge operators --- quantum-line endpoints \ensuremath{\leftrightarrow} string endpoints (A\down{c} violations); line encounter/annihilation/regeneration \ensuremath{\leftrightarrow} string-net branching/closure/fluctuation (proposition-level statement, extension of \S{}5.6 A1);
\item \textbf{(2) rigorous proof of \ensuremath{\lambda}=0 string-net condensation}\mtag{Derivation + Numerical}: H\down{0}=\ensuremath{-}\ensuremath{\Sigma}\down{c} A\down{c} (A\down{c}=\ensuremath{\Pi}\down{f\in{}\partial{}c}Z\down{f}) --- the A\down{c} commute and A\down{c}\up{2}=1 \ensuremath{\implies} spectrum = \ensuremath{-}\ensuremath{\Sigma}\ensuremath{\lambda}\down{c}; ground state has all A\down{c}=1 (closed-string-net condensation) \ensuremath{\implies} \textbf{E0 = \ensuremath{-}N\down{c} exactly}; \textbf{ground-state degeneracy = 2\up{n-N\down{c}}} (n faces, N\down{c} independent constraints; double cube 11 faces/2 cubes \ensuremath{\implies} 512=2\up{9} numerically verified \checkmark); \textbf{excitation gap \ensuremath{\Delta} = 2} (any single A\down{c} flip = string-endpoint/loop excitation = \textbf{emergent particle}, rigorous); at \ensuremath{\lambda}=0.1, \ensuremath{\langle}A\down{c}\ensuremath{\rangle}=0.877 (condensation maintained); subsystem entropy grows with boundary (non-trivial entanglement);
\item \textbf{(3) rigorous duality mapping}\mtag{Derivation + Citation}: 3+1D Z\down{2} string-net \ensuremath{\leftrightarrow} 3D Z\down{2} lattice gauge theory (Fradkin--Shenker line) --- \textbf{string-net condensed phase (small \ensuremath{\lambda}) \ensuremath{\leftrightarrow} gauge confined phase} (A\down{c}=1 \ensuremath{\leftrightarrow} B\down{plaq}=1; string endpoints \ensuremath{\leftrightarrow} charges); the \ensuremath{\lambda}=0 rigorous results (E0/degeneracy/gap) transfer by duality to rigorous statements in the confined phase;
\item \textbf{(4) emergent U(1) path made explicit}\mtag{Citation + Open}: Wen string-net bosonization (photon = string-net phase mode) --- rigorization in three steps: duality (rigorous \checkmark) \ensuremath{\to} continuum limit Z\down{2}\ensuremath{\to}U(1) (open) \ensuremath{\to} Maxwell L\down{eff}=\ensuremath{-}(1/4)F\up{2} verification (open); on the framework side, pairing order n=1 \ensuremath{\implies} U(1) already rigorous (\S{}5.1); \textbf{spin-2 (graviton) emergence rigorization = world-class problem}\mtag{Open, honest note}(same order as the 4D mass gap, not claimed).
\item \textbf{\mtag{v1.29 emergent U(1) continuum limit and spin 2 (follow-up item (2) remainder; cf. framework spin-2 mechanism)}}:
\begin{itemize}
\item \textbf{spin 2 = second-order-tensor rotation eigenvalues \ensuremath{\pm}2 (rigorous verification)}\mtag{Derivation + Numerical}: the second-order traceless symmetric tensor (5 components m=\ensuremath{\pm}2,\ensuremath{\pm}1,0) under SO(3) rotation h\ensuremath{\to}R h R\up{T} has 5-dimensional representation eigenvalues {e\up{\pm{}2i\theta{}}, e\up{\pm{}i\theta{}}, 1} (numerically: D\up{2} spectrum {0.6967\ensuremath{\times}2, 0.9211\ensuremath{\times}2, 1.0} = {cos2\ensuremath{\theta}, cos\ensuremath{\theta}, 1} \checkmark) --- \textbf{spin 2 is a geometric necessity of any second-order-tensor fluctuation mode}\mtag{Framework definition D-2}, with representation-theoretic force;
\item \textbf{emergent U(1) path completed}: duality (v1.28 rigorous) \ensuremath{\to} \textbf{Villain compactification} (Z\down{2} loop-gas defects dissolve \ensuremath{\to} compact U(1) lattice gauge theory) \ensuremath{\to} \textbf{Coulomb-phase linearized Maxwell}: photon dispersion \textbf{\ensuremath{\omega}=|k|} (\ensuremath{\omega}/k=1.0000 verified), 2 transverse polarizations\mtag{Citation, standard Villain/compact-U(1) results}; the rigorous Z\down{2}\ensuremath{\to}U(1) embedding theorem remains open (same order as a Millennium problem, not claimed);
\item \textbf{three-layer argument for spin 2}: framework layer \textbf{rigorous} (rotation-eigenvalue theorem) + string-net layer semi-rigorous (Wen 2006, collective mode = metric-type perturbation) + \textbf{observational layer: the NJU FQH "graviton mode" (chiral spin \ensuremath{\pm}2) = projection of second-order flux fluctuation modes onto a 2D electron sea}\mtag{Framework definition D-2 + observational projection} --- spin-2 emergence advanced from "world-class problem" to "mechanism clear + observational projection support", partial rigorization;
\item \textbf{testable prediction}\mtag{Framework prediction}: \textbf{FQH graviton-mode decay rate \ensuremath{\Gamma} = \ensuremath{\Gamma}\down{Coulomb} + \ensuremath{\varepsilon}\ensuremath{\cdot}\ensuremath{\Gamma}\down{flux}} (\ensuremath{\varepsilon}\ensuremath{\sim}10\up{-20} leakage term, linewidth excess) --- measurable on extremely-low-temperature high-magnetic-field platforms; standard condensed matter has no such term (\S{}9 S-E5).
\end{itemize}
\item \textbf{\mtag{v1.30 Z\down{2}\ensuremath{\to}U(1) rigorous embedding theorem and 4D spin-2 emergence (follow-up item (2), deepest; honest positioning: same order as Millennium problems, no complete proof claimed)}}\mtag{Proposition + Numerical}:
\begin{itemize}
\item \textbf{Theorem A (algebraically rigorous)}: Z\down{2}={1,\ensuremath{-}1} \ensuremath{\subset} U(1) discrete-subgroup embedding (multiplicatively closed \checkmark, injective homomorphism \checkmark) \ensuremath{\implies} \textbf{Z\down{2} string-net gauge sector \ensuremath{\subset} compact U(1) gauge field};
\item \textbf{Proposition B (analytic derivation)}: compact-U(1) lattice gauge cos(\ensuremath{\theta}\down{p}) = 1\ensuremath{-}\ensuremath{\theta}\up{2}/2+\ensuremath{\cdots} (error 2\ensuremath{\times}10\up{-5}) linearization \ensuremath{\implies} S\down{eff} = \ensuremath{\int}(1/4g\up{2})F\up{2} = \textbf{Maxwell};
\item \textbf{Numerical verification C}: Coulomb-phase photon lattice dispersion \ensuremath{\omega}(k)=2\ensuremath{\sqrt{\Sigma{}sin\up{2}(k/2)}}, long-wavelength \textbf{\ensuremath{\omega}/|k|=0.9999} (speed of light c=1 \checkmark);
\item \textbf{spin-2 criterion chain}: (1) Fierz--Pauli uniqueness (citation: massless spin 2 \ensuremath{\implies} linearized GR) (2) rotation eigenvalues \ensuremath{\pm}2 (rigorous, v1.29) (3) Weinberg soft-graviton theorem (cited criterion, soft limit \ensuremath{\propto}1/q measurable) (4) string-net low-energy spectrum (semi-rigorous, Wen) --- \textbf{the only gap = Z\down{2}\ensuremath{\to}U(1) continuum-limit RG / 4D massless spin-2 field-theory emergence (world-class, honestly not claimed)}.
\end{itemize}
\end{itemize}

\begin{itemize}
\item \textbf{\mtag{v1.34 string-net low-energy spectrum and spin-2 criterion evidence (action-list P1-1, criterion (4) numericalized)}}\mtag{Numerical}: use \textbf{chain-axis fourfold rotational symmetry (C4)} to decompose the Z\down{2} string-net low-energy spectrum into three helicity sectors ([R,H]=0, R\up{4}=I exactly verified \checkmark): scalar (C4=+1) / \textbf{helicity \ensuremath{\pm}2 (C4=\ensuremath{-}1, spin-2 criterion signature)} / vector (C4=\ensuremath{\pm}i, photon candidate). Systems: double cube (11 faces, 2,048 states, dense full spectrum) + triple-cube chain (16 faces, 65,536 states, Lanczos). (i) \textbf{\ensuremath{\lambda}=0 exact level structure fully verified}: 2-cube 512/1,024/512, 3-chain 8,192/24,576/24,576/8,192 --- consistent with the analytic 2\up{n-N\down{c}}\ensuremath{\cdot}C(N\down{c},k) \checkmark; \textbf{gapped spectrum} (gaps increase monotonically with \ensuremath{\lambda} from 0.04 to 1.92; no zero-energy collective mode); (ii) \textbf{spin-2 gap}: at \ensuremath{\lambda}=0.3, helicity-\ensuremath{\pm}2 sector \textbf{\ensuremath{\Delta}\ensuremath{\pm}2=0.4919 (double cube)/0.4864 (triple chain)} --- size-independent (1\% variation), tensor/vector modes \ensuremath{\sim}19\% heavier than scalar (\ensuremath{\Delta}\down{scalar}=0.41); (iii) \textbf{carrier evidence}: the second-order-tensor operator space Z\down{fZ}\down{g} carries C4=\ensuremath{-}1 (helicity \ensuremath{\pm}2) eigenstates \textbf{30/120-dim (3-chain)/14/55 (double cube)}, far exceeding the first-order space (3/16) --- \textbf{spin 2 = second-order-tensor sector} (operator-level alignment with \mtag{Framework definition D-2} "spin 2 = second-order flux-tensor fluctuations").\mtag{Honest conclusion} Z\down{2}-level has no massless spin 2 (Abelian gapped, expected physics --- the graviton requires the U(1) 4D continuum limit); this evidence gives the \textbf{spectroscopic positioning} of spin-2 modes in the string-net spectrum (gapped, tensor sector, size-independent), not completed emergence.
\end{itemize}

\begin{itemize}
\item \textbf{\mtag{v1.36 3+1D \ensuremath{\lambda}\down{c} transition + Wilson RG numerical path (action-list P1-3/P1-4; condensation-phase stability boundary + continuum-limit RG layer)}}\mtag{Numerical}: \textbf{(a) \ensuremath{\lambda}\down{c} transition} --- generic face construction build\down{grid}(nx,ny,nz) (chains 2/3/4\ensuremath{\times}1\ensuremath{\times}1 + plate 2\ensuremath{\times}2\ensuremath{\times}1, 2,048--2.1\ensuremath{\times}10\up{6} states, LinearOperator Lanczos) + order parameter \ensuremath{\langle}A\down{c}\ensuremath{\rangle} + \ensuremath{\chi}=\ensuremath{-}d\ensuremath{\langle}A\ensuremath{\rangle}/d\ensuremath{\lambda}: coarse scan (step 0.2) gives \ensuremath{\chi} peaks all at \ensuremath{\lambda}=0 (grid artifact); \textbf{fine scan (step 0.02) locks the \ensuremath{\chi} peaks of all four sizes at \ensuremath{\lambda}\down{\chi{}}=0.140} (\ensuremath{\chi}\down{max} 2.24--2.28 nearly constant) \ensuremath{\implies} \textbf{3+1D Z\down{2} string-net condensation\ensuremath{\to}polarization transition point \ensuremath{\lambda}\down{c}\ensuremath{\approx}0.14} (pseudo-critical value; condensed phase \ensuremath{\lambda}<0.14 stable); \textbf{dimension--critical-coupling decreasing sequence}: 3+1D string-net \ensuremath{\lambda}\down{c}\ensuremath{\approx}0.14 vs 2+1D toric code \ensuremath{\kappa}\down{c}\ensuremath{\approx}0.25--0.32 (P0-3b) vs Ising \ensuremath{\beta}\down{c} 2D/3D=0.44/0.22 --- higher dimension \ensuremath{\implies} smaller critical coupling (stronger fluctuations, easier decondensation); \ensuremath{\lambda}\down{c}\ensuremath{\approx}0.14 is of the same order as the dual 3D Ising \ensuremath{\beta}\down{c}=0.2217 (v1.28 duality self-consistent);\mtag{Framework mechanism: transition = flux-field turbulent-state \ensuremath{\leftrightarrow} isotropic-steady-state condensation point}. \textbf{(b) Wilson RG path} (Z\down{2}\ensuremath{\to}U(1) continuum-limit dynamics) --- (1) U(1) free-field (Gaussian) partition-function RG invariance (exact det spectral decomposition): scaling law \textbf{\ensuremath{\beta}\ensuremath{^\prime}=\ensuremath{\beta}\ensuremath{\cdot}b\up{2-d}} (d=2 critical scale-invariant, d=3 g*=0 Gaussian fixed point IR-free, d=1 flows to strong coupling) --- \textbf{Maxwell = Gaussian fixed point g*=0}; (2) \textbf{1D Ising exact decimation}: \ensuremath{\beta}\ensuremath{^\prime}=atanh(tanh\up{2}\ensuremath{\beta}), unique fixed point \ensuremath{\beta}*=0 (flows to Gaussian/continuum limit \checkmark); (3) \textbf{2D Ising block RG (Metropolis + majority rule, 16\ensuremath{\to}8\ensuremath{\to}4) --- core demonstration of the Z\down{2}\ensuremath{\to}U(1) path}: \ensuremath{\beta}=0.3 (<\ensuremath{\beta}\down{c}=0.44) \ensuremath{\implies} 0.193\ensuremath{\to}0.162\ensuremath{\to}\textbf{0.063 (flows to 0 = U(1) continuum limit)}; \ensuremath{\beta}=0.5 (>\ensuremath{\beta}\down{c}) \ensuremath{\implies} 0.694\ensuremath{\to}0.677\ensuremath{\to}\textbf{1.900 (flows to \ensuremath{\infty} = condensed phase)} --- the Z\down{2} string-net discretization RG-flows to the U(1) Gaussian fixed point for \ensuremath{\beta}<\ensuremath{\beta}\down{c}.\mtag{Honest note} lattice block-averaging produces next-nearest-neighbor corrections (partition-function invariance drifts 0.35--0.61 nonzero; nearest-neighbor correlations diluted within blocks --- real-space RG approximation; \ensuremath{\beta}\ensuremath{^\prime}=\ensuremath{\beta}b\up{2-d} is the continuum-limit leading term); the Gaussian model is a free field (trivial RG; true U(1) requires non-perturbative HOTRG); \ensuremath{\lambda}\down{c} is a small-system pseudo-critical value (\ensuremath{\lambda}\down{c}(\ensuremath{\infty}) requires a 3D box = MPS/DMRG). Four-layer closure with v1.30: Theorem A (group embedding) \ensuremath{\to} Proposition B (Maxwell linearization) \ensuremath{\to} Numerical verification C (photon \ensuremath{\omega}=|k|) \ensuremath{\to} RG flow (this work) --- "Z\down{2}\ensuremath{\to}U(1) continuum limit" upgraded from a single open item to four-layer evidence; connection to v1.34: the 3+1D condensed phase (\ensuremath{\lambda}<\ensuremath{\lambda}\down{c}) is the RG-semantic "basin of attraction of the Gaussian fixed point".
\end{itemize}

\textbf{Two-path cross-validation}: the statistical-mechanics path of gravitational emergence (\S{}4, complete \checkmark) and the string-net path (emergent spin 2, semi-rigorous \textbf{!}; condensation/duality layers rigorized in v1.28) --- \textbf{two independent paths point to "gravity = emergent"} --- if equivalent, string-net collective modes = statistical fluctuations of the flux field (pictorial correspondence, not a rigorous proof).


\section{6 Dynamics (Least Action)}

\subsection{6.1 Action functional}

\mtag{Derivation} Minimal Lagrangian density (multi-representation version):

\[\mathcal{L} = \sqrt{-g}\Big[\frac{1}{16\pi G}R - \tfrac14 F_{\mu\nu}F^{\mu\nu} - \tfrac12(\partial\Phi)^2 - V(\Phi) + i\bar\psi\gamma^\mu\nabla_\mu\psi - m\bar\psi\psi - g\Phi\bar\psi\psi - eA_\mu\bar\psi\gamma^\mu\psi\Big]\]

\begin{itemize}
\item Pinning term g\ensuremath{\Phi}\ensuremath{\bar{\psi}}\ensuremath{\psi}: environment-dependent effective mass m\down{eff}(x) = m + g\ensuremath{\Phi}(x) --- "the particle is pinned by the storm";
\item \ensuremath{\Phi} is the condensate representative of annihilation--regeneration noise (auxiliary field).
\end{itemize}

\subsection{6.2 Field equations (variational derivation)}

\begin{table}[htbp]
\centering
\small
\begin{tabularx}{\textwidth}{@{}lXX@{}}
\toprule
Channel & Field equation & Correspondence \\
\midrule
scalar & \ensuremath{\square}\ensuremath{\Phi} = V\ensuremath{^\prime}(\ensuremath{\Phi}) + g\ensuremath{\bar{\psi}}\ensuremath{\psi} & \textbf{concrete realization of D[\ensuremath{\Phi}]=S} \\
vector & \ensuremath{\partial}\down{\mu{}F}\up{\mu{}\nu{}} = eJ\up{\nu{}} & Maxwell \\
tensor & G\down{\mu{}\nu{}} = 8\ensuremath{\pi}G T\down{\mu{}\nu{}} & Einstein \\
\bottomrule
\end{tabularx}
\end{table}

\textbf{Degeneracy checks}: g=0, V=\ensuremath{\tfrac12}m\up{2}\ensuremath{\Phi}\up{2} \ensuremath{\to} free Dirac + scalar + GR; \ensuremath{\psi}=0, \ensuremath{\Phi}=const \ensuremath{\to} GR + cosmological constant; weak-field low-velocity \ensuremath{\to} g\down{00} \ensuremath{\approx} \ensuremath{-}(1+2\ensuremath{\Phi}\down{N}/c\up{2}).

\subsection{6.3 Noether conservation laws}

\mtag{Derivation} time translation \ensuremath{\to} energy, spatial translation \ensuremath{\to} momentum, rotation \ensuremath{\to} angular momentum, internal U(1) \ensuremath{\to} charge (4 conservation laws); local U(1) absent (single scalar without A\down{\mu{}}; completed in the multi-representation version).

\subsection{6.4 Spontaneous symmetry breaking (\ensuremath{\Phi}\up{4} thermal transition)}

\[V(\Phi) = -\tfrac12\mu^2\Phi^2 + \tfrac{\lambda}{4}\Phi^4,\qquad V_T = V(\Phi) + \tfrac12 cT^2\Phi^2\]

\begin{itemize}
\item T > T\down{c} = \ensuremath{\mu}/\ensuremath{\sqrt{c}}: \ensuremath{\langle}\ensuremath{\Phi}\ensuremath{\rangle}=0 (unified phase); T < T\down{c}: Mexican hat (broken);
\item complex-scalar upgrade \ensuremath{\to} Anderson--Higgs \ensuremath{\to} W/Z masses = "draft-pipe collective mode".
\end{itemize}

\textbf{Physical mapping}: "Higgs = flux-field phase transition" \ensuremath{\iff} order parameter \ensuremath{\langle}\ensuremath{\Phi}\ensuremath{\rangle} rising from 0 to v; "breaking = mode-boundary emergence" \ensuremath{\iff} temperature (flux density) dropping below T\down{c}.

\subsection{6.5 Geodesic equation}

\mtag{Citation} Variating the point-particle action S\down{pp} = \ensuremath{-}m\ensuremath{\int}d\ensuremath{\tau} gives the \ensuremath{\Gamma} equation; Newtonian limit \ensuremath{\vec{a}} = \ensuremath{-}\ensuremath{\nabla}\ensuremath{\Phi}\down{N}. "The body is readjusted at each instant to minimize the pressure difference" \ensuremath{\iff} \ensuremath{\Phi}\down{N} = flux negative-pressure potential.



\section{7 Foundations of Quantum Mechanics}

\subsection{7.1 The problem: classical event flow cannot yield Bell violation}

\mtag{Citation + Derivation} A classical hidden-variable flow field can only give CHSH \ensuremath{\leq} 2, contradicting the experimental CHSH = 2\ensuremath{\sqrt{2}} --- \textbf{the framework must be non-commutativized}. The original work was already aware of this ("for this to survive, you must be able to derive 2\ensuremath{\sqrt{2}} \ldots a structural product of source events").

\subsection{7.2 Non-commutative event flow}

\mtag{Derivation}\textbf{Proposition (non-commutative event flow \ensuremath{\Rightarrow} quantum correlations)}: upgrade the event-flow state from a classical probability distribution (commutative algebra L\up{\infty{}}) to an operator algebra (C* algebra), with "encounter--annihilation" events corresponding to non-commuting projection operators (Clifford-type {\ensuremath{\Pi}\down{\mu{}}, \ensuremath{\Pi}\down{\nu{}}}\ensuremath{\neq}0). Then:
(i) joint measurement of entangled pairs gives CHSH = 2\ensuremath{\sqrt{2}} (Tsirelson bound);
(ii) complementarity V\up{2} + I\up{2} \ensuremath{\leq} 1 holds automatically;
(iii) uncertainty \ensuremath{\Delta}A\ensuremath{\cdot}\ensuremath{\Delta}B \ensuremath{\geq} \ensuremath{\tfrac12}|\ensuremath{\langle}[A,B]\ensuremath{\rangle}| holds automatically.

\textbf{Non-commutative elements already in the framework}: Clifford algebra (micro-variations do not commute), path integral (\ensuremath{-}1)\up{N\down{exchange}} (exchange sign), geometric origin of superposition states (global equivalence of the centers of positive/anti storms), quantum lines (events of common origin).

\textbf{Numerical verification}: Bell state |\ensuremath{\psi}\up{-}\ensuremath{\rangle} + optimal measurement basis \ensuremath{\implies} S = \ensuremath{-}2.8284 = \ensuremath{-}2\ensuremath{\sqrt{2}} \checkmark; path-state inner product gives V\up{2}+I\up{2} = 1 \checkmark.

\subsection{7.3 Quantum walk: dynamics of the non-commutative event flow}

\mtag{Derivation} Discrete-time quantum walk U = S(C\ensuremath{\otimes}I): coin flip C (Hadamard, non-commutative) + conditional shift S. Numerics: quantum \ensuremath{\sigma} \ensuremath{\sim} t (ballistic) vs classical \ensuremath{\sqrt{t}} (diffusive) --- non-commutativization preserves coherence.

\subsection{7.4 Positioning clarification}

\mtag{Derivation} Once non-commutativized, the framework's quantum part \textbf{mathematically coincides} with standard QM. The framework's unique contribution = providing an \textbf{ontological motivation for "why quantum mechanics"} (in the spirit of Gel'fand--Naimark: commutative C* \ensuremath{\iff} classical probability; non-commutative C* \ensuremath{\iff} quantum states). The positioning is \textbf{'t Hooft cellular-automaton-like}: underlying deterministic event flow, non-commutative observable algebra.



\section{8 Cosmology}

\subsection{8.1 Energy-scale dependence of the leakage mechanism (\ensuremath{\varepsilon}\down{A})}

\mtag{Derivation} \ensuremath{\varepsilon}\down{A}(E) = \ensuremath{\varepsilon}\down{0}(M\down{Pl}/E)\up{\beta{},} two boundary conditions:
\begin{itemize}
\item E \ensuremath{\sim} 1 TeV \ensuremath{\to} \ensuremath{\varepsilon} \ensuremath{\sim} 6\ensuremath{\times}10\up{-10} (baryon asymmetry \ensuremath{\eta} \ensuremath{\approx} \ensuremath{\varepsilon});
\item E \ensuremath{\sim} 1 meV \ensuremath{\to} \ensuremath{\varepsilon} \ensuremath{\sim} 1 (dark energy \ensuremath{\rho}\down{\Lambda{}} \ensuremath{\sim} (meV)\up{4});
\item \ensuremath{\beta} = log(6\ensuremath{\times}10\up{-10})/log(10\up{-15}) \ensuremath{\approx} \textbf{0.615}, \ensuremath{\varepsilon}\down{0} \ensuremath{\approx} 7.6\ensuremath{\times}10\up{-20}.
\end{itemize}

\textbf{Cutoff scale}\mtag{Derivation}: probability normalization \ensuremath{\varepsilon}\down{A}(E\down{c}) = 1 \ensuremath{\implies} E\down{c} = M\down{Pl}\ensuremath{\cdot}\ensuremath{\varepsilon}\down{0}\up{1/\beta{}} = \textbf{1.01 meV} --- of the same order as the dark-energy characteristic scale 2.3 meV ((E\down{c})\up{4} differs by a factor 27; O(1) coefficient undetermined).

\subsection{8.2 Statistical decay factor (\ensuremath{\varepsilon}\down{B}) and dark matter}

\mtag{Derivation} Planck-fluctuation raw density \ensuremath{\rho}\down{raw} = 2m\down{P}/l\down{P}\up{3} \ensuremath{\approx} 10\up{97} kg/m\up{3}; statistical decay \ensuremath{\varepsilon}\down{B} = (l\down{P}/R\down{U})\up{2} \ensuremath{\approx} 10\up{-120} (geometric origin, not (1/3)\up{n} --- n is actually 251); dark matter \ensuremath{\rho}\down{DM} = \ensuremath{\varepsilon}\down{B}\ensuremath{\cdot}\ensuremath{\rho}\down{raw} \ensuremath{\approx} 1.6\ensuremath{\times}10\up{-25} kg/m\up{3} (differs from the observed 2.3\ensuremath{\times}10\up{-27} by \ensuremath{\sim}70\ensuremath{\times} --- \textbf{\mtag{To be verified}} requires f\down{space} \ensuremath{\approx} 0.014 or mechanism correction).

\textbf{Topological origin of \ensuremath{\varepsilon}\down{B} (H1--H4)}: quantification of "\ensuremath{\varepsilon} is the topologically forced 'minimal vortex-core density'"\mtag{Proposition chain H1--H4, based on Poincar\'{e}--Hopf}:
\begin{itemize}
\item \textbf{H1 (S\up{2} rigorous)}: \ensuremath{\chi}(S\up{2}) = 2 \ensuremath{\implies} sum of zero indices of a tangent vector field = 2 \ensuremath{\implies} direction space has at least 2 vortex cores\mtag{Topological argument};
\item \textbf{H2 (vortex-core fraction)}: \ensuremath{\varepsilon}\down{topo}(R) = (l\down{P}/R)\up{2}/(2\ensuremath{\pi}) --- probability that an event direction falls in a vortex-core cell;
\item \textbf{H3 (component structure)}: \ensuremath{\varepsilon}\down{total} = \ensuremath{\varepsilon}\down{topo} + \ensuremath{\varepsilon}\down{A} --- topological dominance at Planck scale (0.159), dynamical dominance above TeV (\ensuremath{\varepsilon}\down{A}); \ensuremath{\varepsilon}\down{0} = 7.67\ensuremath{\times}10\up{-20} is the dynamical-component anchor;
\item \textbf{H4 (significant finding)}: \ensuremath{\varepsilon}\down{B}(cosmos) = 1.35\ensuremath{\times}10\up{-123} \ensuremath{\approx} \ensuremath{\varepsilon}\down{topo}(cosmos)\ensuremath{\times}2\ensuremath{\pi} --- \textbf{the dark-matter decay factor lies exactly at the order of the topologically forced minimum; the geometric origin of \ensuremath{\varepsilon}\down{B} gains topological endorsement};
\item \textbf{S\up{3} clarification\mtag{Correction}}: \ensuremath{\chi}(S\up{3})=0 and S\up{3} is parallelizable (trivial tangent bundle); the hairy-ball theorem does not apply to 3D spatial slices (early-version correction) --- vortex cores in spatial slices require Kakeya + leakage supplementation.
\end{itemize}

\subsection{8.3 Dark-matter superposition model}

\mtag{Derivation}\textbf{Dark matter = this-branch Planck-fluctuation remnants (\ensuremath{\varepsilon}\down{B} mechanism) + mirror-branch gravitational projection (cross-branch mechanism)}:

\[\rho_{DM} = \varepsilon_B\,\rho_{raw}\,f_{space} + \kappa_{mirror}\,\rho_{mirror}\]

\begin{itemize}
\item Mirror-branch projection: \ensuremath{\kappa}\down{extra} \ensuremath{\sim} \ensuremath{\Omega}\down{mirror}/\ensuremath{\Omega}\down{total} \ensuremath{\approx} 0.5 (gravity shared across branches; EM/strong/weak filtered by the branch boundary) \ensuremath{\implies} \ensuremath{\rho} \ensuremath{\sim} 4\ensuremath{\times}10\up{-27} kg/m\up{3} \textbf{same order as observation} --- \textbf{the bulk of dark matter may come from mirror-branch gravitational projection}\mtag{To be verified};
\item Testable prediction: \textbf{void-CPT-\ensuremath{\kappa} anomaly} (void regions with stronger weak lensing + mass--light decoupling, most prominent at 10--30 Mpc/h scales; candidate objects CDG-2, J0613+52).
\end{itemize}

\subsection{8.4 Dark energy and w(z)}

\mtag{Derivation} Dark energy = contribution of the \ensuremath{\varepsilon}\down{A} saturation region \ensuremath{\rho}\down{\Lambda{}} \ensuremath{\sim} (E\down{c})\up{4}; the saturation scale drifts slowly with cosmic expansion E\down{c}(z) = E\down{c},0(1+z)\up{\varepsilon{}} \ensuremath{\implies}

\[w = -1 + \frac{4\varepsilon}{3} = -1 + 2\nu\]

Observational constraint (DESI 2024+CMB: w = \ensuremath{-}1.028\ensuremath{\pm}0.031): |\ensuremath{\varepsilon}| < 0.023, |\ensuremath{\nu}| < 0.0155. \textbf{The framework predicts w = const (wa=0)} --- directly testable with DESI DR2/Euclid. \textbf{\mtag{v2.7 supplement: DESI DR2 (2025-03) discrimination status (referee M5)}}: DESI DR2+CMB exclude \ensuremath{\Lambda}CDM at \textbf{3.1\ensuremath{\sigma}} (+Pantheon+ 2.8\ensuremath{\sigma} / +Union3 3.8\ensuremath{\sigma} / +DESY5 4.2\ensuremath{\sigma}), with CPL best fit \textbf{w\down{0}\ensuremath{\approx}\ensuremath{-}0.76\ensuremath{\pm}0.06, wa\ensuremath{\approx}\ensuremath{-}0.77\ensuremath{\pm}0.23} --- exerting \textbf{\ensuremath{\sim}3.3\ensuremath{\sigma} pressure on the framework's wa=0 prediction} (the framework's dark-energy prediction is being discriminated by \textbf{existing data}: if the DDE signal is confirmed \ensuremath{\implies} this prediction is falsified; if a systematic effect \ensuremath{\implies} the framework survives). \textbf{Verdict timeline}: DESI final analysis 2026-04 + Euclid first cosmological data 2026-10 --- \textbf{final verdict within 0-2 years}. Honest note: the DDE-vs-systematics debate is unresolved (ACT DR6/SPT-3G CMB support \ensuremath{\Lambda}CDM); the w\down{0} discrimination is a cross-model wCDM/CPL comparison (not directly comparable; the wa=0 test is authoritative).

\subsection{8.5 Record of original error corrections}

\mtag{Corrections} (1) "meV scale = Hubble scale" correspondence error (off by \ensuremath{\sim}30 orders of magnitude; meV wavelength 1.24 mm); (2) \ensuremath{\rho}\down{raw} dimensional typo (kg/(m\up{3}\ensuremath{\cdot}s) should be kg/m\up{3}); (3) "120" origin = (l\down{P}/R\down{U})\up{2}, not (1/3)\up{120.}


\section{9 Testable Predictions}

\begin{table}[htbp]
\centering
\small
\begin{tabularx}{\textwidth}{@{}lXXX@{}}
\toprule
\# & Prediction & Observational window & Status \\
\midrule
1 & \textbf{B1--B5 black-hole strong-field prediction package}: extreme-Kerr Hawking-radiation angular distribution deviates from axisymmetry (B1); primordial black-hole entropy volume correction \ensuremath{\Delta}S/S\ensuremath{\sim}l\down{P}/R\ensuremath{\sim}10\up{-8} (B2); Hawking-radiation--infalling-matter entanglement (B3, information-flip verdict signature); ER-bridge packing bound \ensuremath{\to} BBH angular-mode distortion (B4); non-thermal correction \ensuremath{\propto}1/S\down{BH} (B5) & EHT/LIGO, primordial black-hole \ensuremath{\gamma}-rays, LISA era, ET/CE & theoretically testable (see \S{}4.10) \\
2 & \textbf{void-CPT-\ensuremath{\kappa} anomaly} (void regions with stronger weak lensing, mass--light decoupling) & DES/Y3, DESI, Euclid, Rubin LSST & testable (analysis scheme given) \\
3 & \textbf{constant deviation w = \ensuremath{-}1+4\ensuremath{\varepsilon}/3 (wa=0)} & \textbf{DESI DR2 (existing, 2025-03)}, Euclid 2026-10 & \textbf{under discrimination: wa=0 under 3.3\ensuremath{\sigma} pressure (DDE vs systematics unresolved), final verdict 2026} \\
4 & \textbf{SU(3) \ensuremath{\beta}-function staircase jump} (fourth-generation threshold) & FCC-hh, LHAASO/CTA & theoretically computable, numerics to be derived \\
5 & traces of a fourth fermion generation (>1 TeV, decays at z>1000) & LHC, cosmic rays, CMB-S4 & tension with LHC non-detection \\
6 & \textbf{Hawking-radiation non-thermal correction} \ensuremath{\sim}1/S (primordial black holes \ensuremath{\sim}10\up{-20}) & primordial black-hole evaporation & near window edge \\
6\ensuremath{^\prime} & \textbf{firewall-thickness non-thermal correction} \ensuremath{\sim}\ensuremath{\sqrt{l\down{P}/R}} (primordial black holes \ensuremath{\sim}10\up{-9}\ensuremath{\sim}10\up{-4}, F4 emission-layer prediction) & primordial black-hole \ensuremath{\gamma}-rays & theoretically testable \\
7 & PeV neutrino flavor-ratio anomaly & IceCube-Gen2, GRAND & qualitative prediction \\
8 & entanglement-entropy area-law coefficient (framework contains ln 2 factor) & quantum simulation / many-body systems & theoretically testable \\
9 & \textbf{topological-order excitations (quantum double / string-net)}: torus ground-state degeneracy \ & G\ & \up{2}, anyonic fractional statistics, string-endpoint excitation gap \ensuremath{\Delta}\ensuremath{\sim}1 (\S{}5.6--5.8) \\
10 & \textbf{confined string-tension numerics + transition-point localization (v1.20--v1.23)}: charge pair \ensuremath{\sigma}\down{E}\ensuremath{\approx}2\ensuremath{\kappa} (X field, thermodynamic-limit convergence <0.02\%) and flux pair \ensuremath{\sigma}\down{m}\ensuremath{\approx}2\ensuremath{\lambda} (Z field, constrained sector: 0.21/1.09/3.04 at \ensuremath{\lambda}=0.5/1.0/2.0); \ensuremath{\lambda}=0/\ensuremath{\kappa}=0 strictly deconfined; \textbf{transition localization (P0-3/P0-3b)}: open-boundary FSS 0.23--0.27 \ensuremath{\to} torus half-polarization point \ensuremath{\kappa}*\ensuremath{\approx}0.20\ensuremath{\pm}0.01 (size-independent) \ensuremath{\to} true transition \textbf{\ensuremath{\kappa}\down{c}\ensuremath{\approx}0.25--0.32} (consistent with literature X-field value h\down{xc}\ensuremath{\approx}0.32); string-tension numerics in the confined phase \checkmark & quantum simulation (Z\down{2} lattice-gauge implementation), superconducting circuits & numerically reproducible (scripts \down{scale}\down{extrap3}/\down{p01}\down{sector}/\down{p03}\down{fss}/\down{p03}\down{torus}*) \\
\textbf{S1} & \textbf{\ensuremath{\alpha} = 1.00\ensuremath{\pm}0.06} (storm-overload shadow compatible with GR, v1.15) & EHT M87*/Sgr A* & \checkmark already compatible \\
\textbf{S2} & \textbf{QNM frequency f = f\down{GR}/\ensuremath{\alpha}} (deviation \ensuremath{\pm}5--11\% for \ensuremath{\alpha}\ensuremath{\neq}1) & ET/CE future & testable (\ensuremath{\pm}0.01) \\
\textbf{S3} & \textbf{firewall spread \ensuremath{\delta}/R \ensuremath{\sim} 10\up{-9} (primordial black holes, framework-unique)} & primordial black-hole evaporation/gravitational waves & framework-unique signature \\
\textbf{S4} & \textbf{overload region without singularity + sound-speed vanishing + optical transparency} (pinned EOS P=P\down{P}\ensuremath{\cdot}tanh(\ensuremath{\rho}/\ensuremath{\rho}\down{P}), v1.24 micro-derivation: P\ensuremath{\leq}\ensuremath{\rho}c\up{2} everywhere, \ensuremath{\partial}P/\ensuremath{\partial}\ensuremath{\rho}\ensuremath{\geq}0 no collapse; c\down{s}\up{2}=c\up{2}\ensuremath{\cdot}sech\up{2}(\ensuremath{\rho}/\ensuremath{\rho}\down{P}), no elastic waves for \ensuremath{\rho}\ensuremath{\geq}\ensuremath{\rho}\down{P}; v1.25 uniform paste \ensuremath{\bar{\rho}}=3c\up{6}/(32\ensuremath{\pi}G\up{3}M\up{2}) + optical depth \ensuremath{\tau}\ensuremath{\sim}10\up{-40} (macroscopic)/0.06 (Planck window); \textbf{v1.31 accretion-disk radiative transfer: standard thin-disk bulk consistent with GR + overload region fully transparent + firewall non-thermal correction (primordial window)} --- framework-unique) & accretion flow/overload-region radiation, Planck-mass black-hole evaporation, primordial black-hole \ensuremath{\gamma}-rays & \checkmark fluid model + micro-derivation + radial structure + disk transfer (\S{}4.6) \\
\textbf{S5} & shadow-edge Planck spread (macroscopic \ensuremath{\sim}10\up{-19}\ensuremath{\sim}10\up{-24} unmeasurable) & ultra-long-baseline interferometry & theoretical prediction \\
\textbf{S-E1} & \textbf{EMRI perihelion precession rate higher \ensuremath{\delta}=0.30\ensuremath{\cdot}x\down{p}} (1--3\%, 2PN g\down{rr} 33\% bifurcation signature; v1.26 proposal + v1.27 Fisher + v1.32 Kerr Bayesian + \textbf{v1.35 a-\ensuremath{\delta} joint}) & LISA EMRI (M\ensuremath{\sim}10\up{5}--10\up{7} M\ensuremath{\odot}) & \textbf{after v1.35 correction: \ensuremath{\Delta}lnL criterion 27\ensuremath{\sigma} (2yr) / 28\ensuremath{\sigma} (4yr), robust with a-\ensuremath{\varphi}0 fully free (\ensuremath{\sigma}\down{\delta{}}=1.6\ensuremath{\times}10\up{-6}; v1.32's 9.7\ensuremath{\times}10\up{-4} was undersampling-inflated by 600\ensuremath{\times}) --- single-source verdict} \\
\textbf{S-E2} & \textbf{EMRI gravitational-wave phase deviation from GR} (precession-modulated phase difference \ensuremath{\sim}2\ensuremath{\times}\ensuremath{\Delta}\ensuremath{\Phi}\down{prec}) & LISA & testable \\
\textbf{S-E3} & \textbf{EMRI orbital decay / chirp phase difference} (3PN g\down{00} 11\% bifurcation) & LISA & testable (\ensuremath{\Delta}\ensuremath{\Phi}\down{chirp}\ensuremath{\sim}2\ensuremath{\times}10\up{2} rad) \\
\textbf{S-E4} & \textbf{precession-difference--radius scaling law} \ensuremath{\delta}\ensuremath{\propto}x\down{p} (0.6\% @ p=50M \ensuremath{\to} 3.0\% @ p=10M) --- framework-unique, multi-source joint test & LISA & testable \\
\textbf{S-E5} & \textbf{FQH graviton-mode decay rate contains \ensuremath{\varepsilon} term} (\ensuremath{\Gamma}=\ensuremath{\Gamma}\down{Coulomb}+\ensuremath{\varepsilon}\ensuremath{\cdot}\ensuremath{\Gamma}\down{flux}, \ensuremath{\varepsilon}\ensuremath{\sim}10\up{-20} leakage; framework prediction, v1.29) & extremely-low-temperature high-field FQH platforms (linewidth excess) & testable (framework-unique) \\
\textbf{S-E6} & \textbf{glueball = non-Abelian flux self-bound state} (framework translation[BESIII 2026]; BESIII X(2370) confirmed: 0\up{-+}, 2395 MeV, flavor singlet, v1.33 anchor) & BESIII/BEPCII, J/\ensuremath{\psi} decays, future hadron colliders & framework-level interpretation (QCD fact + framework translation) \\
\bottomrule
\end{tabularx}
\end{table}



\section{10 Conclusions and Open Problems}

\subsection{10.1 Assessment of the framework's status}

Starting from the four axioms, this framework provides an \textbf{emergent statistical-mechanical argument} for the unification of the four fundamental forces, with completeness \ensuremath{\sim}90\%:

\begin{itemize}
\item \textbf{Gravity side} (most complete): \ensuremath{\gamma}=1, \ensuremath{\beta}=1 all passing (PPN two-parameter verdicts); micro-variation-weight proposition; Boltzmann distribution (thermal-equilibrium proposition); Jacobson emergence path; S=A/4 microstate counting (k=2\up{1/4}); \textbf{packing proposition package} (new proof of the k\ensuremath{\geq}4 lower bound + Tammes regular tetrahedron + geometric origin of the 1/4 coefficient); \textbf{Proposition 6} (interior-space encoding priority, framework-internal deduction); \textbf{observer model O1--O4} (mutual-information framework + rigorous uniform prior); \textbf{present-plane selection N1--N4} (vortex dynamics + symmetry breaking); \textbf{initial-perturbation origin P1--P4} ("necessary contingency", sharing one origin with the cosmology side) --- \textbf{"gravity = statistical-mechanical emergence of the flux field" acquires a complete argument chain from axioms to equations, from micro to observer};
\item \textbf{Gauge side} (v1.19 substantially strengthened + v1.20 numerical closure): "2 vs 3" pairing orders (Lemmas A--D rigorous at the representation-theory level); \textbf{quantum-double proposition package A1--A6} (gauge structure = discrete-algebra emergence, \S{}5.6); \textbf{framework Hamiltonian B1--B5} (emergent gauge symmetry [H\down{F},A]=[H\down{F},B]=0 rigorous + 2+1D mass gap + torus degeneracy |G|\up{2}, \S{}5.7); \textbf{confined string-tension numerics} (X-field charge \ensuremath{\sigma}\down{E}\ensuremath{\approx}2\ensuremath{\kappa} thermodynamic-limit convergence + Z-field flux \ensuremath{\sigma}\down{m}\ensuremath{\approx}2\ensuremath{\lambda} constrained sector + deconfined-degeneracy rigorous verification, \S{}5.7) \textbf{+ transition-point localization} (v1.23: torus half-polarization \ensuremath{\kappa}*\ensuremath{\approx}0.20\ensuremath{\pm}0.01, true transition \ensuremath{\kappa}\down{c}\ensuremath{\approx}0.25--0.32 consistent with literature h\down{xc}\ensuremath{\approx}0.32, P0-3b, \S{}5.7); \textbf{string-net 3+1D route} (quantum lines = string edges, four forces = string-net-order emergence, gravity two-path cross-validation, \S{}5.8) --- \textbf{the gauge side is upgraded from "structural heuristic" to "proposition package + explicit Hamiltonian + emergent-symmetry propositions + confinement numerics"};
\item \textbf{Quantum side}: non-commutativization closure (CHSH=2\ensuremath{\sqrt{2}}, complementarity); 't Hooft-like positioning;
\item \textbf{Cosmology side}: \ensuremath{\varepsilon} division of labor clear, w(z) parameterized; dark-matter superposition model (quantitative closure incomplete); initial perturbation, baryon asymmetry, and dark energy share one origin (F26).
\end{itemize}

\subsection{10.1b Numerical reliability assessment (v2.3, response to referee M3)}

\textbf{Criterion calibration (\ensuremath{\kappa}\down{c})}: fine-grid (52 points) 4-size torus data analysis --- internal criteria (m=0.5 / \ensuremath{\chi} peak / \ensuremath{\Delta} minimum) give 0.11--0.17 by 1/L extrapolation, \textbf{systematically low by 47--67\%} relative to the literature exact value h\down{c}\ensuremath{\approx}0.3285 (iPEPS) (the toy-model scale gap is quantitatively confirmed); \textbf{decisive diagnosis: m(h\down{c}) extrapolates to 0.93 \ensuremath{\neq} 0 with L \ensuremath{\implies} \ensuremath{\langle}X\ensuremath{\rangle} (single-site magnetization) is not the transition order parameter} (mechanism-level confirmation of the v1.23 "m=0.5 is a half-polarization point" criterion correction); \ensuremath{\chi}\down{max} \ensuremath{\propto} L\up{1.12} (2D-Ising-like \ensuremath{\gamma}/\ensuremath{\nu}=1.75) supports the existence of the transition, broadened by small lattices. \textbf{Calibration conclusion}: the true transition point is taken as the literature value \textbf{\ensuremath{\kappa}\down{c}\ensuremath{\approx}0.25--0.32} (framework symmetric convention equivalent to the literature, v1.23), with the criterion basis upgraded to "literature calibration + string order parameter \ensuremath{\langle}W\ensuremath{\rangle} + Binder crossing + large-lattice MPS/DMRG".

\textbf{\ensuremath{\lambda}\down{c} finite-size boundary}: 4-size \ensuremath{\chi}\down{max} \ensuremath{\propto} L\up{-0.04} (not growing) \ensuremath{\implies} the \ensuremath{\chi} peak is a broadened pseudo-peak; \ensuremath{\lambda}\down{\chi{}} shows no drift (all 0.140) \ensuremath{\implies} \textbf{conservative interval \ensuremath{\lambda}\down{c}(\ensuremath{\infty}) \ensuremath{\in} [0.140, 0.140]; precise localization requires 3D-box MPS/DMRG} (open numerical problem, honestly retained).

\textbf{EMRI systematic-error upper bound}: quadrupole waveform-truncation systematic phase error \ensuremath{\delta}\ensuremath{\Phi}\down{sys}\ensuremath{\approx}1 rad (literature order, Huerta \& Gair 2011 / Babak 2017) \ensuremath{\implies} \textbf{\ensuremath{\sigma}\down{\delta{}},tot\ensuremath{\approx}2.8\ensuremath{\times}10\up{-6} (systematic and statistical errors of the same order, \ensuremath{\sigma}\down{sys}/\ensuremath{\sigma}\down{stat}\ensuremath{\approx}1.4)} --- the parametric precision of \ensuremath{\sigma}\down{\delta{}} is limited by quadrupole (Teukolsky could reduce to \ensuremath{\sim}2\ensuremath{\times}10\up{-8}); but \ensuremath{\delta}=2\ensuremath{\times}10\up{-2} \ensuremath{\gg} \ensuremath{\sigma}\down{\delta{}},tot (7108\ensuremath{\sigma}) \ensuremath{\implies} \textbf{the 27\ensuremath{\sigma}/28\ensuremath{\sigma} framework discrimination is not threatened by systematic error}.

\textbf{\ensuremath{\to} toy-model domain of validity clarified}: transition criteria (\ensuremath{\kappa}\down{c}) depend on literature calibration and order-parameter upgrade; \ensuremath{\lambda}\down{c}(\ensuremath{\infty}) is an open numerical problem; EMRI discrimination is robust but its precision is limited by the waveform approximation (see the M3 report).

\subsection{10.2 Open problems (in order of importance)}

\begin{enumerate}
\item \sout{Dynamical proof of interior-space encoding priority} \textbf{\checkmark solved (v1.2)}: Proposition-6 three-lemma chain closed (isotropy erases exterior-space information + event point-like-ness excludes directional degrees of freedom + spinor survival carries interior-space information) --- all three premises are direct axioms, no new assumptions; \textbf{"direction discretization" assumption \checkmark solved (v1.21, B4)}: Planck tube diameter\mtag{Framework assumption: direction discretization}\ensuremath{\implies} direction cluster k=4 (Wang Hong/Tammes) \ensuremath{\implies} Bekenstein bound recovered (\S{}4.9, \S{}4.10);
\item \textbf{complete emergence of the non-Abelian components}: SU(2)\ensuremath{\times}SU(3) from the dynamics of "color--direction" coupling of the flux field (including mass gap, chirality) --- \checkmark advanced (B1: f\up{abc} topology + Jacobi = tube network + m\down{g}\ensuremath{\approx}\ensuremath{\sqrt{\sigma{}}}); \textbf{v1.33 (P0-2)}: S\down{3} non-Abelian 3+1D face-model numerics (46,656 states) --- \ensuremath{\lambda}=0 condensation E0=\ensuremath{-}N\down{c}, degeneracy 6\up{n-1}=7776, \textbf{excitations split into two conjugacy classes} (non-Abelian fingerprint) + glueball experimental anchor (BESIII X(2370), \S{}5.7); complete 4D SU(3) emergence still open;
\item \textbf{strong-field metric form}: the framework's strong field is the "overload region / flip surface" (mechanism complete, seven-loop chain) --- \textbf{v1.21 update}: shadow 5.436GM/c\up{2} of the isotropic exponential metric (\S{}4.5b) compatible with EHT (correcting v1.14's "no metric commitment" to "geometric predictions compatible with observation"); \textbf{v1.24 update}: overload-region pinned EOS micro-derived (P=P\down{P}\ensuremath{\cdot}tanh(\ensuremath{\rho}/\ensuremath{\rho}\down{P}) + sound-speed profile c\down{s}\up{2}=c\up{2}\ensuremath{\cdot}sech\up{2}(\ensuremath{\rho}/\ensuremath{\rho}\down{P}), \S{}4.6); complete derivation of horizon geometry from saturated modes still open (\S{}4.6);
\item \textbf{statistics\ensuremath{\to}geometry rigorization}: ergodicity derived from framework dynamics, chaotic-propagation applicability, quantitative cutoff --- \checkmark advanced (B2: ergodicity from Axiom 1, cutoff r\down{c}=GM/c\up{2}=horizon, B2-2 collision-kernel level); \textbf{v2.6 (referee M1)}: numerical proof of Boltzmann emergence (DSMC temperature 0.00\%/distribution 7.0\%/H theorem \ensuremath{-}2.27/ergodicity 0.00\%; hard-disk control H \ensuremath{-}2.38) --- thermal equilibrium = long-time-limit emergence, not assumption (circularity risk resolved); H\down{event} three-layer candidates (classical collision/spectral action/quantum Lindblad) constructed; LHC energy-scale compatibility argued (predicts no new LHC physics); complete annihilation-type BBGKY proof still open (honestly retained);
\item \sout{observational-limitation\ensuremath{\to}information-capacity rigorization} \textbf{\checkmark solved (v1.3)}: O2 gives the mutual-information framework S\down{access}(R) = I(R;\ensuremath{\Sigma}) (\S{}4.9); \textbf{v1.21 strengthened}: mutual-information area law I=H(\ensuremath{\Sigma}) derived from bulk = surface projection (B5, \S{}4.9);
\item \textbf{higher-order PPN parameters beyond \ensuremath{\beta}\down{ppn}}: non-linear corrections of flux occlusion --- \checkmark solved (v1.21, B3): complete PPN parameter set (\ensuremath{\gamma}=\ensuremath{\beta}=1, all others 0) + 2PN/3PN bifurcation quantified (\S{}4.5b);
\item \sout{dynamics of the two-dimensional time curling} \textbf{\checkmark solved (v1.21, B4\ensuremath{^\prime})}: "exactly two dimensions curled" = 3\ensuremath{-}1=2 definitional closure (three planes equivalent \ensuremath{\to} exactly one locked (N1--N4) \ensuremath{\to} remaining two relatively curled, no free parameter); remaining rigorization: N2's 4\ensuremath{\pi} phase-locking compatibility equation (deeply open);
\item \sout{hairy-ball topology \ensuremath{\implies} quantitative \ensuremath{\varepsilon} lower bound} \textbf{\checkmark solved (v1.7)}: proposition chain H1--H4 --- \ensuremath{\varepsilon}\down{topo}(R) = (l\down{P}/R)\up{2}/(2\ensuremath{\pi}) (Poincar\'{e}--Hopf supported); significant finding \ensuremath{\varepsilon}\down{B}(cosmos) \ensuremath{\approx} \ensuremath{\varepsilon}\down{topo}(cosmos)\ensuremath{\times}2\ensuremath{\pi} (dark-matter decay factor = topological component at cosmic scale); S\up{3} parallelizability clarification (hairy ball not applicable, early-version correction) (\S{}8.2);
\item \sout{dynamical proof of "fluctuations unavoidable"} \textbf{\checkmark solved (v1.21, B5\ensuremath{^\prime})}: three-layer structure --- statistical layer (density fluctuations \ensuremath{\sigma}/\ensuremath{\rho}\ensuremath{\sim}1/\ensuremath{\sqrt{N}} Poisson) + topological layer (hairy-ball Poincar\'{e}--Hopf \ensuremath{\chi}=2) + leakage layer (\ensuremath{\varepsilon} Poisson-random \ensuremath{\implies} fluctuation directions isotropic) --- "absolute uniformity" is not stationary at finite N (\S{}4.9, P2-remainder closed);
\item \textbf{why n=3 (gauge-side A6)}: gauge-side unification of Clifford interior-space 4 components / number-line classification / Fermat solvability (\S{}5.6 A6 --- no Tammes-type uniqueness theorem available; harder than the gravity side) --- \checkmark advanced (B6: numbering misalignment unified + five-layer argument chain); \textbf{v2.5 (referee M6)}: three-criterion structural screening of pairing order \ensuremath{\to} gauge group (S1 dimension matching d\down{min}=n / S2 same-sign pairing order n\ensuremath{^\prime} / S3 third-order antisymmetry \ensuremath{\varepsilon}\down{ijk}) --- 10 candidates eliminated one by one, unique survivor (n=1\ensuremath{\to}U(1), n=2\ensuremath{\to}SU(2), n=3\ensuremath{\to}SU(3)); \textbf{SO(3) hard-excluded because 3 is self-dual (n\ensuremath{^\prime}=2)} (3\ensuremath{\otimes}3 contains trivial \ensuremath{\to} no confinement forcing); "exactly SU(3) rather than SU(4)" gains structural-screening support (n=4 excluded by the 3D Kakeya dimension); residual contingency = micro-origin of vortex winding-mode numbers n=1,2,3 (\S{}5.2, honestly retained);
\item \textbf{origin of the J, K coupling strengths}: the coupling strengths of the framework Hamiltonian H\down{F}=\ensuremath{-}J\ensuremath{\Sigma}A\down{v}\ensuremath{-}K\ensuremath{\Sigma}B\down{p} must be derived from event-flow micro-dynamics (\S{}5.7 B1's "shape" exists, strengths not derived) --- \textbf{v1.21 (B7)}: same-origin argument (J=pairing, K=closure, both storm-flux interactions) \ensuremath{\implies} magnitude J/K\ensuremath{\sim}O(1); toric-code transition \textbf{h\down{xc}\ensuremath{\approx}0.32} (X-field/\ensuremath{\sigma}\up{x} version; 0.328 is the \ensuremath{\sigma}\up{z-field} value) \ensuremath{\implies} J/K\ensuremath{\sim}3 reference (framework side: torus half-polarization \ensuremath{\kappa}*\ensuremath{\approx}0.20, true transition \ensuremath{\kappa}\down{c}\ensuremath{\approx}0.25--0.32, P0-3b, cross-validated \checkmark); no first-principles values (still open);
\item \textbf{Hamiltonian form of the \ensuremath{\varepsilon} leakage term}: \ensuremath{-}\ensuremath{\lambda}\ensuremath{\Sigma}(\ensuremath{-}1)\up{g\neq{}e} mass-term-driven topological-order\ensuremath{\to}confinement transition --- \textbf{\checkmark landed (v1.20)}: \ensuremath{\varepsilon}-term Hamiltonian constructed and numerically studied (\S{}5.7), string tensions \ensuremath{\sigma}\down{m}/\ensuremath{\sigma}\down{E} both versions + deconfined-degeneracy rigorous verification; \textbf{\checkmark transition localization (v1.22 P0-3 open-boundary FSS; v1.23 P0-3b torus refinement)}: no finite-size transition (degeneracy immediately broken for \ensuremath{\lambda}>0, string density decreasing continuously) is a finite-lattice manifestation; torus (periodic boundary) gives half-polarization point \textbf{\ensuremath{\kappa}*\ensuremath{\approx}0.20\ensuremath{\pm}0.01} (size-independent); \textbf{criterion correction} --- m=0.5 crossing is a half-polarization point, not the transition point (torus measures m(h=0.32)\ensuremath{\approx}0.84); true transition \textbf{\ensuremath{\kappa}\down{c}\ensuremath{\approx}0.25--0.32}, consistent with the literature X-field (\ensuremath{\sigma}\up{x}) value \textbf{h\down{xc}\ensuremath{\approx}0.32} (0.328 is the \ensuremath{\sigma}\up{z-field} value, corrected in v1.23; exact value awaits MPS/DMRG, \S{}5.7);
\item \textbf{3+1D string-net rigorization}: general theorem of emergent Maxwell/GR (\S{}5.8 semi-rigorous\ensuremath{\to}rigorous, stage-C remainder) + \textbf{4D YM mass gap} (= Millennium problem, stage D, not claimed) --- \textbf{v1.21 (B8)}: candidate roadmap located (quantum-line = string-net-edge mapping rigorization), 4D not claimed; \textbf{v1.28 (follow-up item (2))}: (1) quantum-line--string-net-edge Hopf bridge landed; (2) \textbf{rigorous proof of \ensuremath{\lambda}=0 string-net condensation} (E0=\ensuremath{-}N\down{c}, degeneracy 2\up{n-N\down{c}}, \ensuremath{\Delta}=2); (3) rigorous duality mapping; (4) emergent U(1) path made explicit (duality rigorous, continuum limit open); \textbf{v1.29}: emergent U(1) path completed (Villain compactification + Coulomb-phase Maxwell \ensuremath{\omega}=|k|, literature citation); \textbf{spin 2 = rigorous verification of second-order-tensor rotation eigenvalues \ensuremath{\pm}2} (framework spin-2 mechanism) + three-layer argument (framework rigorous / string-net semi-rigorous / NJU FQH projection evidence) + \textbf{FQH \ensuremath{\varepsilon}-term prediction} (\S{}9 S-E5); \textbf{v1.30}: \textbf{Theorem A (rigorous)} Z\down{2} \ensuremath{\subset} U(1) group embedding + \textbf{Proposition B (analytic)} cos linearization = Maxwell + \textbf{Numerical verification C} photon \ensuremath{\omega}=|k| (long-wavelength 0.9999) + \textbf{spin-2 criterion chain} (Fierz--Pauli + rotation eigenvalues \ensuremath{\pm}2 + Weinberg soft theorem + string-net low-energy spectrum); \textbf{v1.34 (P1-1)}: \textbf{criterion (4) numericalized} --- C4 helicity decomposition (scalar/helicity-\ensuremath{\pm}2/vector three sectors), spin-2 mode gapped (\ensuremath{\Delta}\ensuremath{\pm}2\ensuremath{\approx}0.49 size-independent), second-order-tensor sector carries 30 helicity-\ensuremath{\pm}2 eigenstates (aligned with \mtag{Framework definition D-2}); \textbf{v1.36 (P1-3/P1-4)}: \textbf{3+1D \ensuremath{\lambda}\down{c}\ensuremath{\approx}0.14} (four-size \ensuremath{\chi}-peak lock, dimension--critical-coupling decreasing sequence with 2D toric code) + \textbf{Wilson RG path} (U(1) Gaussian-fixed-point scaling law \ensuremath{\beta}\ensuremath{^\prime}=\ensuremath{\beta}b\up{2-d} + 2D Ising block RG flowing to the continuum limit for \ensuremath{\beta}<\ensuremath{\beta}\down{c} --- Theorem A + Proposition B + Numerical verification C supplemented by the 4th-layer RG-flow evidence, four-layer closure); \textbf{remaining open}: Z\down{2}\ensuremath{\to}U(1) continuum-limit RG rigorization (world-class, RG numerical path demonstrated), 4D massless spin-2 field-theory emergence (world-class), \ensuremath{\lambda}\down{c}(\ensuremath{\infty}) exact value (requires 3D-box MPS/DMRG), rigorous Maxwell-spectrum verification.
\end{enumerate}

\subsection{10.3 Relation to mainstream physics}

\begin{itemize}
\item \textbf{Does not overthrow GR/QM}: the framework positions itself as the "statistical-mechanical foundation" of both (GR is the geometric description, the framework is the statistical-fluid description; the non-commutative structure of QM gains an ontological motivation);
\item \textbf{Relation to candidate paths}: structural correspondences with Sakharov induced gravity, Verlinde entropic force, 't Hooft cellular automata, causal-set theory, loop quantum gravity (area spectrum/ln 2), string theory (KK/holography) --- the framework provides the "event-flow" language unifying these threads;
\item \textbf{Honesty statement}: no \mtag{To be verified} item may be used as a basis for derivation; the mathematization of the framework has completed its main structure, but several micro-derivations still await rigorization.
\end{itemize}



\section*{Appendix A Notation and conventions}

\begin{itemize}
\item Metric signature (\ensuremath{-},+,+,+); natural units \ensuremath{\hbar}=c=k\down{B}=1 (unless stated);
\item M\down{Pl} = 1.22\ensuremath{\times}10\up{19} GeV; l\down{P} = 1.62\ensuremath{\times}10\up{-35} m; t\down{P} = 5.4\ensuremath{\times}10\up{-44} s;
\item \ensuremath{\varepsilon}\down{A} = energy-scale-dependent leakage probability; \ensuremath{\varepsilon}\down{B} = statistical decay factor;
\item \mtag{Axiom} = fundamental postulate; \mtag{Assumption} = physical assumption introduced by this framework; \mtag{Derivation} = framework-internal derivation; \mtag{Citation} = established result; \mtag{To be verified} = conjecture requiring verification.
\end{itemize}

\section*{Appendix B Sources and AI-assistance statement}

The physical ideas of this framework originate from the author's own research exploration; during conception and derivation, large language models were used as computational and drafting aids (2026). \textbf{All scientific claims, derivations, approximations, and numerical results were proposed and verified by the author}; the language models served only as derivation and writing assistants and do not constitute a scientific authority. The numerical scripts (attachments to the P0--P3 reports) reproduce all numerical results in this paper.



\section*{Appendix C Version notes (development record)}

\begin{quote}
This appendix is the complete record of version changes and history of the v1.1--v1.36 development process. The published body (\S{}1--\S{}10) is the complete statement of the current scientific content and does not depend on this appendix.
\end{quote}

\subsection*{C.1 Change log (v1.21--v1.36)}

\begin{quote}
\textbf{v1.21}: B3 (exponential metric + complete PPN + shadow EHT-compatible) + B4 (direction discretization + Bekenstein bound double closure) + B5 (observer mutual-information area law + 16-fold compression) + B-branch closure (two-dimensional curling / three-layer fluctuation structure / J-K / string-net) integrated into the gravity-side argument
\textbf{v1.22}: P0-3 (\ensuremath{\kappa}\down{c} finite-size scaling) --- framework X-field topological-order\ensuremath{\to}confinement transition \textbf{\ensuremath{\kappa}\down{c} \ensuremath{\approx} 0.25\ensuremath{\pm}0.05} (\S{}5.7); "no transition at finite size" picture gains thermodynamic-limit localization; blueprint P0 stage fully closed
\textbf{v1.23}: P0-3b (torus periodic-boundary refinement) --- (1) torus m=0.5 half-polarization point \textbf{\ensuremath{\kappa}*\ensuremath{\approx}0.20\ensuremath{\pm}0.01} (size-independent, open-boundary pollution removed); (2) \textbf{criterion correction}: m=0.5 crossing is a half-polarization point, not the transition (torus measures m(h=0.32)\ensuremath{\approx}0.84), true transition \textbf{\ensuremath{\kappa}\down{c}\ensuremath{\approx}0.25--0.32}; (3) \textbf{literature-comparison correction}: X-field (\ensuremath{\sigma}\up{x}) version compared with h\down{xc}\ensuremath{\approx}0.32 (0.328 is the \ensuremath{\sigma}\up{z-field} value, mixed up in P0-3, corrected) (\S{}5.7)
\textbf{v1.24}: pinned-EOS micro-derivation (candidate D) --- overload-region EOS P(\ensuremath{\rho})=P\down{P}\ensuremath{\cdot}tanh(\ensuremath{\rho}/\ensuremath{\rho}\down{P}) fully emergent from \textbf{debris--storm collision dynamics} (\S{}4.6 v1.18 section upgraded): (1) \ensuremath{\rho}/\ensuremath{\rho}\down{P} = Planck filling ratio (m\down{P}/l\down{P}\up{3}); (2) tanh = bidirectional net momentum flow (reflection/transmission) logistic saturation; (3) P\down{P}=E\down{P}/l\down{P}\up{3} unit maximum pressure (Axiom 1 + \mtag{Axiom-1 corollary: maximum event frequency}+\mtag{Assumption H-7: finite storm pressure}); \textbf{new prediction: sound-speed profile c\down{s}\up{2}=c\up{2}\ensuremath{\cdot}sech\up{2}(\ensuremath{\rho}/\ensuremath{\rho}\down{P})} (sound-speed vanishing for \ensuremath{\rho}\ensuremath{\gg}\ensuremath{\rho}\down{P} \ensuremath{\implies} structureless paste), energy condition P \ensuremath{\leq} \ensuremath{\rho}c\up{2} everywhere (no collapse)
\textbf{v1.25}: overload-region radial structure and radiative transfer (follow-up item (3), items (2)(3) closed) --- \textbf{isothermal sphere (linear-region hydrostatic equilibrium) gives constant overweight M=4M at r\down{h} \ensuremath{\implies} no net gravitational gradient inside (horizon flux=0 suppression) \ensuremath{\implies} uniform paste \ensuremath{\bar{\rho}}=3c\up{6}/(32\ensuremath{\pi}G\up{3}M\up{2})} (exactly consistent with \ensuremath{\rho}\down{avg}); radiative transfer: macroscopic black-hole optical depth \ensuremath{\tau}\ensuremath{\sim}10\up{-40} (\textbf{optically transparent}, quantitative reason the overload region neither absorbs nor emits), Planck-mass black hole \ensuremath{\tau}\ensuremath{\approx}0.06 (critical window) (\S{}4.6)
\textbf{v1.26}: LISA EMRI 3PN test proposal (candidate E) --- framework 2PN g\down{rr} 33\% bifurcation \ensuremath{\implies} precession difference \ensuremath{\delta}=0.30\ensuremath{\cdot}x\down{p} \textbf{phase-accumulation amplified} in EMRI: \ensuremath{\Delta}\ensuremath{\Phi}=N\ensuremath{\cdot}\ensuremath{\Delta}\ensuremath{\omega}\down{GR}\ensuremath{\cdot}\ensuremath{\delta} reaching 10\up{2}--10\up{3} rad (6 sources / full parameter space testable, margin 1128--513240); 3PN g\down{00} 11\% bifurcation \ensuremath{\implies} chirp phase difference \ensuremath{\sim}2\ensuremath{\times}10\up{2} rad --- EMRI is the \textbf{first verdict window} for the framework's 2PN/3PN bifurcations (complementary to PTA) (\S{}4.5b + \S{}9 S-E1--S-E4)
\textbf{v1.27}: EMRI full waveform template + Fisher precision (follow-up item (4)) --- Peters--Mathews adiabatic evolution + quadrupole waveform (precession difference \ensuremath{\delta} explicitly injected) + LISA noise + \textbf{7-parameter Fisher matrix}: representative source \ensuremath{\delta}=2\% measurable at \textbf{24\ensuremath{\sigma} (2yr) / 167\ensuremath{\sigma} (4yr)} (\ensuremath{\sigma}\down{\delta{}}=8.3\ensuremath{\times}10\up{-4}/1.2\ensuremath{\times}10\up{-4}), GR model excluded \textbf{\ensuremath{\Delta}lnL=288/13930 --- single-source verdict (in-template)} (\S{}4.5b + \S{}9 S-E1)
\textbf{v1.28}: 3+1D string-net rigorization (follow-up item (2), \S{}5.8 semi-rigorous\ensuremath{\to}partially rigorous) --- (1) \textbf{quantum-line--string-net-edge Hopf bridge} (B8 landed, operator-level correspondence); (2) \textbf{rigorous proof of \ensuremath{\lambda}=0 string-net condensation}: E0=\ensuremath{-}N\down{c}, ground-state degeneracy 2\up{n-N\down{c}} (double cube 512 \checkmark), excitation gap \textbf{\ensuremath{\Delta}=2} (emergent particles); (3) \textbf{rigorous duality mapping} (condensed phase \ensuremath{\leftrightarrow} 3D Z\down{2} lattice-gauge confined phase, Fradkin--Shenker); (4) \textbf{emergent U(1) path made explicit} (duality rigorous + continuum-limit candidate); spin 2 still open (world-class) (\S{}5.8 + \S{}10.2 \#13)
\textbf{v1.29}: emergent U(1) continuum limit and spin 2 (follow-up item (2) remainder) --- (1) \textbf{spin 2 = rigorous verification of second-order traceless-symmetric-tensor rotation eigenvalues \ensuremath{\pm}2} (D\up{2} spectrum {cos2\ensuremath{\theta},cos\ensuremath{\theta},1}); (2) emergent U(1) path completed (rigorous duality + Villain compactification + Coulomb-phase Maxwell \ensuremath{\omega}=|k|, literature citation); (3) spin-2 three-layer argument (framework rigorous / string-net semi-rigorous / \textbf{NJU FQH projection observational evidence}); (4) \textbf{testable prediction: FQH graviton-mode decay rate contains \ensuremath{\varepsilon} term} (framework prediction) (\S{}5.8 + \S{}9 S-E5 + \S{}10.2 \#13)
\textbf{v1.30}: Z\down{2}\ensuremath{\to}U(1) rigorous embedding theorem and 4D spin-2 emergence (follow-up item (2), deepest; honest positioning: same order as Millennium problems, not claimed) --- \textbf{Theorem A (rigorous)} Z\down{2} \ensuremath{\subset} U(1) group embedding; \textbf{Proposition B (analytic derivation)} cos(\ensuremath{\theta}\down{p}) linearization = Maxwell; \textbf{Numerical verification C} Coulomb-phase photon \ensuremath{\omega}=|k| (long-wavelength 0.9999); \textbf{spin-2 criterion chain} (Fierz--Pauli uniqueness + rotation eigenvalues \ensuremath{\pm}2 + Weinberg soft theorem + string-net low-energy spectrum) --- the only gap = continuum-limit RG / 4D spectrum (world-class) (\S{}5.8 + \S{}10.2 \#13)
\textbf{v1.31}: accretion-disk radiative transfer (follow-up item (3), item (3) closed) --- standard thin disk (Shakura--Sunyaev multi-temperature blackbody, T\ensuremath{\sim}1.2\ensuremath{\times}10\up{5} K @ 10\up{6} M\ensuremath{\odot} 0.01 \ensuremath{\dot{M}}\down{Edd}) as disk bulk \textbf{consistent with GR} (outer region indistinguishable at 1PN); \textbf{overload region fully transparent} (\ensuremath{\tau}\ensuremath{\sim}10\up{-40}, no secondary absorption of infalling radiation); \textbf{framework-unique differences at the horizon surface} --- firewall non-thermal correction \ensuremath{\sim}\ensuremath{\sqrt{l\down{P}/R}} (primordial black-hole 10\up{-9}--10\up{-4} window) + Hawking--infalling-matter entanglement (\S{}4.6 + \S{}9 S4)
\textbf{v1.32}: Kerr EMRI waveform and Bayesian analysis (follow-up item (4)) --- static template generalized to \textbf{Kerr (a=0.9)} + Fisher\ensuremath{\to}Bayesian (profile likelihood+MCMC): \textbf{\ensuremath{\delta}\down{hat}=0.0200 exactly recovered, \ensuremath{\sigma}\down{\delta{}}=9.7\ensuremath{\times}10\up{-4}, \ensuremath{\Delta}lnL=338 \ensuremath{\implies} GR excluded \ensuremath{\sim}26\ensuremath{\sigma} (robust for Kerr)}; GR control no false positive; \ensuremath{\sigma}\down{\delta{}} same order as static Fisher (8.3\ensuremath{\times}10\up{-4}) (\S{}4.5b + \S{}9 S-E1)
\textbf{v1.33}: non-Abelian 3+1D numerical realization (action-list P0-2) --- \textbf{S\down{3} face model (non-Abelian string-net, 6 faces, 46,656 states)}: \ensuremath{\lambda}=0 string-net condensation E0=\ensuremath{-}N\down{c}, degeneracy \textbf{6\up{n-1}=7776} (group order generalizes Z\down{2}'s 2\up{n-1}); \textbf{non-Abelian fingerprint: excitations split into two conjugacy classes} (transpositions 23,328 / 3-cycles 15,552; absent in Abelian); gap opens monotonically with \ensuremath{\lambda}; \textbf{glueball experimental anchor[BESIII 2026]: BESIII X(2370) (0\up{-+}, 2395 MeV, flavor singlet) confirms glueball = observable instance of framework "flux self-binding"} (lattice-QCD 2.3--3.0 GeV comparison; infrared slavery = string tension\mtag{Framework mechanism: confinement}) (\S{}5.7 + \S{}9 S-E6 + \S{}10.2 \#2)
\end{quote}

\begin{quote}
\textbf{v1.34}: string-net low-energy spectrum and spin-2 criterion evidence (action-list P1-1) --- \textbf{C4 chain-axis fourfold-rotation helicity decomposition} ([R,H]=0, R\up{4}=I exact): scalar/helicity-\ensuremath{\pm}2/vector three-sector projected spectra (2,048--65,536 states); \ensuremath{\lambda}=0 exact level structure fully verified (2\up{n-N\down{c}}\ensuremath{\cdot}C(N\down{c},k) \checkmark); \textbf{spin-2 mode gapped} (\ensuremath{\Delta}\ensuremath{\pm}2\ensuremath{\approx}0.49, double cube 0.4919 / triple chain 0.4864, size-independent); \textbf{carrier evidence: second-order-tensor operator space Z\down{fZ}\down{g} carries 30 helicity-\ensuremath{\pm}2 eigenstates} (\ensuremath{\gg} first-order 3) --- spin 2 = second-order-tensor sector (aligned with \mtag{Framework definition D-2}); honest conclusion: no massless spin 2 at Z\down{2} level (Abelian gapped, expected physics) (\S{}5.8 + \S{}10.2 \#13 criterion (4) numericalized)
\end{quote}

\begin{quote}
\textbf{v1.35}: EMRI a-\ensuremath{\delta} joint MCMC (action-list P1-2, v1.32's "a fixed" note landed; \textbf{v1.32's severe numerical defects discovered and corrected}) --- (1) \textbf{undersampling} (v1.32 dt/carrier-period=34.6 \ensuremath{\to} aliased \ensuremath{\delta}-likelihood comb): analytic inspiral + Nyquist sampling; (2) \textbf{orbital-phase comb} (a enters 2\ensuremath{\lambda}, 4.3\ensuremath{\times}10\up{5} rad in 2yr): \ensuremath{\varphi}0-profile likelihood eliminates; (3) \textbf{Lense--Thirring precession added} (LT\ensuremath{\propto}x\up{3/2} vs 1PN\ensuremath{\propto}x, r-scaling difference breaks \ensuremath{\delta}-a degeneracy) --- \textbf{\ensuremath{\sigma}\down{\delta{}} corrected by 600\ensuremath{\times}: 9.7\ensuremath{\times}10\up{-4} \ensuremath{\to} 1.6\ensuremath{\times}10\up{-6} (2yr; physical \ensuremath{\sigma}\down{\delta{}}\ensuremath{\approx}1/(2wGR\down{T}\ensuremath{\cdot}SNR))}; with a-\ensuremath{\varphi}0 fully free \textbf{\ensuremath{\Delta}lnL=357/396 \ensuremath{\implies} 27\ensuremath{\sigma}/28\ensuremath{\sigma}} (v1.32's 4yr "167\ensuremath{\sigma}" = aliasing inflation); \textbf{main degeneracy \ensuremath{\delta}-\ensuremath{\varphi}0 (\ensuremath{\rho}\ensuremath{\approx}\ensuremath{-}0.85), not \ensuremath{\delta}-a}; GR control no false positive (\S{}4.5b + \S{}9 S-E1)
\end{quote}

\begin{quote}
\textbf{v1.36}: 3+1D \ensuremath{\lambda}\down{c} transition + Wilson RG numerical path (action-list P1-3/P1-4) --- \textbf{\ensuremath{\lambda}\down{c}\ensuremath{\approx}0.14} (\ensuremath{\chi} peaks of four sizes 2\ensuremath{\times}1\ensuremath{\times}1/3\ensuremath{\times}1\ensuremath{\times}1/2\ensuremath{\times}2\ensuremath{\times}1/4\ensuremath{\times}1\ensuremath{\times}1 all locked at \ensuremath{\lambda}\down{\chi{}}=0.140, 2,048--2.1\ensuremath{\times}10\up{6} states; \textbf{dimension--critical-coupling decreasing sequence} with 2D toric code \ensuremath{\kappa}\down{c}\ensuremath{\approx}0.25--0.32, same order as dual 3D Ising \ensuremath{\beta}\down{c}=0.2217; \mtag{Framework mechanism: transition = condensation point}); \textbf{Wilson RG path}: U(1) Gaussian-fixed-point scaling law \ensuremath{\beta}\ensuremath{^\prime}=\ensuremath{\beta}b\up{2-d} (d=2 critical / d=3 IR-free) + 1D Ising exact decimation (\ensuremath{\beta}*=0) + \textbf{2D Ising block RG: \ensuremath{\beta}<\ensuremath{\beta}\down{c} flows to 0 (=U(1) continuum limit), \ensuremath{\beta}>\ensuremath{\beta}\down{c} flows to \ensuremath{\infty} (condensed phase)} --- Z\down{2}\ensuremath{\to}U(1) four-layer evidence closed (Theorem A + Proposition B + Numerical verification C + RG flow); honest note: lattice block-averaging next-nearest-neighbor corrections (\S{}5.8 + \S{}10.2 \#13)
\end{quote}

\subsection*{C.2 Version history}

Draft v1.36 complete. Version history: v1.1 withdrew \S{}4.6 strong-field extrapolation; added \S{}4.10 packing theorem package. v1.2 Theorem 6. v1.3 O1--O4. v1.4 N1--N4. v1.5 P1--P4. v1.6 integrated version. v1.7 hairy ball \ensuremath{\implies} \ensuremath{\varepsilon} lower bound (H1--H4). v1.8 firewall-thickness verdict (F1--F4). v1.9 exact flux-diffusion-equation solution (F1\ensuremath{^\prime}). v1.10 3D spherically symmetric exact solution (F1\ensuremath{^{\prime\prime}}). v1.11 finite-R exact solution (F1\ensuremath{^{\prime\prime\prime}}). v1.12 1/r steady state and weak-field g\down{00} coefficient alignment (factor-2 resolution). v1.13 strong-field bifurcation verdict (PPN \ensuremath{\beta}: Boltzmann \ensuremath{\beta}=1 passes). v1.14 flux-saturation strong-field metrics mathematically excluded. v1.15 storm-overload observational effects (one-parameter family \ensuremath{\alpha} + S1--S5). v1.16 \ensuremath{\alpha} micro-origin (entropy consistency forces \ensuremath{\alpha}=1). v1.17 pressure-critical area locking micro-mechanism (M\up{2}/A Planck criticality). v1.18 overload-region fluid model (Planck-pinned EOS). v1.19 gauge-side emergence argument integration (\S{}5.6 quantum double / \S{}5.7 Hamiltonian / \S{}5.8 string-net). v1.20 confined string-tension numerics (X-field \ensuremath{\sigma}\down{E}\ensuremath{\approx}2\ensuremath{\kappa} + Z-field \ensuremath{\sigma}\down{m}\ensuremath{\approx}2\ensuremath{\lambda}, P0-1 review). \textbf{v1.21 gravity-side completion (B-branch integration)} --- (1) \S{}4.5b: isotropic exponential metric identification (three derivations converging + literature exponential-metric connection) + complete PPN parameter set (\ensuremath{\gamma}=\ensuremath{\beta}=1, others 0) + 2PN/3PN bifurcation quantified (precession difference \ensuremath{\sim}30x\down{p}); (2) \S{}4.6: \textbf{shadow 5.436GM/c\up{2} in the isotropic convention compatible with EHT (+4.55\%), correcting v1.14's "no strong-field metric commitment" to "geometric predictions compatible with observation"}; (3) \S{}4.9: direction discretization (Planck tube diameter \ensuremath{\implies} k=4 \ensuremath{\implies} Bekenstein bound double closure, B4) + observer mutual-information area law (bulk = surface projection \ensuremath{\implies} I=H(\ensuremath{\Sigma}), 16-fold compression, B5); (4) \S{}10.2 updates: \#1 direction discretization \checkmark, \#6 higher-order PPN \checkmark, \#7 two-dimensional curling \checkmark (3\ensuremath{-}1=2 definitional), \#9 fluctuations \checkmark (three-layer structure), \#11 J/K advanced (order O(1)), \#13 string-net confirmed (candidate roadmap). \textbf{v1.22 P0-3 (\ensuremath{\kappa}\down{c} open-boundary FSS 0.23--0.27, \S{}5.7). v1.23 P0-3b (torus refinement: half-polarization \ensuremath{\kappa}*\ensuremath{\approx}0.20\ensuremath{\pm}0.01, true transition \ensuremath{\kappa}\down{c}\ensuremath{\approx}0.25--0.32, literature-comparison correction h\down{xc}\ensuremath{\approx}0.32, \S{}5.7). v1.24 pinned-EOS micro-derivation (P=P\down{P}\ensuremath{\cdot}tanh(\ensuremath{\rho}/\ensuremath{\rho}\down{P}) emergent from debris--storm collisions + sound-speed profile c\down{s}\up{2}=c\up{2}\ensuremath{\cdot}sech\up{2}(\ensuremath{\rho}/\ensuremath{\rho}\down{P}) + energy condition P\ensuremath{\leq}\ensuremath{\rho}c\up{2}, \S{}4.6). v1.25 overload-region radial structure and radiative transfer (isothermal-sphere overweight 4M \ensuremath{\implies} uniform paste \ensuremath{\bar{\rho}}=3c\up{6}/(32\ensuremath{\pi}G\up{3}M\up{2}) + optical transparency \ensuremath{\tau}\ensuremath{\sim}10\up{-40} / Planck window 0.06, \S{}4.6). v1.26 LISA EMRI 3PN test proposal (2PN precession difference \ensuremath{\delta}=0.30\ensuremath{\cdot}x\down{p}, phase accumulation 10\up{2}--10\up{3} rad, 6 sources all testable, S-E1--S-E4, \S{}4.5b/\S{}9). v1.27 EMRI full waveform template + Fisher precision (precession difference \ensuremath{\delta} measurable 24\ensuremath{\sigma}/167\ensuremath{\sigma}, GR excluded \ensuremath{\Delta}lnL 288--13930, single-source verdict, \S{}4.5b/\S{}9 S-E1). v1.28 3+1D string-net rigorization (\ensuremath{\lambda}=0 condensation rigorous proof E0=\ensuremath{-}N\down{c} / degeneracy 2\up{n-N\down{c}} / \ensuremath{\Delta}=2 + rigorous duality + emergent U(1) path made explicit, \S{}5.8/\S{}10.2 \#13). v1.29 emergent U(1) continuum limit and spin 2 (spin 2 = second-order-tensor rotation eigenvalues \ensuremath{\pm}2 rigorous verification + Villain compactification + Coulomb-phase Maxwell \ensuremath{\omega}=|k| + FQH \ensuremath{\varepsilon}-term prediction S-E5, \S{}5.8/\S{}9/\S{}10.2 \#13). v1.30 Z\down{2}\ensuremath{\to}U(1) rigorous embedding theorem and 4D spin 2 (Theorem A/B/C: rigorous group embedding + rigorous Maxwell linearization + photon \ensuremath{\omega}=|k| numerics + spin-2 criterion chain, \S{}5.8/\S{}10.2 \#13). v1.31 accretion-disk radiative transfer (standard thin-disk bulk consistent with GR + overload region fully transparent + firewall non-thermal correction primordial window, \S{}4.6/\S{}9 S4). v1.32 Kerr EMRI waveform and Bayesian (\ensuremath{\delta}\down{hat}=0.0200 exactly recovered, \ensuremath{\sigma}\down{\delta{}}=9.7\ensuremath{\times}10\up{-4}, GR excluded 26\ensuremath{\sigma} robust for Kerr, GR control no false positive, \S{}4.5b/\S{}9 S-E1). v1.33 non-Abelian 3+1D numerics (S\down{3} face model 46,656 states: degeneracy 6\up{n-1}=7776 + two conjugacy-class excitation fingerprints + BESIII glueball X(2370) experimental anchor, \S{}5.7/\S{}9 S-E6/\S{}10.2 \#2). v1.34 string-net low-energy spectrum (C4 helicity decomposition, \ensuremath{\Delta}\ensuremath{\pm}2\ensuremath{\approx}0.49 size-independent, second-order-tensor sector carrier, criterion (4) numericalized, \S{}5.8/\S{}10.2 \#13). v1.35 EMRI a-\ensuremath{\delta} joint MCMC (Nyquist sampling repair + \ensuremath{\varphi}0-profile + Lense--Thirring; \ensuremath{\sigma}\down{\delta{}} corrected 600\ensuremath{\times} 9.7\ensuremath{\times}10\up{-4}\ensuremath{\to}1.6\ensuremath{\times}10\up{-6}, 27\ensuremath{\sigma}/28\ensuremath{\sigma}, \ensuremath{\delta}-\ensuremath{\varphi}0 main degeneracy, GR control no false positive, \S{}4.5b/\S{}9 S-E1). v1.36 3+1D \ensuremath{\lambda}\down{c} transition + Wilson RG (\ensuremath{\lambda}\down{c}\ensuremath{\approx}0.14 four-size lock + dimension sequence; U(1) Gaussian fixed point \ensuremath{\beta}\ensuremath{^\prime}=\ensuremath{\beta}b\up{2-d} + 2D Ising block RG to continuum limit --- four-layer closure Theorem A + Proposition B + Numerical verification C + RG flow, \S{}5.8/\S{}10.2 \#13).}



\section*{References (draft; to be completed to 50+ in the published version)}
\begin{thebibliography}{30}

\bibitem{ref1} Papapetrou A. (1954) Z. Phys. 139, 518; 141, 293. 
\bibitem{ref2} Yilmaz H. (1958) Phys. Rev. 111, 1417; (1971) Phys. Rev. Lett. 27, 1399; (1973) Found. Phys. 3, 195. 
\bibitem{ref3} Rosen N. (1962) Ann. Phys. 17, 269. 
\bibitem{ref4} Makukov M., Mychelkin E. (2018) Phys. Rev. D 98, 064016. 
\bibitem{ref5} Boonserm P., Ngampitipan T., Simpson A., Visser M. (2021) Phys. Rev. D 103, 024019. 
\bibitem{ref6} Ibison M. (2006) "The Yilmaz cosmology".

\bibitem{ref7} Sakharov A.D. (1968) Dokl. Akad. Nauk 177, 70. 
\bibitem{ref8} Jacobson T. (1995) Phys. Rev. Lett. 75, 1260. 
\bibitem{ref9} Padmanabhan T. (2010) arXiv:1003.5665; (2012) arXiv:1206.4916. 
\bibitem{ref10} Verlinde E. (2011) JHEP 04, 029. 
\bibitem{ref11} Bekenstein J.D. (1973) Phys. Rev. D 7, 2333. 
\bibitem{ref12} Hawking S.W. (1975) Commun. Math. Phys. 43, 199. 
\bibitem{ref13} Unruh W.G. (1976) Phys. Rev. D 14, 870. 
\bibitem{ref14} 't Hooft G. (1993) arXiv:gr-qc/9310026; (1999) arXiv:gr-qc/9903084. 
\bibitem{ref15} Susskind L. (1995) J. Math. Phys. 36, 6377.

\bibitem{ref16} Bombelli L. et al. (1987) Phys. Rev. Lett. 59, 521. 
\bibitem{ref17} Malament D. (1977) J. Math. Phys. 18, 1399.

\bibitem{ref18} Kitaev A.Yu. (2003) Ann. Phys. 303, 2. 
\bibitem{ref19} Wen X.-G. (2003) Phys. Rev. D 68, 065003. 
\bibitem{ref20} Levin M., Wen X.-G. (2005) Rev. Mod. Phys. 77, 871. 
\bibitem{ref21} Fradkin E., Shenker S.H. (1979) Phys. Rev. D 19, 3682. 
\bibitem{ref22} Tupitsyn I.S. et al. (2010) Phys. Rev. B 82, 085114. 
\bibitem{ref23} Tammes P.M.L. (1930) Proc. Sect. Sci. K. Akad. Wet. Amsterdam 33, 31. 
\bibitem{ref24} Wilson K.G. (1974) Phys. Rev. D 10, 2445.

\bibitem{ref25} Babak S. et al. (2017) Phys. Rev. D 95, 103012. 
\bibitem{ref26} Huerta E.A., Gair J.R. (2011) Phys. Rev. D 84, 064023. 
\bibitem{ref27} Barack L., Cutler C. (2004) Phys. Rev. D 69, 082005. 
\bibitem{ref28} Peters P.C., Mathews J. (1963) Phys. Rev. 131, 435. 
\bibitem{ref29} Will C.M. (2014) Living Rev. Relativ. 17, 4. 
\bibitem{ref30} Poisson E., Will C.M. (2014) Gravity (Cambridge).

\end{thebibliography}

\end{document}
